% Electromagnetic Waves — Physics Upper Sixth — Cours signé OnBuch+
% Official MINESEC syllabus. Compiler avec : tectonic PHY06-electromagnetic-waves.tex
\newcommand{\DOCMATIERE}{Physics}
\newcommand{\DOCNIVEAU}{Upper Sixth}
\newcommand{\DOCMODULE}{Upper Sixth — Physics}
\newcommand{\DOCLECON}{6}
\newcommand{\DOCTITRE}{Electromagnetic Waves}
% OnBuch+ — préambule « cours signés » ENRICHI (compilé avec tectonic / XeLaTeX)
% Système visuel « Pop » d'OnBuch+ : fond crème (jamais blanc), bordures encre
% épaisses, ombres « dures » décalées, titres Archivo Black en capitales.
% Version MATHÉMATIQUES : graphes (pgfplots/tikz), boîtes pédagogiques
% (définition, propriété, méthode, attention, le savais-tu, exercice résolu,
% exercice, TP, bilan), page de garde et signature OnBuch+ ancrée en bas.
%
% Macros attendues dans main.tex AVANT \input{preamble} :
%   \DOCMATIERE  \DOCNIVEAU  \DOCMODULE  \DOCLECON (numéro)  \DOCTITRE
\documentclass[11pt]{article}

\usepackage{fontspec}
\usepackage[english]{babel}
\usepackage[a4paper,margin=2cm,top=3.4cm,bottom=2.8cm,headheight=1.4cm,footskip=1.3cm]{geometry}
\usepackage{amsmath,amssymb,amsfonts}
\usepackage{mathtools} % \xmapsto, \coloneqq... souvent utilisés par les modèles
\usepackage{cancel} % simplifications barrées (bilans de forces)
% \abs{z} (module, valeur absolue) et \norm{u} (norme), utilisés par les modèles
\providecommand{\abs}{}\let\abs\relax\DeclarePairedDelimiter{\abs}{\lvert}{\rvert}
\providecommand{\norm}{}\let\norm\relax\DeclarePairedDelimiter{\norm}{\lVert}{\rVert}
\usepackage[table]{xcolor} % [table] : fournit aussi \rowcolors (alternance de lignes)
\usepackage{pagecolor}
\usepackage{tikz}
\usetikzlibrary{arrows.meta,positioning,calc,shapes.geometric,shapes.misc,shapes.arrows,decorations.pathmorphing,decorations.pathreplacing,decorations.markings,patterns,shadows,mindmap,trees,fit,backgrounds,matrix,intersections,angles,quotes}
\usepackage{pgfplots}
\usepackage[version=4]{mhchem} % équations nucléaires \ce{^{235}_{92}U}
\usepackage[european,siunitx]{circuitikz} % schémas électriques
\pgfplotsset{compat=1.18}
\usepgfplotslibrary{fillbetween,groupplots} % aires entre courbes (calcul intégral)
\usepackage{enumitem}
\usepackage{fancyhdr}
% [most] charge toutes les bibliothèques tcolorbox (listings, minted, raster,
% documentation...) et gonfle inutilement la mémoire TeX sur un document long
% (~300 boîtes) : on ne charge que ce qui est réellement utilisé.
\usepackage[skins,breakable]{tcolorbox}
\usepackage{graphicx}
\usepackage{colortbl}
\usepackage{booktabs}
\usepackage{tabularx}
\usepackage{diagbox} % case d'en-tête diagonale des tableaux à double entrée
% Colonnes extensibles centrée / gauche / droite (C, L, R), souvent utilisées par les modèles
\newcolumntype{C}{>{\centering\arraybackslash}X}
\newcolumntype{L}{>{\raggedright\arraybackslash}X}
\newcolumntype{R}{>{\raggedleft\arraybackslash}X}
\usepackage{array}
\usepackage{multicol}
\usepackage{siunitx}
% parse-numbers=false : accepte aussi \SI{2,5 \times 10^{-3}}{...}
\sisetup{locale=UK,output-decimal-marker={.},group-minimum-digits=4,parse-numbers=false}
\DeclareSIUnit\ohm{\ensuremath{\symup{\Omega}}} % U+2126 absent de la police
\DeclareSIUnit\division{div} % réglages d'oscilloscope (ms/div, V/div)
\DeclareSIUnit\tr{tr} % tours (vitesse de rotation, ex. \SI{1500}{\tr\per\minute})
\usepackage{qrcode}
% \degree : commande fréquemment utilisée par les modèles pour un angle ou une
% température en dehors de \SI{}{} (non fournie par les paquets chargés).
\providecommand{\degree}{^{\circ}}
% \qed (fin de démonstration), utilisé par les modèles sans amsthm
\providecommand{\qed}{\hfill$\square$}
% \barre : double barre des valeurs interdites dans un tableau de variations (tkz-tab)
\providecommand{\barre}{\ensuremath{\|}}
% environnement proof (amsthm non chargé) utilisé par les modèles
\ifdefined\proof\else\newenvironment{proof}[1][Proof]{\par\noindent\textit{#1.}\ }{\qed\par}\fi
\usepackage{lastpage}
\usepackage{hyperref}
\hypersetup{hidelinks}

% ============================================================
% Palette « Pop » OnBuch+ — mêmes hex que la classe P (plus_ui.dart)
% ============================================================
\definecolor{poporange}{HTML}{F59321}
\definecolor{popdark}{HTML}{E07A0C}
\definecolor{popcream}{HTML}{FFF6E8}
\definecolor{popink}{HTML}{1C1714}
\definecolor{poppurple}{HTML}{7A5AE0}
\definecolor{popgreen}{HTML}{2E9E6A}
\definecolor{poppink}{HTML}{E0457B}
\definecolor{popblue}{HTML}{2F7DE1}
\definecolor{popgold}{HTML}{C98A12}
\definecolor{popmuted}{HTML}{8A7F76}
\definecolor{popline}{HTML}{EADFD2}
\definecolor{popgray}{HTML}{8A7F76}
\colorlet{popgrey}{popgray}
% Alias tolérants (le modèle nomme parfois les couleurs différemment)
\colorlet{popwhite}{white}
\colorlet{popblack}{popink}
\colorlet{popred}{poppink}
\colorlet{popyellow}{popgold}
% Teintes claires pour fonds de tableaux / zones
\colorlet{popblueL}{popblue!10!white}
\colorlet{poporangeL}{poporange!14!white}
\colorlet{popgreenL}{popgreen!12!white}
\colorlet{poppurpleL}{poppurple!12!white}
\colorlet{poppinkL}{poppink!12!white}
\colorlet{popgoldL}{popgold!14!white}

\pagecolor{popcream}

% ============================================================
% Typographie
% ============================================================
\defaultfontfeatures{Ligatures=TeX}
\setmainfont{PlusJakartaSans-Medium.ttf}[Path=../../../fonts/,BoldFont=PlusJakartaSans-Bold.ttf]
% Maths en sans-serif (Fira Math) avec chiffres et lettres droites de Plus
% Jakarta, pour que les formules se fondent dans le texte.
\usepackage[math-style=ISO,bold-style=ISO]{unicode-math}
\setmathfont{FiraMath-Regular.otf}[Path=../../../fonts/]
\setmathfont{PlusJakartaSans-Medium.ttf}[Path=../../../fonts/,range={up/num,up/{Latin,latin},\mathup}]
% (icomma retiré : décimales avec point en anglais)
% Caractères mathématiques Unicode absents des polices de texte (Plus Jakarta,
% Archivo Black) : rendus par leur équivalent mathématique, en texte comme en
% formule — sinon ils s'affichent en carré vide (ex. « Arithmétique dans ℤ »).
\usepackage{newunicodechar}
\newunicodechar{ℝ}{\ensuremath{\mathbb{R}}}
\newunicodechar{ℤ}{\ensuremath{\mathbb{Z}}}
\newunicodechar{ℂ}{\ensuremath{\mathbb{C}}}
\newunicodechar{ℕ}{\ensuremath{\mathbb{N}}}
\newunicodechar{ℚ}{\ensuremath{\mathbb{Q}}}
\newunicodechar{𝔻}{\ensuremath{\mathbb{D}}}
\newunicodechar{α}{\ensuremath{\alpha}}
\newunicodechar{β}{\ensuremath{\beta}}
\newunicodechar{γ}{\ensuremath{\gamma}}
\newunicodechar{δ}{\ensuremath{\delta}}
\newunicodechar{ε}{\ensuremath{\varepsilon}}
\newunicodechar{θ}{\ensuremath{\theta}}
\newunicodechar{λ}{\ensuremath{\lambda}}
\newunicodechar{μ}{\ensuremath{\mu}}
\newunicodechar{σ}{\ensuremath{\sigma}}
\newunicodechar{τ}{\ensuremath{\tau}}
\newunicodechar{φ}{\ensuremath{\varphi}}
\newunicodechar{ψ}{\ensuremath{\psi}}
\newunicodechar{ω}{\ensuremath{\omega}}
\newunicodechar{Σ}{\ensuremath{\Sigma}}
% majuscules produites par \MakeUppercase dans les titres (α-aminés -> Α)
\newunicodechar{Α}{\ensuremath{\alpha}}
\newunicodechar{Β}{\ensuremath{\beta}}
\newunicodechar{Γ}{\ensuremath{\Gamma}}
\newunicodechar{Δ}{\ensuremath{\Delta}}
\newunicodechar{Ε}{\ensuremath{\varepsilon}}
\newunicodechar{Θ}{\ensuremath{\Theta}}
\newunicodechar{Λ}{\ensuremath{\Lambda}}
\newunicodechar{Μ}{\ensuremath{\mu}}
\newunicodechar{Τ}{\ensuremath{\tau}}
\newunicodechar{Φ}{\ensuremath{\Phi}}
\newunicodechar{Ψ}{\ensuremath{\Psi}}
\newunicodechar{Ω}{\ensuremath{\Omega}}
\newunicodechar{↦}{\ensuremath{\mapsto}}
\newunicodechar{∈}{\ensuremath{\in}}
\newunicodechar{∉}{\ensuremath{\notin}}
\newunicodechar{∧}{\ensuremath{\wedge}}
\newunicodechar{⇔}{\ensuremath{\Leftrightarrow}}
\newunicodechar{⟺}{\ensuremath{\Longleftrightarrow}}
\newunicodechar{⇒}{\ensuremath{\Rightarrow}}
\newunicodechar{⟹}{\ensuremath{\Longrightarrow}}
\newunicodechar{∘}{\ensuremath{\circ}}
\newunicodechar{∪}{\ensuremath{\cup}}
\newunicodechar{∩}{\ensuremath{\cap}}
\newunicodechar{≡}{\ensuremath{\equiv}}
\newunicodechar{⊂}{\ensuremath{\subset}}
\newunicodechar{⊥}{\ensuremath{\perp}}
\newunicodechar{∃}{\ensuremath{\exists}}
\newunicodechar{∀}{\ensuremath{\forall}}
\newunicodechar{∛}{\ensuremath{\sqrt[3]{\,}}}
\newunicodechar{∜}{\ensuremath{\sqrt[4]{\,}}}
\newunicodechar{⃗}{\ensuremath{{}^{\to}}}
\newunicodechar{ⁿ}{\ifmmode^{n}\else\textsuperscript{\ensuremath{n}}\fi}
\newunicodechar{ˣ}{\ifmmode^{x}\else\textsuperscript{\ensuremath{x}}\fi}
\newunicodechar{ᵇ}{\ifmmode^{b}\else\textsuperscript{\ensuremath{b}}\fi}
\newunicodechar{ʳ}{\ifmmode^{r}\else\textsuperscript{\ensuremath{r}}\fi}
\newunicodechar{ᵘ}{\ifmmode^{u}\else\textsuperscript{\ensuremath{u}}\fi}
\newunicodechar{ʸ}{\ifmmode^{y}\else\textsuperscript{\ensuremath{y}}\fi}
\newunicodechar{ᵃ}{\ifmmode^{a}\else\textsuperscript{\ensuremath{a}}\fi}
\newunicodechar{ᵉ}{\ifmmode^{e}\else\textsuperscript{\ensuremath{e}}\fi}
\newunicodechar{ᵏ}{\ifmmode^{k}\else\textsuperscript{\ensuremath{k}}\fi}
\newunicodechar{ᵖ}{\ifmmode^{p}\else\textsuperscript{\ensuremath{p}}\fi}
\newunicodechar{ᴺ}{\ifmmode^{N}\else\textsuperscript{\ensuremath{N}}\fi}
\newunicodechar{ᶜ}{\ifmmode^{c}\else\textsuperscript{\ensuremath{c}}\fi}
\newunicodechar{ʰ}{\ifmmode^{h}\else\textsuperscript{\ensuremath{h}}\fi}
\newunicodechar{ᵠ}{\ifmmode^{q}\else\textsuperscript{\ensuremath{q}}\fi}
\newunicodechar{ₙ}{\ifmmode_{n}\else\textsubscript{\ensuremath{n}}\fi}
\newunicodechar{ₐ}{\ifmmode_{a}\else\textsubscript{\ensuremath{a}}\fi}
\newunicodechar{ᵢ}{\ifmmode_{i}\else\textsubscript{\ensuremath{i}}\fi}
\newunicodechar{ᵣ}{\ifmmode_{r}\else\textsubscript{\ensuremath{r}}\fi}
\newunicodechar{ₖ}{\ifmmode_{k}\else\textsubscript{\ensuremath{k}}\fi}
\newunicodechar{ᵦ}{\ifmmode_{\beta}\else\textsubscript{\ensuremath{\beta}}\fi}
\newunicodechar{ₓ}{\ifmmode_{x}\else\textsubscript{\ensuremath{x}}\fi}
\newfontfamily\popdisplay{ArchivoBlack-Regular.ttf}[Path=../../../fonts/]
\newfontfamily\poplogo{SpaceGrotesk-Bold.ttf}[Path=../../../fonts/]
\newfontfamily\popsemi{PlusJakartaSans-SemiBold.ttf}[Path=../../../fonts/]
\newfontfamily\popxbold{PlusJakartaSans-ExtraBold.ttf}[Path=../../../fonts/]
\newfontfamily\popxboldls{PlusJakartaSans-ExtraBold.ttf}[Path=../../../fonts/,LetterSpace=3.0]

\setlist[enumerate]{leftmargin=1.7em,itemsep=5pt,topsep=4pt}
\setlist[itemize]{leftmargin=1.7em,itemsep=4pt,topsep=4pt,label={\color{poporange}\textbullet}}
\setlength{\parindent}{0pt}
\setlength{\parskip}{5pt}
\renewcommand{\arraystretch}{1.25}

% pgfplots : style Pop par défaut
\pgfplotsset{
  every axis/.append style={axis line style={popink,line width=1pt},tick style={popink},
    grid style={popline,line width=0.5pt},label style={font=\small},tick label style={font=\footnotesize}},
  popcurve/.style={poporange,line width=1.8pt,no markers,smooth},
}
% Même style disponible dans un \draw TikZ simple (hors pgfplots)
\tikzset{popcurve/.style={poporange,line width=1.8pt}}
\tikzset{popgrid/.style={popline,very thin}}
\colorlet{popcurve}{poporange} % le nom du style est parfois utilisé comme couleur

% ============================================================
% Pill / squiggle
% ============================================================
\newcommand{\poppill}[3]{%
  \tikz[baseline=-0.55ex]\node[draw=popink,line width=1.2pt,rounded corners=2.6pt,%
    fill=#2,inner xsep=6.5pt,inner ysep=3pt]{%
    \popxboldls\fontsize{7}{7}\selectfont\color{#3}\MakeUppercase{#1}};%
}
\newcommand{\popsquiggle}[1]{%
  \tikz[baseline=-0.4ex]{
    \draw[#1,line width=2.6pt,line cap=round]
      plot [smooth] coordinates {(0,0.09) (0.34,0.01) (0.68,0.21) (1.02,0.01) (1.36,0.10)};
  }%
}
% Mot-clé surligné (marqueur orange sous le texte)
\newcommand{\cle}[1]{{\popxbold\color{popdark}#1}}

% ============================================================
% En-tête / pied
% ============================================================
\pagestyle{fancy}
\fancyhf{}
\renewcommand{\headrulewidth}{0pt}
\renewcommand{\footrulewidth}{0pt}
\lhead{\raisebox{-0.15ex}{\poplogo\fontsize{13}{13}\selectfont\color{popink}OnBuch\color{poporange}+}}
\rhead{\poppill{\DOCMATIERE\ \textbullet\ \DOCNIVEAU\ \textbullet\ Chapter \DOCLECON}{white}{popink}}
\lfoot{\popxbold\fontsize{7.6}{7.6}\selectfont\color{popmuted}\MakeUppercase{Signed course OnBuch+}\ \textbullet\ {\popsemi\DOCTITRE}}
\rfoot{\tikz[baseline=-0.6ex]\node[rounded corners=8.5pt,fill=popink,minimum height=17pt,minimum width=17pt,inner xsep=5pt,inner sep=0]{\popxbold\fontsize{8}{8}\selectfont\color{popcream}\thepage\,/\,\pageref*{LastPage}};}

% ============================================================
% Titres
% ============================================================
\newcommand{\chaptitre}[2]{%
  \begin{tcolorbox}[enhanced,colback=poporange,colframe=popink,boxrule=2.6pt,arc=11pt,
      drop shadow=popink,left=15pt,right=15pt,top=12pt,bottom=11pt,before skip=3pt,after skip=18pt]
    \poppill{#2}{popink}{popcream}\par\vspace{9pt}
    {\raggedright\hyphenpenalty=10000\exhyphenpenalty=10000\popdisplay\fontsize{22}{24}\selectfont\color{popcream}\MakeUppercase{#1}\par}
  \end{tcolorbox}
}
% Section numérotée : pastille orange avec le numéro + titre Archivo
\newcounter{popsec}
\newcommand{\coursec}[1]{%
  \stepcounter{popsec}\setcounter{popsub}{0}%
  \par\vspace{16pt}\noindent
  \begin{minipage}{\linewidth}
  \tikz[baseline=-0.7ex]\node[draw=popink,line width=1.6pt,fill=poporange,rounded corners=4pt,minimum size=22pt,inner sep=0]{\popdisplay\fontsize{12}{12}\selectfont\color{popcream}\thepopsec};%
  \hspace{8pt}{\popdisplay\fontsize{15}{17}\selectfont\color{popink}\MakeUppercase{#1}}\par
  \vspace{4pt}\noindent{\color{poporange}\rule{36pt}{3.6pt}}
  \end{minipage}\par\nopagebreak\vspace{8pt}
}
\newcounter{popsub}
\newcommand{\courssub}[1]{%
  \stepcounter{popsub}%
  \par\vspace{9pt}\noindent{\popxbold\fontsize{11.5}{13}\selectfont\color{popink}{\color{poporange}\thepopsec.\thepopsub}\hspace{6pt}#1}\par\nopagebreak\vspace{4pt}
}

% ============================================================
% Boîtes pédagogiques — même recette : fond blanc, bordure épaisse colorée,
% ombre dure encre, badge plein. Titre optionnel [#1].
% ============================================================
\tcbset{popbase/.style={breakable,enhanced,colback=white,boxrule=2.3pt,arc=9pt,
  shadow={3pt}{-3pt}{0pt}{popink},left=14pt,right=14pt,top=11pt,bottom=11pt,before skip=11pt,after skip=15pt,
  fontupper=\popsemi\fontsize{10.6}{14.5}\selectfont\color{popink},
  fontlower=\popsemi\fontsize{10.6}{14.5}\selectfont\color{popink}}}
% Coche dessinée (le glyphe ✓ n'existe pas dans les polices du document)
\newcommand{\popcheck}[1]{\tikz[baseline=-0.5ex]\draw[#1,line width=1.6pt,line cap=round,line join=round](-0.13,0.0)--(-0.04,-0.09)--(0.13,0.1);}
\providecommand{\coche}{\popcheck{popgreen}}
\providecommand{\resultat}[1]{\par\smallskip\noindent\fcolorbox{popgreen}{popgreenL}{\parbox{\dimexpr\linewidth-2\fboxsep-2\fboxrule\relax}{\textbf{Result.} #1}}\par\smallskip}
\newcommand{\popboxhead}[3]{\poppill{#1}{#2}{white}\ifx\relax#3\relax\else\hspace{6pt}{\popxbold\fontsize{10.6}{12}\selectfont\color{popink}#3}\fi\par\vspace{7pt}}

\newtcolorbox{aretenir}[1][]{popbase,colframe=popgreen,before upper={\popboxhead{Key points}{popgreen}{#1}}}
\newtcolorbox{exemplebox}[1][]{popbase,colframe=poporange,before upper={\popboxhead{Exemple}{poporange}{#1}}}
\newtcolorbox{definition}[1][]{popbase,colframe=popblue,before upper={\popboxhead{Definition}{popblue}{#1}}}
\newtcolorbox{propriete}[1][]{popbase,colframe=poppurple,before upper={\popboxhead{Property}{poppurple}{#1}}}
\newtcolorbox{methode}[1][]{popbase,colframe=popgold,before upper={\popboxhead{Method}{popgold}{#1}}}
\newtcolorbox{attention}[1][]{popbase,colframe=poppink,before upper={\popboxhead{Warning}{poppink}{#1}}}
\newtcolorbox{experience}[1][]{popbase,colframe=popblue,colback=popblueL,before upper={\popboxhead{Experiment}{popblue}{#1}}}
% « Le savais-tu ? » : carte violette pleine (contexte, culture, Cameroun)
\newtcolorbox{savaistu}[1][]{popbase,colback=poppurple,colframe=popink,
  fontupper=\popsemi\fontsize{10.4}{14.2}\selectfont\color{white},
  before upper={\poppill{Did you know?}{white}{poppurple}\ifx\relax#1\relax\else\hspace{6pt}{\popxbold\color{white}#1}\fi\par\vspace{7pt}}}
% Exercice résolu : énoncé en haut, \tcblower puis la solution sur fond crème
\newcounter{popexo}\newcounter{popexr}
\newtcolorbox{exoresolu}[1][]{popbase,colframe=poporange,colbacklower=popcream,
  segmentation style={popink,line width=1.2pt,dashed},
  before upper={\stepcounter{popexr}\popboxhead{Worked example \thepopexr}{poporange}{#1}},
  before lower={\poppill{Solution}{popink}{popcream}\par\vspace{6pt}}}
% Exercice (non résolu en place) — la correction est dans la partie Corrigés
\newtcolorbox{exercice}[1][]{popbase,colframe=popink,
  before upper={\stepcounter{popexo}\popboxhead{Exercise \thepopexo}{popink}{#1}}}
% Corrigé
\newtcolorbox{corrige}[1][]{popbase,colframe=popgreen,colback=popgreenL,
  before upper={\popboxhead{Solution}{popgreen}{#1}}}
% Objectifs / prérequis / bilan
\newtcolorbox{objectifs}{popbase,colframe=popink,colback=poporangeL,
  before upper={\popboxhead{Chapter objectives}{popink}{}}}
\newtcolorbox{prerequis}{popbase,colframe=popink,colback=popblueL,
  before upper={\popboxhead{Prerequisites}{popblue}{}}}
\newtcolorbox{bilan}{popbase,colframe=popink,colback=white,boxrule=2.8pt,
  before upper={\popboxhead{Fiche bilan}{poporange}{L'essentiel à connaître pour l'examen}}}
% Figure encadrée avec légende
\newtcolorbox{popfigure}[1][]{enhanced,breakable=false,colback=white,colframe=popline,boxrule=1.4pt,arc=8pt,
  left=10pt,right=10pt,top=10pt,bottom=8pt,before skip=10pt,after skip=12pt,halign=center,
  lower separated=false,halign lower=center,
  fontlower=\popsemi\fontsize{9}{11.5}\selectfont\color{popmuted},
  #1}
\newcommand{\legende}[1]{\par\vspace{6pt}{\popsemi\fontsize{9}{11.5}\selectfont\color{popmuted}#1}}

% ============================================================
% Signature OnBuch+
% ============================================================
\newcommand{\poplogoinline}[1]{{\poplogo\fontsize{#1}{#1}\selectfont\color{popink}OnBuch\color{poporange}+}}
\newcommand{\signatureonbuch}{%
  \begin{tcolorbox}[enhanced,colback=popink,colframe=popink,boxrule=2.2pt,arc=10pt,
      drop shadow=poporange,left=14pt,right=14pt,top=10pt,bottom=10pt,before skip=0pt,after skip=0pt]
    \tikz[baseline=-0.7ex]\node[circle,fill=popgreen,draw=popcream,line width=1.3pt,minimum size=22pt,inner sep=0]{\popcheck{white}};%
    \hspace{9pt}\begin{minipage}[c]{0.62\linewidth}
      {\popxbold\fontsize{12}{13}\selectfont\color{popcream}Signed\ \textbullet\ {\poplogo OnBuch\color{poporange}+}}\par\vspace{2pt}
      {\raggedright\popsemi\fontsize{8}{10.5}\selectfont\color{popline}\MakeUppercase{OnBuch+ teaching team}\\\MakeUppercase{Follows the official MINESEC syllabus}\par}
    \end{minipage}\hfill
    {\poplogo\fontsize{22}{22}\selectfont\color{popcream}OnBuch\color{poporange}+}
  \end{tcolorbox}%
}

% ============================================================
% Page de garde
% #1 titre · #2 matière · #3 niveau · #4 module · #5 n° leçon · #6 durée ·
% #7 description / objectifs (texte ou itemize)
% ============================================================
\newcommand{\pagegarde}[7]{%
  \thispagestyle{empty}
  \newgeometry{a4paper,margin=1.7cm,top=1.6cm,bottom=1.4cm}%
  \begin{tikzpicture}[remember picture,overlay]
    \fill[poporange!18] ([xshift=-3.2cm,yshift=-3.6cm]current page.north east) circle (4.6cm);
    \fill[poppurple!14] ([xshift=2.4cm,yshift=7.6cm]current page.south west) circle (3.8cm);
  \end{tikzpicture}%
  {\poplogo\fontsize{30}{30}\selectfont\color{popink}OnBuch\color{poporange}+}\hfill
  \raisebox{0.6ex}{\poppill{Signed course \textbullet\ Premium}{popink}{popcream}}\par
  \vspace{1pt}\popsquiggle{poppurple}\par
  \vspace{8pt}
  \begin{tcolorbox}[enhanced,colback=poporange,colframe=popink,boxrule=2.8pt,arc=13pt,
      drop shadow=popink,left=18pt,right=18pt,top=12pt,bottom=13pt,after skip=10pt]
    \poppill{#2 \textbullet\ #3}{popink}{popcream}\hspace{5pt}\poppill{#4}{white}{popink}\par\vspace{8pt}
    {\popdisplay\fontsize{14}{14}\selectfont\color{popink}CHAPTER #5}\par\vspace{4pt}
    {\raggedright\hyphenpenalty=10000\exhyphenpenalty=10000\popdisplay\fontsize{25}{27}\selectfont\color{popcream}\MakeUppercase{#1}\par}
    \vspace{8pt}
    {\popxbold\fontsize{9.5}{9.5}\selectfont\color{popink}\MakeUppercase{Suggested time: #6}}
  \end{tcolorbox}
  \begin{tcolorbox}[enhanced,colback=white,colframe=popink,boxrule=2.2pt,arc=10pt,
      drop shadow=popink,left=16pt,right=16pt,top=10pt,bottom=10pt,after skip=8pt]
    \poppill{In this chapter}{poppurple}{white}\par\vspace{5pt}
    \popsemi\fontsize{9.3}{12.6}\selectfont\color{popink}#7
  \end{tcolorbox}
  \noindent\begin{minipage}[t]{0.487\linewidth}
    \begin{tcolorbox}[enhanced,colback=popgreen,colframe=popink,boxrule=2.2pt,arc=10pt,
        drop shadow=popink,left=11pt,right=11pt,top=10pt,bottom=10pt,height=3.1cm,valign=center]
      \tcbox[colback=white,colframe=popink,boxrule=1.3pt,arc=5pt,boxsep=0pt,nobeforeafter,
        left=3pt,right=3pt,top=3pt,bottom=3pt,valign=center]{\qrcode[height=1.9cm]{https://play.google.com/store/apps/details?id=com.luvvix.onbuch&hl=en}}%
      \hspace{8pt}\begin{minipage}[c]{0.5\linewidth}
        {\popdisplay\fontsize{11}{12.5}\selectfont\color{white}\MakeUppercase{Download OnBuch}}\par\vspace{4pt}
        {\raggedright\popsemi\fontsize{8.4}{11}\selectfont\color{white}Free \textbullet\ Google Play\\AI tutor, past papers, results\par}
      \end{minipage}
    \end{tcolorbox}
  \end{minipage}\hfill
  \begin{minipage}[t]{0.487\linewidth}
    \begin{tcolorbox}[enhanced,colback=poppurple,colframe=popink,boxrule=2.2pt,arc=10pt,
        drop shadow=popink,left=11pt,right=11pt,top=10pt,bottom=10pt,height=3.1cm,valign=center]
      \tikz[baseline=-0.7ex]\node[circle,fill=white,draw=popink,line width=1.3pt,minimum size=1.6cm,inner sep=0]{\popdisplay\fontsize{24}{24}\selectfont\color{poppurple}+};%
      \hspace{8pt}\begin{minipage}[c]{0.6\linewidth}
        {\popdisplay\fontsize{11}{12.5}\selectfont\color{white}\MakeUppercase{Go OnBuch+}}\par\vspace{4pt}
        {\raggedright\popsemi\fontsize{8.4}{11}\selectfont\color{white}All signed courses, unlimited AI tutor, PDF sheets — OnBuch+ tab in the app\par}
      \end{minipage}
    \end{tcolorbox}
  \end{minipage}
  \vspace{8pt}
  \popcontenu
  \vfill
  \signatureonbuch
  \restoregeometry
  \newpage
}

% Rangée « Ce cours contient » (page de garde)
\newcommand{\poptile}[3]{% #1 couleur #2 gros texte #3 légende
  \begin{tcolorbox}[enhanced,colback=#1,colframe=popink,boxrule=2pt,arc=9pt,shadow={2.5pt}{-2.5pt}{0pt}{popink},
      left=6pt,right=6pt,top=9pt,bottom=9pt,halign=center,valign=center,height=1.75cm,width=\linewidth,nobeforeafter]
    {\popdisplay\fontsize{12.5}{13}\selectfont\color{white}\MakeUppercase{#2}}\par\vspace{3pt}
    {\centering\popsemi\fontsize{7.8}{9.5}\selectfont\color{white}#3\par}
  \end{tcolorbox}}
\newcommand{\popcontenu}{%
  \par\noindent{\popxboldls\fontsize{7.5}{7.5}\selectfont\color{popmuted}THIS SIGNED COURSE CONTAINS}\par\vspace{7pt}
  \noindent\begin{minipage}[t]{0.235\linewidth}\poptile{popblue}{Course}{complete, illustrated, step by step}\end{minipage}\hfill
  \begin{minipage}[t]{0.235\linewidth}\poptile{poporange}{Methods}{and worked examples}\end{minipage}\hfill
  \begin{minipage}[t]{0.235\linewidth}\poptile{poppink}{Exercises}{exam style + solutions}\end{minipage}\hfill
  \begin{minipage}[t]{0.235\linewidth}\poptile{popgreen}{Summary sheet}{the essentials to remember}\end{minipage}\par}

% Sommaire
\newtcolorbox{sommaire}{enhanced,colback=white,colframe=popink,boxrule=2.2pt,arc=10pt,
  drop shadow=popink,left=18pt,right=18pt,top=13pt,bottom=13pt,before skip=6pt,after skip=20pt,
  before upper={\poppill{Contents}{poppurple}{white}\par\vspace{9pt}\popsemi\fontsize{10.6}{14.5}\selectfont\color{popink}}}

% Signature de fin de cours — poussée tout en bas de la dernière page
\newcommand{\findecours}{%
  \par\vfill\nopagebreak
  \begin{center}\popsquiggle{poporange}\end{center}\vspace{4pt}
  \signatureonbuch
}
\AtBeginDocument{\def\llbracket{\mathopen{[\mkern-3.5mu[}}\def\rrbracket{\mathclose{]\mkern-3.5mu]}}} % intervalles d'entiers (glyphe ⟦ absent de la police)
% Glyphes absents de Fira Math : redéfinis à partir de glyphes disponibles
\AtBeginDocument{%
  \def\setminus{\mathbin{\backslash}}\let\smallsetminus\setminus
  \def\vdots{\mathord{\vbox{\baselineskip3.2pt\lineskiplimit0pt\kern4pt\hbox{.}\hbox{.}\hbox{.}}}}%
  \def\ddots{\mathinner{\mkern1mu\raise6.5pt\hbox{.}\mkern2mu\raise3.6pt\hbox{.}\mkern2mu\raise0.7pt\hbox{.}\mkern1mu}}%
  \def\triangle{\mathord{\scalebox{0.9}{$\symup{\Delta}$}}}%
  \def\nabla{\mathord{\raisebox{\depth}{\scalebox{1}[-1]{$\symup{\Delta}$}}}}%
  \def\checkmark{\popcheck{popdark}}%
  \def\star{\mathbin{\ast}}\def\bigstar{\mathord{\ast}}%
  \def\diamond{\mathbin{\scalebox{0.6}{$\Diamond$}}}%
  \def\bigcup{\mathop{\mathchoice{\raisebox{-0.3em}{\scalebox{1.7}{$\cup$}}}{\scalebox{1.25}{$\cup$}}{\cup}{\cup}}\displaylimits}%
  \def\bigcap{\mathop{\mathchoice{\raisebox{-0.3em}{\scalebox{1.7}{$\cap$}}}{\scalebox{1.25}{$\cap$}}{\cap}{\cap}}\displaylimits}%
  \def\lesseqqgtr{\mathrel{\lessgtr}}\def\bigoplus{\mathop{\oplus}\displaylimits}%
}
\providecommand{\ssi}{if and only if} % abréviation fréquente
\AtBeginDocument{\providecommand{\wideparen}{\overparen}} % arc de cercle

%% FIGFIX-BEGIN v5 — correctifs globaux des figures (docs/onbuch-generation/tools/figfix)
\usepackage{adjustbox}\usepackage{etoolbox}\tcbuselibrary{raster}
% toute figure TikZ/pgfplots plus large que la ligne est réduite à la largeur disponible
\BeforeBeginEnvironment{tikzpicture}{\begin{adjustbox}{max width=\linewidth}}
\AfterEndEnvironment{tikzpicture}{\end{adjustbox}}
% légendes pgfplots lisibles par-dessus les courbes
\pgfplotsset{every axis/.append style={legend style={fill=white,fill opacity=0.92,text opacity=1,draw=black!15,inner sep=3pt}}}
% étiquettes d'arêtes (syntaxe edge["..."]) avec fond blanc
\tikzset{every edge quotes/.append style={fill=white,inner sep=1.2pt}}
%% FIGFIX-END
\begin{document}

\pagegarde{Electromagnetic Waves}{Physics}{Upper Sixth}{Upper Sixth — Physics}{6}{9 periods}{This chapter introduces electromagnetic waves, their spectrum, production and uses, then develops the wave behaviour of light: diffraction and interference with gratings, reflection, refraction, dispersion, lenses, prisms and optical instruments. It closes with the energy of waves in free space and the inverse square law, preparing students for GCE Advanced Level Paper 2 and Paper 3.
\vspace{4pt}
\begin{multicols}{2}\small
\begin{itemize}
\item By the end of this chapter I can describe an experiment to measure the speed of sound in air and compare it with the speed of light.
\item By the end of this chapter I can state the characteristics of electromagnetic waves and order the regions of the EM spectrum by wavelength and frequency.
\item By the end of this chapter I can explain the production, detection and uses of each region of the EM spectrum, including X-rays.
\item By the end of this chapter I can explain plane polarisation of EM waves and give its applications.
\item By the end of this chapter I can use the grating equation d sinθ = nλ to determine wavelengths and analyse multiple-slit diffraction.
\item By the end of this chapter I can state and apply the laws of reflection and refraction and use Snell's law with refractive index.
\end{itemize}\end{multicols}}

\begin{sommaire}
\begin{multicols}{2}
\begin{enumerate}
\item Section 1: Measuring the Speed of Sound and the Nature of Electromagnetic Waves
\item Section 2: The Electromagnetic Spectrum, X-rays and Polarisation
\item Section 3: Light Sources, Diffraction Gratings and Multiple-Slit Diffraction
\item Section 4: Reflection, Refraction and Refractive Index
\item Section 5: Lenses, Prisms and Optical Instruments
\item Section 6: Energy of Waves in Free Space and the Inverse Square Law
\item Integration activity
\item Methods and worked examples
\item Exercises
\item Solutions to the exercises
\item Summary sheet
\end{enumerate}
\end{multicols}
\end{sommaire}

\begin{objectifs}
\begin{itemize}
\item By the end of this chapter I can describe an experiment to measure the speed of sound in air and compare it with the speed of light.
\item By the end of this chapter I can state the characteristics of electromagnetic waves and order the regions of the EM spectrum by wavelength and frequency.
\item By the end of this chapter I can explain the production, detection and uses of each region of the EM spectrum, including X-rays.
\item By the end of this chapter I can explain plane polarisation of EM waves and give its applications.
\item By the end of this chapter I can use the grating equation d sinθ = nλ to determine wavelengths and analyse multiple-slit diffraction.
\item By the end of this chapter I can state and apply the laws of reflection and refraction and use Snell's law with refractive index.
\item By the end of this chapter I can explain dispersion, total internal reflection and the critical angle, with applications such as optical fibres.
\item By the end of this chapter I can use the thin lens formula and magnification to predict images formed by a converging lens, and use the dioptre.
\item By the end of this chapter I can describe the passage of light through a prism and the working of a compound microscope and an astronomical telescope.
\item By the end of this chapter I can apply conservation of energy to a point source and use the inverse square law for intensity.
\end{itemize}
\end{objectifs}

\begin{prerequis}
\begin{itemize}
\item Wave basics from Lower Sixth: amplitude, wavelength, frequency, period, v = fλ
\item Trigonometry: sine, tangent, small-angle approximations
\item Reflection of light at plane mirrors and image formation (Form 5 / Lower Sixth)
\item Simple harmonic motion and progressive waves (Lower Sixth)
\item Basic electricity: potential difference, electron acceleration (useful for X-ray production)
\end{itemize}
\end{prerequis}

\newpage

\coursec{Section 1: Measuring the Speed of Sound and the Nature of Electromagnetic Waves}

Every evening in Yaound\'e, CRTV broadcasts reach remote villages, hospital X-ray machines in Douala reveal broken bones without surgery, and a fisherman on the Wouri estuary sees a fish that is not exactly where it appears. All these phenomena --- radio signals, X-rays, the bending of light in water, the rainbow after a tropical downpour --- are governed by one family of waves: \cle{electromagnetic waves}. How can the same physical entity carry a phone call, cook food in a microwave, and let a surgeon see inside the human body? And why does a straw in a glass of palm wine look broken at the surface? This chapter answers these questions by exploring the nature, propagation, and manipulation of electromagnetic waves, with light as our main laboratory.

Before we can handle electromagnetic waves with confidence, we must sharpen our tools on a wave we already know well: \cle{sound}. Measuring the speed of sound in air is a classic GCE practical, and it teaches us the universal wave relation $v = f\lambda$ in a concrete setting. We then contrast sound with light and arrive at Maxwell's beautiful picture of the electromagnetic wave.

\courssub{Recap of Progressive Waves}

\begin{definition}[Progressive wave]
A \cle{progressive wave} is a disturbance that travels through a medium (or through vacuum, for electromagnetic waves), transferring \textbf{energy} from one point to another \emph{without any net transfer of matter}.
\end{definition}

Think of the ``wave'' that travels round a football stadium in Douala: each spectator simply stands up and sits down, yet the disturbance travels right round the stand. The spectators are the medium; only energy and information move.

Three quantities describe any progressive wave:
\begin{itemize}
\item the \cle{wavelength} $\lambda$: the shortest distance between two points vibrating in phase (e.g.\ two successive crests), measured in metres;
\item the \cle{frequency} $f$: the number of complete vibrations per second made by any point of the medium, measured in hertz (\SI{}{Hz});
\item the \cle{speed} $v$: the distance travelled by the disturbance per unit time, in \SI{}{m.s^{-1}}.
\end{itemize}

\begin{propriete}[The wave equation]
For any progressive wave,
\[ v = f\lambda. \]
\end{propriete}

\textbf{Justification.} In one period $T = 1/f$, the source completes exactly one oscillation and the disturbance advances by exactly one wavelength $\lambda$. Hence
\[ v = \frac{\text{distance}}{\text{time}} = \frac{\lambda}{T} = f\lambda. \]

Sound in air is a \cle{longitudinal} wave: the air molecules oscillate \emph{parallel} to the direction of propagation, producing alternating compressions and rarefactions. Light, as we shall see, is \cle{transverse}. The relation $v = f\lambda$ applies to both.

\courssub{Measuring the Speed of Sound in Air}

\textbf{Method 1: The echo method.}

\begin{experience}[The echo (clap--echo) experiment]
Stand at a measured distance $d$ (e.g.\ \SI{100}{m}) from a large flat wall on the school field. Clap hands (or strike two wooden blocks) repeatedly and adjust the rhythm until each clap coincides exactly with the echo of the previous one. A partner times $n = 50$ clap intervals with a stopwatch, giving the total time $t$. The time for one round trip is $t/n$, and the sound travels $2d$ in that time:
\[ v = \frac{2d}{t/n} = \frac{2nd}{t}. \]
\textbf{Sources of error:} reaction time of the timekeeper (reduced by timing many claps), wind, and uncertainty in judging coincidence of clap and echo.
\end{experience}

\textbf{Method 2: The resonance tube.}

This is the precise laboratory method. A long glass tube, open at the top, is connected at its base to a reservoir so that the water level --- and hence the length $l$ of the air column --- can be varied. A vibrating tuning fork of known frequency $f$ is held above the open end. The fork sends sound down the tube; the wave reflects at the water surface (a closed end) and a \cle{stationary wave} forms in the air column. When the column length is just right, the air vibrates strongly and a loud sound is heard: this is \cle{resonance}.

At resonance the air column has a node at the water surface and an antinode just \emph{above} the open end. The distance from the open end to the antinode is the \cle{end correction} $e$. Starting from a high water level and lowering it slowly, the first (shortest) resonance occurs at length $l_1$ and the second at length $l_2$:
\[ l_1 + e = \frac{\lambda}{4}, \qquad l_2 + e = \frac{3\lambda}{4}. \]
Subtracting eliminates the unknown end correction:
\[ l_2 - l_1 = \frac{\lambda}{2} \quad\Longrightarrow\quad \lambda = 2(l_2 - l_1). \]

\begin{methode}[Speed of sound by the resonance tube]
\begin{enumerate}
\item Strike the tuning fork of known frequency $f$ and hold it over the open end of the tube.
\item Lower the water level slowly from the top; note the length $l_1$ of the air column at the \textbf{first} (loudest) resonance.
\item Continue lowering the water; note the length $l_2$ at the \textbf{second} resonance.
\item Compute the wavelength: $\lambda = 2(l_2 - l_1)$. The end correction cancels automatically.
\item Compute the speed: \[ \boxed{\,v = 2f(l_2 - l_1)\,} \]
\item Repeat with forks of different frequencies and average, or plot $v$ values, to improve accuracy. Note the air temperature, as $v$ varies with $\sqrt{T}$.
\end{enumerate}
\end{methode}

\begin{popfigure}
\begin{tikzpicture}[scale=0.9]
% tube
\draw[thick, popink] (0,0) -- (0,7);
\draw[thick, popink] (1.4,0) -- (1.4,7);
% water at first resonance level (for illustration)
\fill[popblueL] (0,0) rectangle (1.4,2.6);
\draw[popblue, thick] (0,2.6) -- (1.4,2.6);
\node[popblue, right] at (1.5,2.6) {water surface (node)};
% air column
\node[popmuted] at (0.7,5.2) {air column};
% reservoir
\draw[thick, popink] (1.4,0.4) -- (2.6,0.4);
\draw[thick, popink] (2.6,0.4) -- (2.6,2.0);
\draw[thick, popink] (2.6,2.0) -- (3.4,2.0);
\draw[thick, popink] (3.4,2.0) -- (3.4,0.0);
\fill[popblueL] (2.6,0.4) rectangle (3.4,1.6);
\node[popmuted, right, align=center] at (3.5,1.0) {movable\\reservoir};
% tuning fork
\draw[thick, poporange] (-0.25,7.6) -- (-0.25,8.4);
\draw[thick, poporange] (-0.25,8.4) arc (180:360:0.25);
\draw[thick, poporange] (0.25,7.6) -- (0.25,8.4);
\draw[thick, poporange] (0,7.6) -- (0,7.15);
\node[poporange, left] at (-0.3,8.0) {tuning fork, $f$};
% l1 marker
\draw[<->, popgreen, thick] (1.7,2.6) -- (1.7,5.6);
\node[popgreen, right] at (1.8,4.1) {$l_1$ (1st resonance)};
\draw[dashed, popgreen] (0,5.6) -- (1.7,5.6);
% l2 marker (second water level at 0.8)
\draw[<->, poppurple, thick] (-0.4,0.8) -- (-0.4,7.0);
\node[poppurple, left, align=right] at (-0.5,3.9) {$l_2$ (2nd\\resonance)};
\draw[dashed, poppurple] (0,0.8) -- (-0.4,0.8);
\node[popmuted, right] at (1.5,0.8) {2nd water level};
% standing wave hint for first resonance (quarter wave)
\draw[popgold, domain=2.6:5.6, samples=40, smooth] plot (\x,{0.35*sin(deg(3.14159*(\x-2.6)/3.0)) + 0.7});
\node[popgold, above] at (4.1,1.05) {$\lambda/4$};
\end{tikzpicture}
\legende{Resonance tube: the water level is adjusted until the air column resonates with the tuning fork. The first and second resonance lengths satisfy $l_2 - l_1 = \lambda/2$. The standing wave shown corresponds to the first resonance (quarter-wave).}
\end{popfigure}

\begin{exoresolu}[Resonance tube calculation]
A tuning fork of frequency $f = \SI{512}{Hz}$ is held over a resonance tube. Resonances are heard when the air column is $l_1 = \SI{16.0}{cm}$ and $l_2 = \SI{49.2}{cm}$ long. Calculate the speed of sound in air.
\tcblower
\textbf{Solution.} Convert lengths to SI units: $l_1 = \SI{0.160}{m}$, $l_2 = \SI{0.492}{m}$.
\[ \lambda = 2(l_2 - l_1) = 2(0.492 - 0.160) = 2 \times 0.332 = \SI{0.664}{m}. \]
\[ v = f\lambda = 512 \times 0.664 \approx \SI{340}{m.s^{-1}}. \]
This agrees well with the accepted value of about \SI{340}{m.s^{-1}} at room temperature.
\end{exoresolu}

\begin{attention}[Common mistakes with the resonance tube]
\begin{itemize}
\item Do \textbf{not} write $\lambda = 4l_1$ and stop there: the antinode lies slightly \emph{above} the open end, so $l_1 + e = \lambda/4$ contains the unknown end correction $e$. Using $l_1$ alone gives a systematically low value of $v$.
\item The clean formula is $v = 2f(l_2 - l_1)$: subtracting the two resonance positions \textbf{eliminates the end correction}.
\item Do not confuse the first resonance ($\lambda/4$) with the second ($3\lambda/4$); the difference is $\lambda/2$, not $\lambda$.
\item Quote $v$ in \SI{}{m.s^{-1}}: convert centimetres to metres before substituting.
\end{itemize}
\end{attention}

\begin{methode}[Exam tip (GCE): describing the echo experiment]
If asked to \emph{describe an experiment to measure the speed of sound by the echo method}, structure your answer in four parts:
\begin{enumerate}
\item \textbf{Labelled diagram}: observer, measuring tape, distance $d$, large reflecting wall.
\item \textbf{Measurements}: distance $d$ with a tape; time $t$ for $n$ clap--echo intervals with a stopwatch, clapping in rhythm so each clap coincides with the previous echo.
\item \textbf{Calculation}: $v = \dfrac{2nd}{t}$ (the sound makes a \emph{return} journey, hence the factor 2).
\item \textbf{Sources of error and precautions}: reaction time (time many intervals), wind along the path, choice of a large flat wall, repeat and average.
\end{enumerate}
\end{methode}

\courssub{Sound versus Light: Two Kinds of Wave}

Why do we bother contrasting sound and light? Because the contrast reveals what is truly remarkable about electromagnetic waves.

\begin{center}
\begin{tabularx}{\linewidth}{|l|X|X|}
\hline
\rowcolor{popblueL} \textbf{Feature} & \textbf{Sound} & \textbf{Light (EM wave)} \\
\hline
Nature & Mechanical wave & Electromagnetic wave \\
\hline
Medium & Needs a material medium (air, water, steel) & Travels through vacuum \\
\hline
Vibration & Longitudinal (parallel to propagation) & Transverse (perpendicular to propagation) \\
\hline
Speed in air/vacuum & $\approx \SI{340}{m.s^{-1}}$ in air & $c = \SI{3.00e8}{m.s^{-1}}$ in vacuum \\
\hline
Demonstration of transverse nature & None (cannot be polarised) & Polarisation \\
\hline
\end{tabularx}
\end{center}

That light travels through the vacuum of space --- we see the Sun and the stars --- while sound does not (explosions in space films are pure fiction!) already tells us light is \emph{not} a mechanical wave. What, then, is ``waving''?

\courssub{Maxwell's Prediction and the Nature of Electromagnetic Waves}

Around 1865, James Clerk Maxwell united the laws of electricity and magnetism into four equations and made a stunning prediction. A changing electric field creates a magnetic field (as a changing current does), and a changing magnetic field creates an electric field (electromagnetic induction). Maxwell realised these two processes can sustain each other: once an electric charge oscillates --- say electrons surging up and down a transmitter aerial --- it launches a self-propagating disturbance in which oscillating electric and magnetic fields continually regenerate one another, travelling outward even through empty space.

\begin{definition}[Electromagnetic wave]
An \cle{electromagnetic wave} is a self-propagating disturbance consisting of oscillating electric and magnetic fields, which requires \textbf{no material medium}. The electric field $\vec{E}$ and the magnetic field $\vec{B}$ vibrate \textbf{perpendicular to each other} and both \textbf{perpendicular to the direction of propagation}: the wave is \cle{transverse}.
\end{definition}

\begin{popfigure}
\begin{tikzpicture}[scale=1.0]
% axes
\draw[->, thick, popink] (0,0) -- (8.6,0) node[right] {propagation $x$};
\draw[->, thick, popink] (0,0) -- (0,2.2) node[above] {$y$};
\draw[->, thick, popink] (0,0) -- (-1.3,-1.3) node[below left] {$z$};
% E field (vertical, in xy plane)
\draw[popblue, very thick, domain=0:8, samples=120, smooth] plot (\x, {1.5*sin(deg(2*3.14159*\x/4))});
\node[popblue] at (6.3,1.9) {$\vec{E}$ (electric field)};
% B field (in xz plane, drawn with depth)
\draw[poporange, very thick, domain=0:8, samples=120, smooth] plot (\x, {-0.75*sin(deg(2*3.14159*\x/4)) - 0.5});
\node[poporange] at (6.6,-1.6) {$\vec{B}$ (magnetic field)};
% vectors at one point
\draw[->, popblue, thick] (2,0) -- (2,1.5);
\draw[->, poporange, thick] (2,0) -- (1.25,-0.75);
\draw[->, popgreen, very thick] (2,0) -- (3.1,0) node[below right, popgreen] {$\vec{v}$};
% labels of perpendicularity
\node[popmuted, align=center] at (4,-2.2) {$\vec{E} \perp \vec{B} \perp$ direction of propagation};
\end{tikzpicture}
\legende{An electromagnetic wave: the electric field $\vec{E}$ and magnetic field $\vec{B}$ oscillate in phase, perpendicular to each other and to the direction of propagation.}
\end{popfigure}

Maxwell's equations also predicted the \emph{speed} of these waves purely from electric and magnetic constants of vacuum:
\[ c = \frac{1}{\sqrt{\varepsilon_0 \mu_0}} \approx \SI{3.00e8}{m.s^{-1}}, \]
exactly the measured speed of light. Maxwell concluded that \textbf{light itself is an electromagnetic wave} --- one of the greatest unifications in the history of physics. Hertz confirmed the prediction experimentally about twenty years later by producing and detecting radio waves.

\begin{propriete}[Speed of electromagnetic waves]
All electromagnetic waves, from radio waves to gamma rays, travel in vacuum at the same speed
\[ c = \SI{3.00e8}{m.s^{-1}}, \]
and obey the universal wave relation
\[ c = f\lambda. \]
Consequently, high-frequency waves have short wavelengths and low-frequency waves have long wavelengths.
\end{propriete}

\begin{exemplebox}[Wavelength of a CRTV radio signal]
A radio station broadcasts at $f = \SI{100}{MHz} = \SI{1.00e8}{Hz}$. Its wavelength in air (essentially vacuum) is
\[ \lambda = \frac{c}{f} = \frac{3.00\times 10^{8}}{1.00\times 10^{8}} = \SI{3.00}{m}. \]
This is why FM aerials are of the order of a metre long: the aerial must be comparable with the wavelength to absorb the wave efficiently.
\end{exemplebox}

\courssub{The Transverse Nature of Light: a Preview of Polarisation}

How do we \emph{know} the oscillations of an EM wave are transverse? The decisive evidence is \cle{polarisation}. A longitudinal wave, like sound, vibrates along the direction of travel: rotating a filter around the beam can never block it. But a transverse wave vibrates \emph{across} the direction of travel, and a polarising filter (such as the lenses of polaroid sunglasses, or the polariser of a calculator display) transmits only the component of $\vec{E}$ aligned with its axis. Rotate a second filter through $90^{\circ}$ and the light is completely extinguished --- impossible for a longitudinal wave. Polarisation therefore proves that light is transverse, exactly as Maxwell's picture requires. We shall study polarisation quantitatively in Section 2.

\begin{aretenir}[Key points of Section 1]
\begin{itemize}
\item Every progressive wave obeys $v = f\lambda$.
\item Echo method: $v = \dfrac{2nd}{t}$; the factor 2 comes from the return journey.
\item Resonance tube: two successive resonances give $v = 2f(l_2 - l_1)$; subtracting eliminates the end correction.
\item Sound is mechanical and longitudinal and needs a medium; light is electromagnetic and transverse and travels through vacuum.
\item An EM wave consists of mutually perpendicular oscillating fields $\vec{E}$ and $\vec{B}$, both perpendicular to the direction of propagation.
\item In vacuum all EM waves travel at $c = \SI{3.00e8}{m.s^{-1}}$ and obey $c = f\lambda$.
\item Polarisation demonstrates the transverse nature of EM waves.
\end{itemize}
\end{aretenir}

\begin{exercice}[Speed of sound by resonance]
A resonance tube is used with a tuning fork of frequency $\SI{480}{Hz}$. The first resonance occurs at $l_1 = \SI{17.2}{cm}$ and the second at $l_2 = \SI{52.4}{cm}$.
\begin{enumerate}
\item Calculate the wavelength of the sound in the tube.
\item Deduce the speed of sound in air.
\item Explain why using $\lambda = 4l_1$ would give a less accurate result.
\end{enumerate}
\end{exercice}

\begin{corrige}[Exercise 1]
\begin{enumerate}
\item $\lambda = 2(l_2 - l_1) = 2(0.524 - 0.172) = 2 \times 0.352 = \SI{0.704}{m}$.
\item $v = f\lambda = 480 \times 0.704 \approx \SI{338}{m.s^{-1}}$.
\item $\lambda = 4l_1$ ignores the end correction $e$: the antinode lies a small distance above the open end, so the true quarter-wavelength is $l_1 + e$, not $l_1$. The result would be systematically too small. Using $(l_2 - l_1)$ cancels $e$ exactly.
\end{enumerate}
\end{corrige}

\begin{savaistu}[From Maxwell to mobile money]
Every mobile phone call, every mobile-money transfer and every radio broadcast in Cameroon relies on Maxwell's 1865 prediction. The oscillating currents in a transmitter mast launch EM waves that regenerate themselves across hundreds of kilometres of empty space and air, until a tiny oscillating current is induced in your phone's aerial. The same equations that describe a CRTV signal also describe the light by which you read this page --- only the frequency differs.
\end{savaistu}

\coursec{Section 2: The Electromagnetic Spectrum, X-rays and Polarisation}

In Section 1 we established that light is only one member of a vast family of \cle{electromagnetic waves} — self-sustaining oscillations of electric and magnetic fields that all travel through vacuum at the same speed $c = \SI{3.00e8}{m.s^{-1}}$ and obey the fundamental relation $c = f\lambda$. We also saw that sound, by contrast, is a mechanical longitudinal wave needing a material medium. In this section we exploit that difference fully: we organise the whole electromagnetic family into the \cle{electromagnetic spectrum}, study how each region is produced, detected and used, examine in detail the production of X-rays in the Coolidge tube, and finish with a property that only transverse waves possess — \cle{polarisation}.

\courssub{The Electromagnetic Spectrum}

\textbf{Ordering the family.} Since $c = f\lambda$ is fixed, the product $f\lambda$ is constant: a wave of long wavelength has a low frequency, and vice versa. Ordering all EM waves by decreasing wavelength (equivalently increasing frequency) gives the seven classical regions:

\begin{center}
\textbf{radio} $\rightarrow$ \textbf{microwave} $\rightarrow$ \textbf{infrared} $\rightarrow$ \textbf{visible} $\rightarrow$ \textbf{ultraviolet} $\rightarrow$ \textbf{X-rays} $\rightarrow$ \textbf{gamma rays}
\end{center}

\begin{popfigure}
\begin{tikzpicture}[x=1cm,y=1cm]
% spectrum bar
\fill[popblue!25]   (0,0) rectangle (2,1);
\fill[popgreen!30]  (2,0) rectangle (4,1);
\fill[poporange!35] (4,0) rectangle (6,1);
\fill[popgold!40]   (6,0) rectangle (7,1);
\fill[poppurple!30] (7,0) rectangle (9,1);
\fill[poppink!35]   (9,0) rectangle (11,1);
\fill[popink!20]    (11,0) rectangle (13,1);
\draw[thick] (0,0) rectangle (13,1);
\node[align=center, font=\footnotesize] at (1,0.5) {Radio};
\node[align=center, font=\footnotesize] at (3,0.5) {Micro- \\ wave};
\node[align=center, font=\footnotesize] at (5,0.5) {Infrared};
\node[align=center, font=\scriptsize] at (6.5,0.5) {Visible};
\node[align=center, font=\footnotesize] at (8,0.5) {UV};
\node[align=center, font=\footnotesize] at (10,0.5) {X-rays};
\node[align=center, font=\footnotesize] at (12,0.5) {$\gamma$-rays};
% wavelength scale (top)
\foreach \x/\lab in {0/{$10^{3}$},2/{$10^{-1}$},4/{$10^{-5}$},6/{$10^{-6}$},7/{$7\times10^{-7}$},9/{$10^{-8}$},11/{$10^{-10}$},13/{$10^{-12}$}}
  \node[above, font=\scriptsize] at (\x,1) {\lab};
\node[above, font=\footnotesize\itshape] at (6.5,1.55) {wavelength $\lambda$ / m (decreasing $\rightarrow$)};
% frequency scale (bottom)
\foreach \x/\lab in {0/{$10^{5}$},2/{$10^{9}$},4/{$10^{13}$},6/{$10^{14}$},7/{$4\times10^{14}$},9/{$10^{16}$},11/{$10^{18}$},13/{$10^{20}$}}
  \node[below, font=\scriptsize] at (\x,0) {\lab};
\node[below, font=\footnotesize\itshape] at (6.5,-0.55) {frequency $f$ / Hz (increasing $\rightarrow$)};
\draw[->, thick, popdark] (0,1.9) -- (13,1.9);
\draw[->, thick, popdark] (0,-0.9) -- (13,-0.9);
\end{tikzpicture}
\legende{The electromagnetic spectrum. Wavelength decreases and frequency increases from radio waves to gamma rays; the visible region is only a narrow band (about $400$--$700\ \mathrm{nm}$).}
\end{popfigure}

\begin{aretenir}[Orders of magnitude to memorise]
\begin{center}
\begin{tabularx}{\linewidth}{|l|X|X|}
\hline
\rowcolor{popblueL} \textbf{Region} & \textbf{Approximate wavelength} & \textbf{Approximate frequency} \\
\hline
Radio waves & $> 0.1\ \mathrm{m}$ & $< 10^{9}\ \mathrm{Hz}$ \\
\hline
Microwaves & $10^{-3}$ to $10^{-1}\ \mathrm{m}$ & $10^{9}$ to $10^{11}\ \mathrm{Hz}$ \\
\hline
Infrared & $7\times10^{-7}$ to $10^{-3}\ \mathrm{m}$ & $10^{11}$ to $4\times10^{14}\ \mathrm{Hz}$ \\
\hline
Visible light & $4\times10^{-7}$ to $7\times10^{-7}\ \mathrm{m}$ & $4\times10^{14}$ to $7.5\times10^{14}\ \mathrm{Hz}$ \\
\hline
Ultraviolet & $10^{-8}$ to $4\times10^{-7}\ \mathrm{m}$ & $7.5\times10^{14}$ to $10^{16}\ \mathrm{Hz}$ \\
\hline
X-rays & $10^{-12}$ to $10^{-8}\ \mathrm{m}$ & $10^{16}$ to $10^{20}\ \mathrm{Hz}$ \\
\hline
Gamma rays & $< 10^{-11}\ \mathrm{m}$ & $> 10^{19}\ \mathrm{Hz}$ \\
\hline
\end{tabularx}
\end{center}
The boundaries overlap slightly; the regions are named according to how the waves are \emph{produced}, not by sharp cut-offs.
\end{aretenir}

\textbf{Photon energy across the spectrum.} Each EM wave can also be pictured as a stream of photons of energy
\[ E = hf = \frac{hc}{\lambda}, \]
where $h = \SI{6.63e-34}{J.s}$ is Planck's constant. Since $f$ increases from radio to gamma, \emph{photon energy increases across the spectrum in the same direction}. Radio photons are far too weak to affect individual molecules; X-ray and gamma photons carry enough energy to \emph{ionise} atoms — this single fact explains all the health hazards discussed below.

\begin{exemplebox}[Photon energy of visible and X-ray radiation]
For green light, $\lambda = 500\ \mathrm{nm}$:
\[ E = \frac{hc}{\lambda} = \frac{6.63\times10^{-34}\times 3.00\times10^{8}}{5.00\times10^{-7}} = \SI{3.98e-19}{J} \approx 2.5\ \mathrm{eV}. \]
For an X-ray of $\lambda = 0.10\ \mathrm{nm}$: $E = \SI{2.0e-15}{J} \approx 12\ \mathrm{keV}$ — about $5000$ times more energetic. Hence X-rays can eject electrons from atoms (ionisation) while visible light cannot.
\end{exemplebox}

\courssub{Production, Detection and Uses of Each Region}

\begin{center}
\begin{tabularx}{\linewidth}{|l|X|X|X|}
\hline
\rowcolor{popblueL} \textbf{Region} & \textbf{Production} & \textbf{Detection} & \textbf{Uses} \\
\hline
Radio & Oscillating charges in transmitter aerials & Receiver aerial + tuned circuit & Radio and TV broadcasting (CRTV), radio astronomy \\
\hline
Microwave & Klystron / magnetron; electronic oscillators & Aerial, diode detector & Mobile phones, satellite links, radar, microwave ovens \\
\hline
Infrared & All warm bodies; heaters, lamps & Skin (warmth), thermistor, IR photodiode, thermal camera & Remote controls, thermal imaging, night vision, physiotherapy \\
\hline
Visible & Hot filaments, discharge tubes, LEDs, the Sun & The eye, photographic film, LDR, photodiode & Vision, photography, illumination, optical fibres \\
\hline
Ultraviolet & The Sun, mercury-vapour lamps, arcs & Fluorescent materials, photographic film, photodiode & Sterilisation, security marking, sun-tan lamps \\
\hline
X-rays & Sudden deceleration of fast electrons (Coolidge tube) & Photographic film, Geiger--M\"uller tube, solid-state detector & Medical radiography, airport security, crystal structure analysis \\
\hline
Gamma & Radioactive decay of unstable nuclei; nuclear reactions & Geiger--M\"uller tube, scintillation counter & Cancer radiotherapy, sterilisation of medical equipment, tracers \\
\hline
\end{tabularx}
\end{center}

A few comments that examiners like to test:
\begin{itemize}
\item \textbf{Radio waves} are produced when electrons oscillate in a metal \cle{aerial}; the same aerial, connected to a tuned circuit, detects them. CRTV broadcasts reach your receiver in Yaound\'e exactly this way.
\item \textbf{Microwave ovens} exploit resonance: microwaves of frequency about $2.45\ \mathrm{GHz}$ are strongly absorbed by water molecules, heating food from within.
\item \textbf{Infrared} is emitted by \emph{every} body above absolute zero; a thermal camera simply maps this emission, which is why rescue teams can locate people in the dark.
\item \textbf{Ultraviolet} causes fluorescence: certain substances absorb UV and re-emit visible light. This is used to reveal security marks on banknotes (including the CFA franc) and to sterilise equipment, since UV kills bacteria.
\end{itemize}

\begin{methode}[Identifying an EM region from its use or detector (exam technique)]
\begin{enumerate}
\item Read the clue: is it a \emph{use} (oven, remote control, radiography), a \emph{detector} (eye, GM tube) or a \emph{source} (aerial, radioactive nucleus)?
\item Match the clue to the region using the table above.
\item If a wavelength or frequency is given, compute the partner quantity with $c = f\lambda$ and locate the region on the spectrum.
\item Always justify your answer in one sentence: ``This is infrared because it is detected as heat by a thermistor.''
\end{enumerate}
\end{methode}

\begin{attention}[Ionising versus non-ionising radiation]
Radio, microwave, infrared, visible and (mostly) ultraviolet photons are \textbf{non-ionising}: they can heat tissue or, for UV, damage skin (sunburn, skin cancer with over-exposure), but they cannot eject electrons from atoms. X-rays and gamma rays are \textbf{ionising}: they break chemical bonds inside living cells and can cause cancer. Hence radiographers stand behind lead screens, and radioactive sources are handled with tongs and stored in lead-lined containers. Never write that ``all EM waves are dangerous'' — precision earns marks.
\end{attention}

\courssub{X-rays: Production in the Coolidge Tube}

\begin{propriete}[Production of X-rays]
\cle{X-rays} are produced when fast-moving electrons are suddenly decelerated by striking a metal target. Part of the kinetic energy lost by the electrons is converted into electromagnetic radiation of very short wavelength (most of the rest heats the target). \emph{(The description of the tube is examinable; learn it with the diagram.)}
\end{propriete}

\begin{popfigure}
\begin{tikzpicture}[x=1cm,y=1cm, >=latex]
% glass envelope
\draw[thick, popdark, rounded corners=6pt] (0,-1.6) rectangle (10,1.6);
\node[font=\footnotesize\itshape, popmuted] at (5,1.95) {evacuated glass envelope (vacuum)};
% cathode
\draw[thick, fill=popink!70] (0.6,-0.5) rectangle (1.1,0.5);
\node[below, font=\footnotesize, align=center] at (0.85,-0.6) {heated cathode \\ (filament)};
% filament supply
\draw[thick] (0.85,-0.5) -- (0.85,-2.2);
\draw[thick] (0.85,-2.2) -- (2.2,-2.2);
\draw (2.2,-2.2) to[battery1, l_={$U_f$}] (3.4,-2.2);
\node[below, font=\scriptsize] at (2.8,-2.5) {filament supply};
% electron beam
\draw[->, very thick, popblue] (1.3,0) -- (7.2,0) node[midway, above, font=\footnotesize] {electron beam};
% target / anode
\draw[thick, fill=poporange!60] (7.3,-0.9) -- (7.9,-0.9) -- (7.9,0.9) -- (7.3,0.9) -- (7.3,0.35) -- (7.55,0.35) -- (7.55,-0.35) -- (7.3,-0.35) -- cycle;
\node[above, font=\footnotesize, align=center] at (7.6,1.0) {tungsten target \\ (anode)};
% cooling fins
\draw[thick, popdark] (7.9,0) -- (9.4,0);
\foreach \y in {-0.5,-0.25,0,0.25,0.5} \draw[thick, popdark] (9.0,\y) -- (9.4,\y);
\node[right, font=\scriptsize, align=left] at (9.45,0) {cooling \\ fins};
% X-rays
\foreach \a in {25,45,65} \draw[->, thick, poppink] (7.55,0.35) -- ++(\a:1.9);
\node[font=\footnotesize, poppink] at (8.1,2.3) {X-rays};
% high voltage
\draw[thick] (0.85,0.5) -- (0.85,2.6) -- (3.6,2.6);
\draw[thick] (7.6,0.9) -- (7.6,2.6) -- (6.0,2.6);
\draw (3.6,2.6) to[sV, l={$V$}] (6.0,2.6);
\node[above, font=\scriptsize] at (4.8,2.75) {high accelerating p.d.\ (tens of kV)};
\node[font=\scriptsize] at (3.4,2.35) {$-$};
\node[font=\scriptsize] at (6.2,2.35) {$+$};
\end{tikzpicture}
\legende{The Coolidge (X-ray) tube. Electrons emitted by the heated cathode are accelerated through a large p.d.\ and decelerated by the tungsten target, emitting X-rays.}
\end{popfigure}

\textbf{How the tube works, step by step.}
\begin{enumerate}
\item A low-voltage supply heats the tungsten \cle{filament} (cathode), which emits electrons by \emph{thermionic emission}.
\item A very high p.d.\ (typically $20$--$100\ \mathrm{kV}$) between cathode and anode accelerates the electrons to high speed. Each electron gains kinetic energy $E_k = eV$.
\item The electrons strike a \cle{target} of a metal with a high melting point (tungsten) and are suddenly decelerated; a small fraction of their kinetic energy becomes X-ray photons (continuous spectrum / bremsstrahlung).
\item The tube is \emph{evacuated} so that electrons are not scattered by air molecules.
\item Because about $99\%$ of the energy appears as heat, the anode is fitted with cooling fins (or rotated and oil-cooled) to prevent melting.
\end{enumerate}

The maximum photon energy corresponds to an electron losing \emph{all} its kinetic energy in one collision (Duane-Hunt law):
\[ eV = hf_{\max} = \frac{hc}{\lambda_{\min}} \quad\Longrightarrow\quad \lambda_{\min} = \frac{hc}{eV}. \]

\begin{savaistu}[Characteristic X-rays]
In addition to the continuous spectrum, if the accelerating voltage is high enough to eject inner-shell electrons from the target atoms, the subsequent de-excitation of outer electrons produces sharp \emph{characteristic X-ray lines} (e.g.\ $K_\alpha$, $K_\beta$) whose wavelengths depend only on the target material, not on the voltage. These are vital for X-ray fluorescence analysis.
\end{savaistu}

\begin{exoresolu}[Minimum wavelength from an X-ray tube]
An X-ray tube operates at an accelerating p.d.\ of $V = 50\ \mathrm{kV}$. Calculate (a) the kinetic energy of an electron reaching the target, (b) the minimum wavelength of the X-rays produced.
\tcblower
(a) Kinetic energy gained by one electron:
\[ E_k = eV = 1.60\times10^{-19}\times 5.0\times10^{4} = \SI{8.0e-15}{J}. \]
(b) The most energetic photon is produced when all of $E_k$ becomes one photon:
\[ \lambda_{\min} = \frac{hc}{eV} = \frac{6.63\times10^{-34}\times 3.00\times10^{8}}{8.0\times10^{-15}} = \SI{2.5e-11}{m}. \]
This lies in the X-ray region ($< 10^{-8}\ \mathrm{m}$), as expected.
\end{exoresolu}

\textbf{Properties and uses of X-rays.}
\begin{itemize}
\item They are EM waves of very short wavelength: they travel in straight lines, are not deflected by electric or magnetic fields (they carry no charge), and ionise gases.
\item They \emph{penetrate} matter, but are absorbed more by dense material: \textbf{bone absorbs X-rays far more than soft tissue}, which is why bones appear white on a radiograph.
\item Uses: \textbf{medical radiography} (fractures, chest images), \textbf{airport security scanners} for luggage, and \textbf{X-ray diffraction by crystals}, which reveals the spacing of atomic planes.
\end{itemize}

\begin{attention}[X-rays $\neq$ gamma rays]
A classic confusion. \textbf{X-rays} are produced \emph{outside the nucleus}, by the sudden deceleration of fast electrons (or by electronic transitions in inner atomic shells). \textbf{Gamma rays} are emitted by \emph{unstable nuclei} during radioactive decay, e.g.\ \ensuremath{\mathrm{^{60}_{27}Co \rightarrow  ^{60}_{28}Ni + ^{0}_{-1}e + \gamma}}. The two may have identical wavelengths — they differ by \emph{origin}, not by nature. Saying ``X-rays come from radioactive decay'' loses the mark.
\end{attention}

\courssub{Plane Polarised Electromagnetic Waves}

\textbf{Intuition.} A transverse wave on a rope can oscillate vertically, horizontally, or in any direction perpendicular to the rope. Ordinary light from a lamp or the Sun is a mixture of waves oscillating in \emph{all} transverse directions at random: it is \cle{unpolarised}. If we pass it through a filter that transmits only one direction of oscillation, the emerging wave oscillates in a single plane.

\begin{definition}[Plane polarised wave]
A \cle{plane polarised wave} is a transverse wave whose oscillations occur in \textbf{one plane only} — the plane containing the direction of propagation and the (single) direction of vibration. For an EM wave, it is the \emph{electric field vector} $\vec{E}$ that is confined to this plane.
\end{definition}

\begin{attention}[Only transverse waves can be polarised]
Polarisation is the decisive experimental proof that light is a \emph{transverse} wave. Sound in air is \emph{longitudinal}: the oscillations are already along the direction of travel, so there is no transverse direction to select. \textbf{Sound cannot be polarised.} Writing ``polarised sound'' in the GCE is an automatic error.
\end{attention}

\begin{popfigure}
\begin{tikzpicture}[x=1cm,y=1cm, >=latex]
% unpolarised light
\draw[->, thick] (-0.5,0) -- (1.2,0);
\foreach \a in {20,65,110,155} \draw[<->, popblue, thick] (0.35,0) ++(\a:-0.45) -- ++(\a:0.9);
\node[font=\footnotesize, align=center] at (0.35,-0.95) {unpolarised \\ light};
% first polaroid
\draw[thick, fill=popgreen!25] (1.6,-1.0) rectangle (1.9,1.0);
\foreach \y in {-0.8,-0.4,0,0.4,0.8} \draw[popgreen!60!black] (1.62,\y) -- (1.88,\y);
\node[font=\footnotesize, align=center] at (1.75,1.35) {polariser \\ (vertical axis)};
% polarised light between
\draw[->, thick] (1.9,0) -- (4.6,0);
\draw[<->, poporange, very thick] (2.6,-0.55) -- (2.6,0.55);
\draw[<->, poporange, very thick] (3.4,-0.55) -- (3.4,0.55);
\draw[<->, poporange, very thick] (4.2,-0.55) -- (4.2,0.55);
\node[font=\footnotesize, align=center] at (3.3,-0.95) {plane polarised \\ (vertical)};
% second polaroid crossed
\draw[thick, fill=poppink!30] (4.6,-1.0) rectangle (4.9,1.0);
\foreach \x in {4.68,4.82} \draw[poppink!70!black] (\x,-0.9) -- (\x,0.9);
\foreach \y in {-0.6,-0.2,0.2,0.6} \draw[poppink!70!black] (4.62,\y) -- (4.88,\y);
\node[font=\footnotesize, align=center] at (4.75,1.35) {analyser \\ (horizontal axis)};
% no light after
\draw[->, thick, dashed, popmuted] (4.9,0) -- (6.3,0);
\node[font=\footnotesize, align=center, popdark] at (5.7,-0.95) {no light \\ transmitted};
\end{tikzpicture}
\legende{Unpolarised light passes through a polaroid (the polariser); the emerging plane polarised light is completely blocked by a second polaroid (the analyser) whose transmission axis is crossed at $90^{\circ}$.}
\end{popfigure}

\textbf{Malus-type behaviour (qualitative).} Let plane polarised light of intensity $I_0$ fall on an analyser whose transmission axis makes an angle $\theta$ with the plane of polarisation. Only the component of $\vec{E}$ along the axis is transmitted, and since intensity is proportional to the \emph{square} of the amplitude,
\[ I = I_0\cos^2\theta. \]
Consequences to remember: $\theta = 0^{\circ}$ gives full transmission; $\theta = 90^{\circ}$ (crossed polaroids) gives \emph{darkness}; rotating the analyser through $360^{\circ}$ produces two maxima and two extinctions. Note also that when \emph{unpolarised} light passes through a single polaroid, the transmitted intensity is exactly half, whatever the orientation.

\begin{exemplebox}[Intensity through two polaroids]
Unpolarised light of intensity $I_i$ falls on a polariser, then on an analyser at $60^{\circ}$ to the polariser. After the polariser: $I_1 = \tfrac{1}{2}I_i$. After the analyser:
\[ I_2 = I_1\cos^2 60^{\circ} = \tfrac{1}{2}I_i \times (0.5)^2 = 0.125\,I_i. \]
Only one eighth of the original intensity emerges.
\end{exemplebox}

\textbf{Applications of polarisation.}
\begin{itemize}
\item \textbf{Polaroid sunglasses}: light reflected by a road or by water is partially polarised horizontally; the vertical transmission axis of the lenses blocks this \emph{glare} while transmitting useful light — a real comfort when driving on the Douala--Yaound\'e highway in the sun.
\item \textbf{Stress analysis}: a transparent plastic model placed between crossed polaroids shows coloured fringes where internal stresses are greatest; engineers use this to test bridges and machine parts before construction.
\item \textbf{LCD screens}: a liquid-crystal display consists of a liquid-crystal layer between two crossed polaroids; applying a voltage rotates the plane of polarisation locally, switching each pixel between bright and dark.
\end{itemize}

\begin{savaistu}[Polarisation in everyday life]
Bees navigate using the polarisation pattern of skylight, even when the Sun is hidden by clouds. Photographers, on the other hand, screw a polarising filter onto their lenses to darken blue skies and remove reflections from shop windows — the same physics as your sunglasses.
\end{savaistu}

\begin{aretenir}[GCE exam tip: one production, one detector, one use per region]
Examiners regularly ask: ``Name the region of the EM spectrum which \dots'' or ``State one use and one detector of \dots''. Memorise for \emph{each} region one production mechanism, one detector and one use (the summary table above is enough). For X-rays, be able to draw and label the Coolidge tube and to state that X-rays arise from the sudden deceleration of fast electrons. For polarisation, remember the three keywords: \emph{transverse only}, \emph{one plane}, $\emph{I} = I_0\cos^2\theta$.
\end{aretenir}

\begin{exercice}[Regions, X-rays and polarisation]
\begin{enumerate}
\item A wave has frequency $f = 6.0\times10^{14}\ \mathrm{Hz}$. Calculate its wavelength in vacuum and identify the region of the spectrum.
\item An X-ray tube operates at $30\ \mathrm{kV}$. Calculate the maximum frequency of the X-rays produced.
\item Explain why the target of an X-ray tube is made of tungsten and why the tube must be evacuated.
\item Plane polarised light of intensity $I_0$ passes through an analyser. At what angle $\theta$ is the transmitted intensity equal to $\tfrac{1}{4}I_0$?
\item State, with a reason, whether sound waves in air can be polarised.
\end{enumerate}
\end{exercice}

\begin{corrige}[Exercise 1]
\begin{enumerate}
\item $\lambda = c/f = (3.00\times10^{8})/(6.0\times10^{14}) = \SI{5.0e-7}{m} = 500\ \mathrm{nm}$: visible light (green).
\item $f_{\max} = eV/h = (1.60\times10^{-19}\times 3.0\times10^{4})/(6.63\times10^{-34}) = \SI{7.2e18}{Hz}$.
\item Tungsten has a very high melting point, so it resists the intense heating produced when $99\%$ of the electron energy becomes heat. The vacuum prevents the electrons from colliding with air molecules, which would scatter them and reduce their energy before they reach the target.
\item $\cos^2\theta = 1/4 \Rightarrow \cos\theta = 1/2 \Rightarrow \theta = 60^{\circ}$.
\item No. Sound in air is longitudinal: its oscillations are parallel to the direction of propagation, so there is no transverse plane to select. Only transverse waves can be polarised.
\end{enumerate}
\end{corrige}

\coursec{Section 3: Light Sources, Diffraction Gratings and Multiple-Slit Diffraction}

In Section 2, we explored the vast electromagnetic spectrum and the production of X-rays. We also introduced the concept of polarisation, a property exclusive to transverse waves. Now, we turn our attention to the \emph{visible} region of the spectrum — light. To understand how we analyse light, we must first understand where it comes from and how it behaves when it encounters a periodic structure with thousands of slits: the diffraction grating. This instrument is the workhorse of spectroscopy, allowing us to identify the chemical composition of stars, the gases in a neon sign, or the wavelength of a laser with remarkable precision.

\courssub{Light Sources and the Nature of Spectra}

Before studying diffraction, we must characterise the light sources we use in the laboratory and in nature.

\begin{definition}[Light Sources and Spectra]
A \cle{light source} is any object that emits electromagnetic radiation in the visible range ($\approx 400$--$700\ \mathrm{nm}$). We classify them by their spectral output:
\begin{itemize}
    \item \textbf{Natural sources:} The Sun (approximately a black body at $5800\ \mathrm{K}$), stars, fireflies, lightning.
    \item \textbf{Artificial sources:} Incandescent filaments (tungsten lamps), gas discharge tubes (sodium, mercury, neon, hydrogen), Light Emitting Diodes (LEDs), Lasers.
\end{itemize}
The \cle{spectrum} of a source is the distribution of intensity as a function of wavelength (or frequency).
\begin{itemize}
    \item \textbf{Continuous spectrum:} Contains all wavelengths over a range with no gaps. Produced by hot, dense solids, liquids, or high-pressure gases (e.g., incandescent bulb filament, the Sun's photosphere).
    \item \textbf{Line spectrum (Emission spectrum):} Consists of bright, sharp lines at specific wavelengths separated by dark gaps. Produced by low-pressure gases in discharge tubes (e.g., sodium vapour lamp, neon sign). Each element has a unique "fingerprint" of lines.
    \item \textbf{Absorption spectrum:} A continuous spectrum crossed by dark lines at specific wavelengths. Produced when white light passes through a cooler, low-pressure gas (e.g., Fraunhofer lines in the solar spectrum).
\end{itemize}
\end{definition}

\begin{savaistu}[Cameroon Context]
In the physics laboratories of our lycées (e.g., Lycée Leclerc Yaoundé, GBHS Bamenda), the standard sources for diffraction experiments are the \textbf{sodium vapour lamp} (quasi-monochromatic yellow doublet at $589.0$ and $589.6\ \mathrm{nm}$) and the \textbf{mercury vapour lamp} (distinct lines: violet $435.8\ \mathrm{nm}$, green $546.1\ \mathrm{nm}$, yellow $577.0/579.1\ \mathrm{nm}$). He-Ne lasers ($\lambda = 632.8\ \mathrm{nm}$) are increasingly common for demonstrating sharp diffraction orders.
\end{savaistu}

\courssub{Coherence: The Prerequisite for Interference and Diffraction}

You know from double-slit interference that a stable pattern requires \cle{coherent sources}. Two sources are coherent if they maintain a \textbf{constant phase difference} and have the \textbf{same frequency}.

\begin{aretenir}[Conditions for Observable Interference/Diffraction Patterns]
\begin{enumerate}
    \item \textbf{Monochromaticity (Same frequency):} $\Delta f \approx 0$. White light produces overlapping patterns of different colours; only the central maximum ($n=0$) remains white.
    \item \textbf{Constant phase relationship:} The phase difference $\Delta \phi$ does not vary randomly with time.
    \item \textbf{Same polarisation state:} For maximum contrast, the electric field vectors must oscillate in parallel planes.
\end{enumerate}
In practice, we derive coherence from a \textbf{single primary source} (e.g., illuminating a grating with light from one sodium lamp). The different slits of the grating then act as secondary coherent sources (Huygens' principle).
\end{aretenir}

\courssub{The Optical Transmission Grating}

A transmission grating consists of a large number $N$ of equally spaced, narrow, parallel slits (or transparent lines) ruled on a glass or plastic substrate.

\begin{popfigure}
\begin{tikzpicture}[scale=1.2, >=stealth]
% Grating substrate
\fill[popblueL!30] (-3, -1.5) rectangle (3, 1.5);
\draw[thick, popdark] (-3, -1.5) -- (3, -1.5) -- (3, 1.5) -- (-3, 1.5) -- cycle;
% Slits (representative few)
\foreach \y in {-1.2, -0.8, -0.4, 0, 0.4, 0.8, 1.2} {
    \fill[white] (-3, \y - 0.08) rectangle (3, \y + 0.08);
    \draw[popline] (-3, \y - 0.08) -- (3, \y - 0.08);
    \draw[popline] (-3, \y + 0.08) -- (3, \y + 0.08);
}
% Label N lines per metre
\node[align=center, text width=3cm, popdark] at (0, -2) {Grating: $N$ lines per metre \\ Grating spacing $d = 1/N$};
% Incident beam
\draw[->, thick, poporange] (-4.5, 0) -- (-3, 0) node[midway, above] {Normal incidence};
% Diffracted orders
\draw[->, thick, popgreen] (3, 0) -- (5, 0) node[midway, above] {$n=0$};
\draw[->, thick, popgreen] (3, 0) -- (5, 1.2) node[midway, above right] {$n=+1$};
\draw[->, thick, popgreen] (3, 0) -- (5, -1.2) node[midway, below right] {$n=-1$};
\draw[->, thick, popgreen!70!black] (3, 0) -- (5, 2.2) node[midway, above right] {$n=+2$};
\draw[->, thick, popgreen!70!black] (3, 0) -- (5, -2.2) node[midway, below right] {$n=-2$};
% Angle theta for n=1
\draw[poppurple, thick] (3, 0) -- (5, 1.2);
\draw[poppurple, <->] (3.3, 0) arc (0:31:0.3);
\node[poppurple] at (3.7, 0.2) {$\theta_1$};
% Normal line
\draw[dashed, popmuted] (3, -2.5) -- (3, 2.5);
\node[popmuted, right] at (3, 2.5) {Normal};
\end{tikzpicture}
\legende{Transmission grating at normal incidence. Parallel rays incident normally are diffracted into discrete orders $n = 0, \pm 1, \pm 2, \dots$ at angles $\theta_n$ satisfying $d \sin\theta_n = n\lambda$.}
\end{popfigure}

\begin{definition}[Grating Spacing]
If the grating has \cle{$N$ lines per unit length} (usually given in lines per millimetre or lines per metre), the distance between the centres of two adjacent slits — the \cle{grating spacing} or \cle{grating element} — is:
\[
d = \frac{1}{N}
\]
\textbf{SI Unit:} metre ($\mathrm{m}$). If $N$ is in lines/mm, convert to lines/m first: $N_{\mathrm{m}} = N_{\mathrm{mm}} \times 10^3$.
\end{definition}

\begin{attention}[Classic GCE Trap: Units of $N$ and $d$]
\emph{The syllabus and examiners frequently test unit conversion.}
\begin{itemize}
    \item A grating marked "$500\ \mathrm{lines/mm}$" means $N = 500 \times 10^3 = 5 \times 10^5\ \mathrm{lines/m}$.
    \item Then $d = \frac{1}{5 \times 10^5} = 2 \times 10^{-6}\ \mathrm{m} = 2000\ \mathrm{nm}$.
    \item \textbf{Never} plug "$500$" directly into $d = 1/N$ expecting metres. Always convert $N$ to $\mathrm{lines/m}$ first.
\end{itemize}
\end{attention}

\courssub{Derivation of the Grating Equation (Normal Incidence)}

Consider a plane wave of wavelength $\lambda$ incident normally on the grating. Each slit acts as a source of Huygens wavelets. We look for directions $\theta$ where wavelets from \emph{all} slits interfere constructively.

\begin{popfigure}
\begin{tikzpicture}[scale=1.5, >=stealth]
% Two adjacent slits
\draw[thick, popdark] (-1, -0.3) -- (-1, 0.3); % Slit A
\draw[thick, popdark] (0, -0.3) -- (0, 0.3);   % Slit B
\node[left] at (-1, 0) {$A$};
\node[right] at (0, 0) {$B$};
% Wavefronts / Rays
\draw[->, thick, poporange] (-2, 0.15) -- (-1, 0.15);
\draw[->, thick, poporange] (-2, -0.15) -- (-1, -0.15);
% Diffracted rays at angle theta
\draw[->, thick, popgreen] (-1, 0.15) -- (1.5, 1.2);
\draw[->, thick, popgreen] (0, -0.15) -- (1.5, 0.05); % shifted for clarity
% Path difference construction
\draw[dashed, popmuted] (-1, 0.15) -- (0, 0.15); % horizontal from A
\draw[dashed, popmuted] (0, -0.15) -- (0, 0.15); % vertical at B
% Right angle mark
\draw[popmuted] (0, 0.15) -- ++(-0.1, 0) -- ++(0, -0.1) -- ++(0.1, 0);
% Labels
\node[above] at (-0.5, 0.18) {$d$};
\node[right, poppurple] at (0, 0.0) {$\Delta x = d\sin\theta$};
\draw[<->, poppurple] (0.05, -0.15) -- (0.05, 0.15);
\node[right] at (0.05, 0) {$d\sin\theta$};
% Angle theta
\draw[poppurple, <->] (-0.8, 0) arc (0:36.87:0.8);
\node[poppurple, above left] at (-0.5, 0.3) {$\theta$};
% Normal
\draw[dashed, popmuted] (-1, -0.5) -- (-1, 0.5);
\end{tikzpicture}
\legende{Geometry of path difference between adjacent slits $A$ and $B$ separated by $d$. For constructive interference, the path difference $\Delta x = d \sin\theta$ must be an integer multiple of $\lambda$.}
\end{popfigure}

From the right-angled triangle in the figure, the path difference between wavelets from adjacent slits $A$ and $B$ reaching a distant screen at angle $\theta$ is:
\[
\Delta x = d \sin\theta
\]
For a \cle{principal maximum} (bright fringe), the waves from \emph{every} slit must arrive in phase. This requires the path difference between \emph{any} two adjacent slits to be an integer multiple of the wavelength:
\[
d \sin\theta = n\lambda \qquad (n = 0, \pm 1, \pm 2, \dots)
\]

\begin{propriete}[Grating Equation at Normal Incidence]
For a transmission grating with spacing $d$, illuminated normally by monochromatic light of wavelength $\lambda$, principal maxima (bright spectral lines) occur at angles $\theta_n$ given by:
\[
\[\boxed{d \sin\theta_n = n\lambda}\]
\]
where $n$ is the \cle{order of the spectrum} ($n=0$ is the central/zeroth order, $n=\pm 1$ first order, etc.).
\textbf{Proof:} Examinable. Derive using path difference $d\sin\theta$ between adjacent slits and the condition for constructive interference ($\Delta x = n\lambda$).
\end{propriete}

\begin{aretenir}[Key Features of the Grating Equation]
\begin{itemize}
    \item \textbf{Angular separation:} Unlike a prism, the grating separates wavelengths \emph{linearly} in $\sin\theta$ ($\sin\theta \propto \lambda$). Longer wavelengths (red) are diffracted \emph{more} ($\theta$ larger).
    \item \textbf{Multiple orders:} Each wavelength appears in several orders ($n=0, 1, 2\dots$) at different angles.
    \item \textbf{Maximum order:} Since $|\sin\theta| \le 1$, the highest observable order $n_{\max}$ for a given $\lambda$ is the largest integer satisfying $n \le d/\lambda$.
\end{itemize}
\end{aretenir}

\begin{propriete}[Maximum Observable Order]
The maximum order number $n_{\max}$ is the integer part of $d/\lambda$:
\[
n_{\max} = \left\lfloor \frac{d}{\lambda} \right\rfloor
\]
If $d/\lambda = 3.7$, orders $n = 0, \pm 1, \pm 2, \pm 3$ are visible; $n=4$ would require $\sin\theta > 1$ (impossible).
\end{propriete}

\courssub{Measuring Wavelength with a Diffraction Grating}

This is a standard practical experiment (Paper 3) and a frequent theory question.

\begin{methode}[Measuring $\lambda$ using a Transmission Grating (Spectrometer Method)]
\begin{enumerate}
    \item \textbf{Setup:} Collimator $\to$ Grating (normal to axis) $\to$ Telescope. Use a sodium lamp or laser.
    \item \textbf{Adjust for normal incidence:} Rotate grating until the reflected image of the slit coincides with the direct view (or use the telescope to view the direct beam $n=0$, then rotate grating by $45^\circ$ and set telescope to $90^\circ$ — standard spectrometer adjustment). \emph{Simpler school method:} Place grating normal to the bench, shine laser/beam normally, measure $\theta$ on a screen at distance $D$.
    \item \textbf{Measure angles:} For a chosen order $n$ (usually $n=1$ or $2$), measure the angular position $\theta_n$ on \textbf{both} sides (left and right). Calculate mean $\theta_n$ to eliminate alignment errors.
    \item \textbf{Determine $d$:} Read $N$ (lines/mm) from grating label. Compute $d = 1/(N \times 10^3)\ \mathrm{m}$.
    \item \textbf{Calculate $\lambda$:} $\lambda = \dfrac{d \sin\theta_n}{n}$.
    \item \textbf{Repeat:} Use different orders $n$ and average results.
\end{enumerate}
\end{methode}

\begin{exoresolu}[Measuring the Wavelength of a Laser]
{A He-Ne laser beam is incident normally on a diffraction grating marked "$300\ \mathrm{lines/mm}$". The first-order maximum ($n=1$) is observed on a screen $2.00\ \mathrm{m}$ away, at a distance $y = 0.420\ \mathrm{m}$ from the central spot. Determine the wavelength of the laser light.}
\tcblower
\textbf{Solution:}
\begin{enumerate}
    \item \textbf{Grating spacing $d$:}
    $N = 300\ \mathrm{lines/mm} = 300 \times 10^3 = 3.00 \times 10^5\ \mathrm{lines/m}$.
    $d = \dfrac{1}{N} = \dfrac{1}{3.00 \times 10^5} = 3.33 \times 10^{-6}\ \mathrm{m}$.
    \item \textbf{Angle $\theta_1$:}
    $\tan\theta_1 = \dfrac{y}{D} = \dfrac{0.420}{2.00} = 0.210$.
    $\theta_1 = \tan^{-1}(0.210) = 11.86^\circ$.
    \item \textbf{Wavelength $\lambda$ (using $n=1$):}
    $\lambda = d \sin\theta_1 = (3.33 \times 10^{-6}) \times \sin(11.86^\circ)$.
    $\sin(11.86^\circ) \approx 0.205$.
    $\lambda = 3.33 \times 10^{-6} \times 0.205 = 6.83 \times 10^{-7}\ \mathrm{m} = \mathbf{683\ nm}$.
    \item \textbf{Comment:} The accepted value for He-Ne is $632.8\ \mathrm{nm}$. The discrepancy suggests experimental error in $y$ or $D$, or the grating marking is approximate. In an exam, use your calculated values consistently.
\end{enumerate}
\end{exoresolu}

\courssub{Multiple-Slit Diffraction: From Two Slits to $N$ Slits}

We know the double-slit pattern: broad interference fringes modulated by the single-slit diffraction envelope. What happens as we increase the number of slits $N$ to $10$, $100$, $10^5$?

\begin{popfigure}
\begin{tikzpicture}
\begin{axis}[
    width=\linewidth, height=7cm,
    xlabel={$\sin\theta$ (or position on screen)}, ylabel={Relative Intensity $I/I_0$},
    xmin=-1.5, xmax=1.5, ymin=0, ymax=1.1,
    xtick={-1, 0, 1}, xticklabels={$-1$, $0$, $+1$},
    ytick={0, 0.5, 1}, legend pos=north east,
    grid=major, grid style={popline!50},
    samples=501, domain=-1.5:1.5,
    declare function={
        single_slit(\x) = (abs(\x) < 0.0001) ? 1 : (sin(deg(5*pi*\x))/(5*pi*\x))^2;
        multi_slit(\N,\x) = (abs(sin(deg(pi*\x))) < 0.0001) ? \N*\N : (sin(deg(\N*pi*\x))/(\N*sin(deg(pi*\x))))^2;
    }
]
% Double slit (N=2)
\addplot[popblue, thick] {single_slit(x) * multi_slit(2,x)};
\addlegendentry{$N=2$ (Double slit)}

% N=4
\addplot[poporange, thick] {single_slit(x) * multi_slit(4,x)};
\addlegendentry{$N=4$}

% N=10 (Grating)
\addplot[popgreen, thick] {single_slit(x) * multi_slit(10,x)};
\addlegendentry{$N=10$ (Grating)}

% Vertical lines for orders
\draw[dashed, popmuted] (axis cs:-1,0) -- (axis cs:-1,1.1);
\draw[dashed, popmuted] (axis cs:0,0) -- (axis cs:0,1.1);
\draw[dashed, popmuted] (axis cs:1,0) -- (axis cs:1,1.1);
\node[popmuted, anchor=north] at (axis cs:-1,1.1) {$n=-1$};
\node[popmuted, anchor=north] at (axis cs:0,1.1) {$n=0$};
\node[popmuted, anchor=north] at (axis cs:1,1.1) {$n=+1$};
\end{axis}
\end{tikzpicture}
\legende{Intensity distribution for multiple slits (normalised). As $N$ increases, principal maxima (at integer $n$) become extremely sharp and intense ($\propto N^2$), while secondary maxima become negligible. The single-slit envelope (dotted) modulates the overall intensity.}
\end{popfigure}

\begin{propriete}[Effect of Increasing Number of Slits $N$]
For a grating with $N$ illuminated slits:
\begin{enumerate}
    \item \textbf{Principal Maxima:} Occur at the same angles $d\sin\theta = n\lambda$. Their \textbf{intensity} is proportional to $N^2$ (amplitude $\propto N$). Their \textbf{angular width} (between first minima) is proportional to $1/N$. They become extremely sharp lines.
    \item \textbf{Minima:} There are $(N-1)$ minima between adjacent principal maxima.
    \item \textbf{Secondary Maxima:} There are $(N-2)$ weak secondary maxima between principal maxima. Their intensity $\propto 1/N^2$ relative to principal maxima $\to$ negligible for large $N$.
    \item \textbf{Resolving Power:} The sharpness allows separation of close wavelengths. Resolving power $R = \lambda/\Delta\lambda = nN$. (Examinable definition).
\end{enumerate}
\end{propriete}

\begin{aretenir}[Why Gratings Give Sharp Lines]
A grating is essentially a "multi-slit interferometer". Constructive interference requires the path difference to be \emph{exactly} $n\lambda$ for \emph{every} pair of adjacent slits. A tiny deviation $\Delta\theta$ from the peak angle introduces a phase shift $\Delta\phi = \frac{2\pi}{\lambda}d\cos\theta\,\Delta\theta$ between neighbours. After $N$ slits, the total phase spread is $N\Delta\phi$. The first minimum occurs when $N\Delta\phi = 2\pi$, giving angular width $\Delta\theta \approx \lambda/(Nd\cos\theta)$. \textbf{Width $\propto 1/N$.}
\end{aretenir}

\courssub{White Light and the Grating Spectrum}

When white light (continuous spectrum) illuminates a grating, each wavelength $\lambda$ is diffracted at its own angle $\theta_n(\lambda) = \sin^{-1}(n\lambda/d)$.

\begin{popfigure}
\begin{tikzpicture}[scale=1.2, >=stealth]
% Grating
\fill[popblueL!30] (-0.3, -2) rectangle (0.3, 2);
\draw[thick, popdark] (-0.3, -2) -- (0.3, -2) -- (0.3, 2) -- (-0.3, 2) -- cycle;
% White incident beam
\draw[->, thick, white, double=popdark, double distance=2pt] (-2, 0) -- (-0.3, 0) node[midway, above] {White light};
% Diffracted orders - Central (n=0)
\draw[->, thick, white, double=popdark, double distance=2pt] (0.3, 0) -- (3, 0) node[midway, above] {$n=0$ (White)};
% n=1 Spectrum (Violet to Red)
\foreach \wavelength/\col in {400/violet, 450/blue, 530/green, 590/yellow, 650/orange, 700/red} {
    \pgfmathsetmacro{\ypos}{1.6*\wavelength/700}
    \draw[->, thick, \col] (0.3, 0) -- (3, \ypos);
}
\node[align=left, text width=2.5cm, popdark] at (3.5, 0.7) {$n=+1$ Spectrum:\\ Violet $\to$ Red};
% n=-1 Spectrum
\foreach \wavelength/\col in {400/violet, 450/blue, 530/green, 590/yellow, 650/orange, 700/red} {
    \pgfmathsetmacro{\ypos}{-1.6*\wavelength/700}
    \draw[->, thick, \col] (0.3, 0) -- (3, \ypos);
}
\node[align=left, text width=2.5cm, popdark] at (3.5, -0.7) {$n=-1$ Spectrum:\\ Violet $\to$ Red};
% n=2 Spectrum (overlap warning)
\foreach \wavelength/\col in {400/violet, 450/blue, 530/green} {
    \pgfmathsetmacro{\ypos}{3.0*\wavelength/700}
    \draw[->, thick, \col, dashed] (0.3, 0) -- (3, \ypos);
}
\node[align=left, text width=2.5cm, popmuted] at (3.5, 1.9) {$n=+2$ (Overlaps $n=1$ Red)};
\end{tikzpicture}
\legende{White light through a grating. Central maximum ($n=0$) is white (all $\lambda$ at $\theta=0$). Higher orders disperse into spectra. Red ($\lambda$ large) deviates \textbf{more} than violet. Note: Order overlap occurs if $n\lambda_{\max} > (n+1)\lambda_{\min}$.}
\end{popfigure}

\begin{attention}[GCE Exam Tip: Grating vs Prism Spectra]
This is a \textbf{classic comparison question}.
\begin{center}
\begin{tabularx}{\linewidth}{|l|X|X|}\hline
\rowcolor{popblueL}
\textbf{Feature} & \textbf{Diffraction Grating} & \textbf{Prism} \\ \hline
\textbf{Mechanism} & Interference / Diffraction & Refraction (Dispersion) \\ \hline
\textbf{Central Image ($n=0$)} & \textbf{White} (undeviated) & \textbf{Deviated} (no "central" order) \\ \hline
\textbf{Order of Colours} & \textbf{Red deviated MORE} ($\theta \propto \lambda$) & \textbf{Violet/Blue deviated MORE} ($n_v > n_r$) \\ \hline
\textbf{Spectrum Type} & Multiple orders ($n=0, \pm1, \pm2\dots$) & Single spectrum per refraction \\ \hline
\textbf{Linearity} & $\sin\theta \propto \lambda$ (Linear in $\sin\theta$) & Non-linear ($n(\lambda)$ curve) \\ \hline
\textbf{Intensity} & Lower (light split into many orders) & Higher (most light in one spectrum) \\ \hline
\textbf{Resolving Power} & High ($R = nN$) & Moderate \\ \hline
\end{tabularx}
\end{center}
\textbf{Memorise:} "Grating: Red bends more. Prism: Violet bends more. Grating central is white."
\end{attention}

\courssub{Worked Example: Overlap of Orders}

\begin{exoresolu}[Order Overlap in a Grating Spectrum]
{A grating has $500\ \mathrm{lines/mm}$. White light ($400\ \mathrm{nm}$ to $700\ \mathrm{nm}$) is incident normally. Show that the second-order spectrum ($n=2$) overlaps with the third-order spectrum ($n=3$). Determine the wavelength range of the overlap.}
\tcblower
\textbf{Solution:}
\begin{enumerate}
    \item \textbf{Grating spacing:}
    $N = 500 \times 10^3 = 5.00 \times 10^5\ \mathrm{m^{-1}}$.
    $d = 1/N = 2.00 \times 10^{-6}\ \mathrm{m} = 2000\ \mathrm{nm}$.
    \item \textbf{Angles for $n=2$:}
    $\sin\theta_2 = 2\lambda/d$.
    For $\lambda_{\min}=400\ \mathrm{nm}$: $\sin\theta_{2,\min} = 2(400)/2000 = 0.400 \Rightarrow \theta \approx 23.6^\circ$.
    For $\lambda_{\max}=700\ \mathrm{nm}$: $\sin\theta_{2,\max} = 2(700)/2000 = 0.700 \Rightarrow \theta \approx 44.4^\circ$.
    Range of $\sin\theta$ for $n=2$: $[0.400, 0.700]$.
    \item \textbf{Angles for $n=3$:}
    $\sin\theta_3 = 3\lambda/d$.
    For $\lambda_{\min}=400\ \mathrm{nm}$: $\sin\theta_{3,\min} = 3(400)/2000 = 0.600 \Rightarrow \theta \approx 36.9^\circ$.
    For $\lambda_{\max}=700\ \mathrm{nm}$: $\sin\theta_{3,\max} = 3(700)/2000 = 1.05 > 1$ (Not visible).
    Max $\lambda$ for $n=3$: $\lambda_{\max,3} = d/3 = 2000/3 \approx 667\ \mathrm{nm}$.
    Range of $\sin\theta$ for $n=3$: $[0.600, 1.000]$ (corresponds to $\lambda \in [400, 667]\ \mathrm{nm}$).
    \item \textbf{Overlap Condition:}
    Overlap occurs where $\sin\theta$ ranges intersect.
    $n=2$ range: $[0.400, 0.700]$.
    $n=3$ range: $[0.600, 1.000]$.
    Intersection: $[0.600, 0.700]$.
    \item \textbf{Wavelength range of overlap:}
    For $n=2$: $\lambda = (d/2)\sin\theta$. At $\sin\theta=0.600$, $\lambda=600\ \mathrm{nm}$. At $\sin\theta=0.700$, $\lambda=700\ \mathrm{nm}$.
    For $n=3$: $\lambda = (d/3)\sin\theta$. At $\sin\theta=0.600$, $\lambda=400\ \mathrm{nm}$. At $\sin\theta=0.700$, $\lambda \approx 467\ \mathrm{nm}$.
    \textbf{Conclusion:} The $n=2$ red/orange light ($600$--$700\ \mathrm{nm}$) overlaps spatially with the $n=3$ violet/blue light ($400$--$467\ \mathrm{nm}$).
\end{enumerate}
\end{exoresolu}

\courssub{Summary and Key Points for Revision}

\begin{aretenir}[Section 3 Checklist]
\begin{itemize}
    \item \textbf{Sources:} Continuous (hot solids), Line (low-pressure gases). Know examples: Na lamp (yellow doublet), Hg lamp (multi-line), Laser (monochromatic, coherent).
    \item \textbf{Coherence:} Same frequency, constant phase difference. Derived from single source.
    \item \textbf{Grating Equation:} $d\sin\theta = n\lambda$. \textbf{Derivation examinable.} $d = 1/N$ ($N$ in lines/m).
    \item \textbf{Orders:} $n=0$ (white, undeviated), $n=\pm 1, \pm 2\dots$ Max order $n_{\max} = \lfloor d/\lambda \rfloor$.
    \item \textbf{Measuring $\lambda$:} Spectrometer or bench method. Measure $\theta$ for $n=1,2$ on both sides. $\lambda = d\sin\theta/n$.
    \item \textbf{Multiple Slits:} Principal maxima sharpen as $N$ increases (width $\propto 1/N$, intensity $\propto N^2$). Resolving power $R = nN$.
    \item \textbf{White Light:} Central max white. Spectra on either side. \textbf{Red deviated MORE than violet.} (Opposite to prism).
    \item \textbf{Overlap:} Check if $n\lambda_{\max} > (n+1)\lambda_{\min}$.
\end{itemize}
\end{aretenir}

\begin{exercice}[Practice Problems]
\begin{enumerate}
    \item A grating with $600\ \mathrm{lines/mm}$ is illuminated normally by a mercury lamp. The green line ($\lambda = 546.1\ \mathrm{nm}$) is observed in the second order ($n=2$). Calculate the diffraction angle $\theta_2$.
    \item In a spectrometer experiment, a student measures the angle for the first-order maximum of a sodium line ($n=1$) as $\theta = 20.50^\circ$ on the left and $\theta = 20.70^\circ$ on the right. The grating has $300\ \mathrm{lines/mm}$. Calculate the mean wavelength.
    \item Explain why the central maximum ($n=0$) is white when white light is used, but the first-order spectra are coloured.
    \item A grating has spacing $d = 1.67 \times 10^{-6}\ \mathrm{m}$. What is the highest order visible for red light ($\lambda = 650\ \mathrm{nm}$)?
    \item Compare the spectrum produced by a prism and a grating for white light. State three differences.
\end{enumerate}
\end{exercice}

\begin{corrige}[Exercise 1]
$N = 600 \times 10^3 = 6 \times 10^5\ \mathrm{m^{-1}}$, $d = 1.667 \times 10^{-6}\ \mathrm{m}$.
$n=2$, $\lambda = 546.1 \times 10^{-9}\ \mathrm{m}$.
$\sin\theta_2 = n\lambda/d = 2(546.1 \times 10^{-9}) / (1.667 \times 10^{-6}) = 0.655$.
$\theta_2 = \sin^{-1}(0.655) = 40.9^\circ$.
\end{corrige}

\begin{corrige}[Exercise 2]
Mean $\theta = (20.50 + 20.70)/2 = 20.60^\circ$.
$d = 1/(300 \times 10^3) = 3.333 \times 10^{-6}\ \mathrm{m}$.
$\lambda = d \sin\theta = 3.333 \times 10^{-6} \times \sin(20.60^\circ) = 3.333 \times 10^{-6} \times 0.3518 = 1.173 \times 10^{-6}\ \mathrm{m} = 1173\ \mathrm{nm}$.
\textbf{Note:} This is in the Infrared! Sodium D-lines are $\approx 589\ \mathrm{nm}$. The angle $20.6^\circ$ for $n=1$ with $300\ \mathrm{lines/mm}$ gives $\lambda \approx 1173\ \mathrm{nm}$. If the question intended sodium light, the grating would need $\approx 600\ \mathrm{lines/mm}$ or the angle would be $\approx 10^\circ$. Always check consistency of data in exam questions.
\end{corrige}

\begin{corrige}[Exercise 3]
At $n=0$, $d\sin\theta = 0 \Rightarrow \theta = 0$ for \textbf{all} wavelengths. All colours arrive at the same angle (undeviated), superimposing to form white light. For $n \neq 0$, $\theta = \sin^{-1}(n\lambda/d)$ depends on $\lambda$, so different colours separate.
\end{corrige}

\begin{corrige}[Exercise 4]
$n_{\max} = \lfloor d/\lambda \rfloor = \lfloor 1.67 \times 10^{-6} / 650 \times 10^{-9} \rfloor = \lfloor 2.57 \rfloor = 2$.
Orders $n=0, \pm 1, \pm 2$ visible. $n=3$ requires $\sin\theta > 1$.
\end{corrige}

\begin{corrige}[Exercise 5]
See the \textbf{Exam Tip} box above: (1) Mechanism (Interference vs Refraction), (2) Central image (White/Undeviated vs Deviated), (3) Order of colours (Red most deviated vs Violet most deviated), (4) Multiple orders vs Single spectrum.
\end{corrige}

\coursec{Section 4: Reflection, Refraction and Refractive Index}

In Section 3 we treated light as a \emph{wave}: diffraction gratings and multiple-slit interference only make sense because light spreads, overlaps and interferes. In this section we change our point of view. When the obstacles and openings are much larger than the wavelength, light travels essentially in straight lines and we may represent it by \cle{rays}. This approximation, called \cle{geometrical optics}, is powerful enough to explain why a straw looks bent in a glass of water, why a swimming pool looks shallower than it is, and how light can be trapped inside a glass fibre thinner than a hair. All of these phenomena follow from two simple laws: the laws of reflection and the laws of refraction.

\courssub{Reflection at a Plane Surface}

\textbf{Observation.} Stand in front of a plane mirror in your house in Yaound\'e or Bamenda. Your image appears to be \emph{behind} the mirror, at the same distance as you are in front of it, and it is the same size as you. Raise your right hand: the image raises the hand on its left side. These everyday facts are completely explained by two laws.

\begin{propriete}[Laws of reflection]
When a light ray strikes a reflecting surface at a point $I$:
\begin{enumerate}
\item The \cle{incident ray}, the \cle{reflected ray} and the \cle{normal} to the surface at $I$ all lie in the same plane (the plane of incidence).
\item The angle of incidence $i$ equals the angle of reflection $r$: $i = r$, both angles being measured from the \emph{normal}.
\end{enumerate}
\end{propriete}

\begin{propriete}[Image formed by a plane mirror]
The image of a point object formed by a plane mirror is:
\begin{itemize}
\item \cle{virtual} (the rays do not actually pass through it; they only \emph{appear} to come from it);
\item situated as far behind the mirror as the object is in front (object distance $=$ image distance);
\item the same size as the object, and \emph{laterally inverted} (left and right interchanged).
\end{itemize}
\end{propriete}

\begin{attention}[Angles are measured from the normal]
In both the laws of reflection and the laws of refraction, angles are measured between the ray and the \textbf{normal} to the surface, never between the ray and the surface itself. If a ray strikes a mirror ``at $30^{\circ}$ to the surface'', the angle of incidence is $i = 90^{\circ} - 30^{\circ} = 60^{\circ}$, and the reflected ray also makes $60^{\circ}$ with the normal. Mixing up the two references is one of the most frequent errors at the GCE.
\end{attention}

\courssub{Refraction at a Plane Interface}

\textbf{Observation.} Put a straight straw in a transparent glass of water and look from the side: the straw appears \emph{bent} at the water surface. Look down into a clear stream in the West Region: the bed seems closer to the surface than it really is. Light changes direction when it crosses the boundary between two transparent media. This bending is called \cle{refraction}, and its physical origin is the change in the \emph{speed} of light from one medium to the other.

\begin{definition}[Absolute refractive index]
The \cle{absolute refractive index} of a transparent medium is the ratio of the speed of light in vacuum to the speed of light in the medium:
\[ n = \frac{c}{v} \]
where $c = \SI{3.00e8}{m.s^{-1}}$ and $v$ is the speed of light in the medium. Since light always travels slower in a material medium than in vacuum, $n \geq 1$ always. For air, $n_{\text{air}} \approx 1.00$; for water $n \approx 1.33$; for ordinary glass $n \approx 1.5$; for diamond $n \approx 2.42$. A medium with a larger index is said to be optically \cle{denser}.
\end{definition}

\begin{definition}[Relative refractive index]
When light passes from medium 1 (index $n_1$) into medium 2 (index $n_2$), the \cle{relative refractive index} of medium 2 with respect to medium 1 is
\[ n_{21} = \frac{n_2}{n_1} = \frac{v_1}{v_2} = \frac{\sin\theta_1}{\sin\theta_2}. \]
This quantity appears directly in Snell's law and is often the one measured experimentally (e.g., glass relative to air).
\end{definition}

\begin{propriete}[Laws of refraction --- Snell's law]
When a ray passes from medium 1 (index $n_1$) into medium 2 (index $n_2$):
\begin{enumerate}
\item The incident ray, the refracted ray and the normal at the point of incidence are coplanar.
\item \textbf{Snell's law} (examinable and must be quoted correctly):
\[ n_1 \sin\theta_1 = n_2 \sin\theta_2 \]
where $\theta_1$ and $\theta_2$ are the angles of incidence and refraction, measured from the normal.
\end{enumerate}
For a ray going from \emph{air} (or vacuum) into a medium of index $n$, this reduces to
\[ n = \frac{\sin i}{\sin r}. \]
Consequences: entering a \emph{denser} medium ($n_2 > n_1$), the ray bends \textbf{towards} the normal ($\theta_2 < \theta_1$); entering a \emph{rarer} medium, it bends \textbf{away} from the normal.
\end{propriete}

\begin{popfigure}
\begin{center}
\begin{tikzpicture}[scale=1.1, >=stealth]
% interface
\fill[popblueL, opacity=0.5] (-4,-2.6) rectangle (4,0);
\draw[thick, popink] (-4,0) -- (4,0);
\draw[dashed, popmuted] (0,-2.6) -- (0,2.6);
\node[popmuted, right] at (0.1,2.45) {normal};
% incident ray
\draw[very thick, poporange, ->] (-2.6,2.6) -- (0,0);
% refracted ray
\draw[very thick, popblue, ->] (0,0) -- (1.55,-2.6);
% angle arcs
\draw[poporange] (-0.85,0.85) arc[start angle=135, end angle=90, radius=1.2];
\node[poporange] at (-0.85,1.55) {$i$};
\draw[popblue] (0.55,-0.92) arc[start angle=-59, end angle=-90, radius=1.1];
\node[popblue] at (0.85,-1.15) {$r$};
% labels
\node[popink] at (-3.4,0.35) {air ($n_1 \approx 1$)};
\node[popink] at (-3.4,-0.4) {glass ($n_2$)};
\node[popink, align=center] at (2.6,-2.2) {denser medium:\\ ray bends\\ towards normal};
\end{tikzpicture}
\end{center}
\legende{Refraction at an air--glass interface: the refracted ray bends towards the normal, so $r < i$.}
\end{popfigure}

\begin{exoresolu}[Applying Snell's law]
A ray of light in air strikes the surface of water ($n = 1.33$) at an angle of incidence $i = 50.0^{\circ}$. Calculate the angle of refraction $r$.
\tcblower
\textbf{Solution.} Apply Snell's law between air and water:
\[ n_{\text{air}}\sin i = n_{\text{water}}\sin r \quad\Longrightarrow\quad \sin r = \frac{1.00 \times \sin 50.0^{\circ}}{1.33} = \frac{0.766}{1.33} = 0.576. \]
Hence $r = \sin^{-1}(0.576) \approx 35.2^{\circ}$. As expected, $r < i$: the ray bends towards the normal on entering the denser medium.
\end{exoresolu}

\textbf{Real and apparent depth.} Refraction explains why objects under water look nearer the surface than they really are. Rays from a point $O$ on the bottom bend \emph{away} from the normal as they leave the water; the eye receives them as if they came from a point $I$ above $O$. For near-normal viewing (small angles), geometry together with Snell's law (using the small-angle approximation $\sin\theta \approx \tan\theta \approx \theta$) gives a simple and very useful result.

\begin{methode}[Finding $n$ by the real/apparent depth method]
\begin{enumerate}
\item Place a mark (e.g.\ a coin or a cross on paper) at the bottom of a transparent block or beaker of liquid.
\item View it from directly above and locate the apparent position of the image (e.g.\ with a travelling microscope or a search pin using the no-parallax method).
\item Measure the \cle{real depth} $d$ (mark to surface) and the \cle{apparent depth} $d'$ (image to surface).
\item Compute the refractive index from
\[ n = \frac{\text{real depth}}{\text{apparent depth}} = \frac{d}{d'}. \]
\end{enumerate}
\emph{Example:} a fish lies $\SI{1.20}{m}$ below the surface of a pond but appears to be only $\SI{0.90}{m}$ deep. Then $n = 1.20/0.90 = 1.33$ --- the index of water.
\end{methode}

\begin{attention}[Who sees whom?]
The formula $n = d/d'$ applies when the observer is in \emph{air} looking into the denser medium, so the object appears \emph{shallower}. Symmetrically, a fish looking up at you sees you \emph{higher} than you really are (apparent height $= n \times$ real height). Always ask: in which medium is the observer?
\end{attention}

\courssub{Refraction Through a Rectangular Glass Block}

When a ray crosses a parallel-sided glass block, it is refracted twice: towards the normal on entering, away from the normal on leaving. Because the two faces are parallel, the two deviations cancel in direction.

\begin{propriete}[Emergent ray and lateral displacement]
For a rectangular block with parallel faces surrounded by the same medium on both sides, the \cle{emergent ray is parallel to the incident ray}: the block produces no net change of direction, only a sideways shift called the \cle{lateral displacement} $d$. The displacement increases with the thickness $t$ of the block, the angle of incidence $i$ and the refractive index $n$. For a block of thickness $t$ in air:
\[ d = t\,\frac{\sin(i-r)}{\cos r}, \]
where $r$ is the angle of refraction inside the block given by $\sin r = \sin i / n$.
\end{propriete}

\begin{popfigure}
\begin{center}
\begin{tikzpicture}[scale=1.0, >=stealth]
% block
\fill[popblueL, opacity=0.5] (-1.2,-0.2) rectangle (3.2,2.6);
\draw[thick, popink] (-1.2,-0.2) rectangle (3.2,2.6);
% normals
\draw[dashed, popmuted] (-0.4,1.9) -- (0.9,1.9);
\draw[dashed, popmuted] (1.6,0.5) -- (2.9,0.5);
% rays
\draw[very thick, poporange, ->] (-2.6,1.9) -- (-0.2,1.9);
\draw[very thick, poporange] (-0.2,1.9) -- (0.0,1.9);
\draw[very thick, popgreen, ->] (-0.2,1.9) -- (2.2,0.5);
\draw[very thick, poporange, ->] (2.2,0.5) -- (4.2,0.5);
% extended incident direction
\draw[dashed, poporange] (-0.2,1.9) -- (3.4,1.9);
% lateral displacement
\draw[<->, poppurple, thick] (3.6,1.9) -- (3.6,0.5);
\node[poppurple, right] at (3.65,1.2) {$d$};
% angles
\node[poporange] at (-0.85,1.65) {$i$};
\node[popgreen] at (0.35,1.55) {$r$};
\node[popink] at (1.0,2.85) {glass block};
\end{tikzpicture}
\end{center}
\legende{A ray through a rectangular glass block: the emergent ray is parallel to the incident ray but laterally displaced by $d$.}
\end{popfigure}

\begin{experience}[Determining $n$ by ray tracing through a glass block]
Place a rectangular glass block on a sheet of paper and trace its outline. Draw an incident ray, stick two pins on it, then look through the block and stick two more pins so that all four appear aligned. Remove the block, join the pin lines, and draw the refracted ray inside the block. Measure $i$ and $r$ with a protractor and compute $n = \sin i / \sin r$. Repeating for several angles of incidence and averaging gives a reliable value (typically $n \approx 1.5$ for ordinary glass). A graph of $\sin i$ against $\sin r$ is a straight line through the origin whose gradient is $n$ --- an excellent way to present the results at the practical examination.
\end{experience}

\begin{exoresolu}[Lateral displacement in a glass block]
A ray of light in air strikes a rectangular glass block of thickness $\SI{5.0}{cm}$ and refractive index $1.50$ at an angle of incidence $i = 40.0^{\circ}$. Calculate the lateral displacement $d$ of the emergent ray.
\tcblower
\textbf{Solution.}
First find the angle of refraction $r$ inside the glass:
\[ \sin r = \frac{\sin 40.0^{\circ}}{1.50} = \frac{0.6428}{1.50} = 0.4285 \quad\Rightarrow\quad r = \sin^{-1}(0.4285) \approx 25.4^{\circ}. \]
Now apply the displacement formula:
\[ d = t\,\frac{\sin(i-r)}{\cos r} = \SI{5.0}{cm} \times \frac{\sin(40.0^{\circ}-25.4^{\circ})}{\cos 25.4^{\circ}} = \SI{5.0}{cm} \times \frac{\sin 14.6^{\circ}}{0.903} = \SI{5.0}{cm} \times \frac{0.252}{0.903} \approx \SI{1.40}{cm}. \]
\end{exoresolu}

\courssub{Dispersion: the Dependence of Refractive Index on Colour}

\textbf{Observation.} After rain, when the Sun reappears over the hills of the North West Region, a rainbow arcs across the sky. White sunlight has been split into colours by millions of raindrops. The reason is that the refractive index of a medium is \emph{not exactly the same for all colours}.

\begin{definition}[Dispersion]
\cle{Dispersion} is the separation of white light into its constituent colours because the refractive index of a medium depends on the wavelength (colour) of the light. For most transparent materials $n$ is slightly larger for \emph{violet} light than for \emph{red} light, so violet is deviated more than red by a prism or a water drop.
\end{definition}

When white light passes through a glass prism, each colour is refracted by a different amount: the emerging beam spreads into a \cle{spectrum}, red deviated least, violet most. In a rainbow, sunlight undergoes refraction, internal reflection and refraction again inside each raindrop; dispersion at the refractions separates the colours, and each observer sees light from millions of drops arranged on a cone of about $42^{\circ}$ around the anti-solar direction.

\begin{savaistu}[Why is the sky blue and the sunset red?]
Dispersion is not the only colour effect in the atmosphere. Air molecules scatter short wavelengths (blue) much more strongly than long ones (red) --- this is Rayleigh scattering, not refraction. The blue of the midday sky and the red of the Sun near the horizon at sunset over the Atlantic at Limbe both come from this scattering, while the rainbow comes from dispersion in raindrops.
\end{savaistu}

\courssub{Total Internal Reflection and the Critical Angle}

\textbf{Observation.} A swimmer under clear water looking up sees the outside world compressed into a bright circle overhead; outside that circle the surface acts like a mirror. Light travelling \emph{inside} the water can be totally reflected at the surface. This is \cle{total internal reflection} (TIR), and it only happens when light tries to pass from a \emph{denser} to a \emph{rarer} medium.

Consider a ray going from glass ($n_1 = n$) to air ($n_2 = 1$). By Snell's law, $\sin r = n\sin i$. As $i$ increases, $r$ increases faster and reaches $90^{\circ}$ first. Beyond that point no refracted ray can exist.

\begin{definition}[Critical angle]
The \cle{critical angle} $C$ is the angle of incidence in the denser medium for which the angle of refraction in the rarer medium is $90^{\circ}$ (the refracted ray grazes along the surface). Setting $r = 90^{\circ}$ in Snell's law $n_1\sin i = n_2\sin r$:
\[ n_1\sin C = n_2 \quad\Longrightarrow\quad \boxed{\ \sin C = \frac{n_2}{n_1}\ } \quad (n_1 > n_2). \]
For the special case where the rarer medium is air ($n_2 \approx 1$), this reduces to $\sin C = 1/n_1$. For $i > C$, \textbf{all} the light is reflected back into the denser medium: total internal reflection.
\end{definition}

\begin{propriete}[Conditions for total internal reflection]
\begin{enumerate}
\item Light must travel from the optically \cle{denser} medium towards the \cle{rarer} medium ($n_1 > n_2$).
\item The angle of incidence must \emph{exceed} the critical angle: $i > C$.
\end{enumerate}
For glass ($n = 1.50$) to air: $C = \sin^{-1}(1/1.50) \approx 41.8^{\circ}$. For water ($n = 1.33$): $C \approx 48.8^{\circ}$. For diamond ($n = 2.42$): $C \approx 24.4^{\circ}$.
\end{propriete}

\begin{popfigure}
\begin{center}
\begin{tikzpicture}[scale=1.05, >=stealth]
% media
\fill[popblueL, opacity=0.5] (-4.5,-2.4) rectangle (4.5,0);
\draw[thick, popink] (-4.5,0) -- (4.5,0);
\draw[dashed, popmuted] (0,-2.4) -- (0,2.4);
\node[popink] at (-3.9,-0.35) {glass ($n_1$)};
\node[popink] at (-3.9,0.35) {air ($n_2$)};
\node[popink] at (0,-2.7) {O};
\fill[popink] (0,-2.4) circle (1.5pt);
% ray 1: refraction
\draw[very thick, popblue, ->] (0,-2.4) -- (-1.1,0);
\draw[very thick, popblue, ->] (-1.1,0) -- (-2.5,1.9);
\node[popblue, left] at (-2.6,1.0) {$i<C$: refraction};
% ray 2: critical
\draw[very thick, popgreen, ->] (0,-2.4) -- (-0.1,0);
\draw[very thick, popgreen, ->] (0,0) -- (2.9,0.06);
\node[popgreen, above] at (2.1,0.12) {$i=C$: grazes surface};
% ray 3: TIR
\draw[very thick, poporange, ->] (0,-2.4) -- (1.5,0);
\draw[very thick, poporange, ->] (1.5,0) -- (3.0,-2.0);
\node[poporange, right, align=center] at (3.1,-1.0) {$i>C$:\\ total internal\\ reflection};
% critical angle arc
\draw[popgreen] (0,-1.2) arc[start angle=90, end angle=90, radius=1.2];
\draw[popgreen] (0,-1.2) -- (0,-1.2);
\node[popgreen] at (0.35,-1.5) {$C$};
\end{tikzpicture}
\end{center}
\legende{Three rays from a source $O$ in glass reaching the glass--air boundary: ordinary refraction ($i<C$), the critical ray ($i=C$, refracted along the surface), and total internal reflection ($i>C$).}
\end{popfigure}

\begin{attention}[TIR is impossible from rare to dense]
The formula $\sin C = n_2/n_1$ and the whole phenomenon of total internal reflection apply \textbf{only} when light travels from the denser medium towards the rarer one. If light goes from air into glass, $\sin r = \sin i / n \leq 1/n < 1$ always has a solution, so a refracted ray \emph{always} exists: TIR can never occur. Writing $\sin C = 1/n$ for an air-to-glass ray is a classic error that costs marks.
\end{attention}

\begin{exoresolu}[Critical angle of water and of diamond]
Calculate the critical angle for (a) water of index $1.33$, (b) diamond of index $2.42$, each in contact with air.
\tcblower
\textbf{Solution.} In each case light goes from the dense medium to air, so $\sin C = 1/n$.\\
(a) Water: $\sin C = 1/1.33 = 0.752$, so $C = \sin^{-1}(0.752) \approx 48.8^{\circ}$.\\
(b) Diamond: $\sin C = 1/2.42 = 0.413$, so $C = \sin^{-1}(0.413) \approx 24.4^{\circ}$.\\
The small critical angle of diamond explains its sparkle: light entering a well-cut diamond strikes the inner faces at angles well above $24.4^{\circ}$, is totally reflected several times, and comes back out towards the eye.
\end{exoresolu}

\courssub{Applications of Total Internal Reflection}

\textbf{Optical fibres.} An optical fibre is a thin strand of very pure glass (core, index $n_1$) surrounded by a cladding of slightly lower index $n_2$. Light entering the core at a suitable angle strikes the core--cladding boundary at $i > C$ and is totally internally reflected again and again, travelling kilometres with very little loss.

\begin{popfigure}
\begin{center}
\begin{tikzpicture}[scale=1.0, >=stealth]
% cladding
\fill[popblueL, opacity=0.4] (0,-0.9) rectangle (9,0.9);
% core
\fill[white] (0,-0.45) rectangle (9,0.45);
\draw[thick, popink] (0,0.45) -- (9,0.45);
\draw[thick, popink] (0,-0.45) -- (9,-0.45);
\draw[thick, popmuted] (0,0.9) -- (9,0.9);
\draw[thick, popmuted] (0,-0.9) -- (9,-0.9);
% zigzag ray
\draw[very thick, poporange, ->] (-0.8,0) -- (0.6,0.45);
\draw[very thick, poporange] (0.6,0.45) -- (1.8,-0.45);
\draw[very thick, poporange] (1.8,-0.45) -- (3.0,0.45);
\draw[very thick, poporange] (3.0,0.45) -- (4.2,-0.45);
\draw[very thick, poporange] (4.2,-0.45) -- (5.4,0.45);
\draw[very thick, poporange] (5.4,0.45) -- (6.6,-0.45);
\draw[very thick, poporange] (6.6,-0.45) -- (7.8,0.45);
\draw[very thick, poporange, ->] (7.8,0.45) -- (9.6,0);
% labels
\node[popink, above] at (4.5,0.95) {cladding ($n_2$)};
\node[popink] at (8.4,-0.72) {core ($n_1 > n_2$)};
\node[poporange, align=center] at (2.2,-1.35) {total internal reflection\\ at each bounce ($i > C$)};
\end{tikzpicture}
\end{center}
\legende{Light guided along an optical fibre by repeated total internal reflection at the core--cladding boundary.}
\end{popfigure}

\begin{itemize}
\item \textbf{Telecommunications.} Fibre-optic cables carry telephone and internet signals as pulses of light. Cameroon is linked to the world network by submarine fibre-optic cables landing on the Atlantic coast, and fibre backbones connect the main towns: the speed and capacity of these links rely entirely on total internal reflection.
\item \textbf{Endoscopy.} In medicine, a flexible bundle of fibres carries light down into a patient's body and brings the image back, allowing a doctor to examine, for example, the stomach without surgery.
\end{itemize}

\textbf{Totally reflecting prisms.} A $45^{\circ}$--$45^{\circ}$--$90^{\circ}$ glass prism has $C \approx 42^{\circ}$. A ray entering normally through one face strikes the hypotenuse at $45^{\circ} > C$ and is totally reflected. Such prisms replace mirrors in binoculars, periscopes and single-lens reflex cameras: unlike a metal mirror, the reflecting surface loses no light and never tarnishes.

\textbf{Mirages.} On a very hot tarred road in the Far North Region, the air just above the tarmac is hotter and optically \emph{rarer} than the air above. Rays from the sky heading downwards are gradually bent and can undergo total internal reflection in this hot layer, reaching the driver's eye from below: the road ahead appears ``wet''. This inferior mirage is refraction and TIR in a medium whose index varies continuously.

\textbf{Sparkling of diamond.} As computed above, diamond's critical angle is only about $24^{\circ}$. A skilfully cut stone traps entering light by multiple total internal reflections and returns it to the observer's eye, while strong dispersion separates the colours --- hence the brilliant, fiery sparkle.

\begin{aretenir}[Key points of Section 4]
\begin{itemize}
\item Reflection: incident ray, reflected ray and normal are coplanar; $i = r$. A plane mirror gives a virtual, upright, same-size, laterally inverted image as far behind the mirror as the object is in front.
\item Refractive index: absolute $n = c/v$; relative $n_{21} = n_2/n_1 = v_1/v_2$. Snell's law $n_1\sin\theta_1 = n_2\sin\theta_2$; from air, $n = \sin i / \sin r$. Angles are measured from the \emph{normal}.
\item Real and apparent depth: $n = d/d'$ (observer in air, object in the medium, near-normal viewing).
\item A parallel-sided block displaces a ray laterally by $d = t\sin(i-r)/\cos r$ but the emergent ray is parallel to the incident ray.
\item Dispersion: $n$ depends on wavelength; violet is deviated more than red (prism, rainbow).
\item Total internal reflection needs a dense-to-rare crossing ($n_1 > n_2$) and $i > C$, with $\sin C = n_2/n_1$; applications: optical fibres, reflecting prisms, mirages, diamond's brilliance.
\end{itemize}
\end{aretenir}

\begin{exercice}[Practice]
\begin{enumerate}
\item A ray in air enters a glass block ($n = 1.52$) at $i = 40.0^{\circ}$. Find the angle of refraction.
\item A swimming pool is $\SI{2.0}{m}$ deep. How deep does it appear to someone looking almost vertically down? Take $n_{\text{water}} = 1.33$.
\item Show that the critical angle for a glass of index $1.60$ in contact with water of index $1.33$ is given by $\sin C = 1.33/1.60$, and calculate $C$. Why is the formula not $\sin C = 1/1.60$ here?
\item A rectangular glass block of thickness $\SI{4.0}{cm}$ and refractive index $1.50$ is placed on a printed page. A ray of light from a point on the page enters the block at an angle of incidence of $30.0^{\circ}$. Calculate the lateral displacement of the emergent ray.
\item Explain, with a ray diagram, why a totally reflecting prism is preferred to a plane mirror in good-quality binoculars.
\end{enumerate}
\end{exercice}

\begin{corrige}[Exercise 1]
\textbf{1.} $\sin r = \sin 40.0^{\circ}/1.52 = 0.643/1.52 = 0.423$, so $r \approx 25.0^{\circ}$.\\
\textbf{2.} Apparent depth $d' = d/n = 2.0/1.33 \approx \SI{1.5}{m}$.\\
\textbf{3.} The second medium is water, not air. Snell's law at the critical angle gives $n_{\text{glass}}\sin C = n_{\text{water}}\sin 90^{\circ}$, hence $\sin C = n_{\text{water}}/n_{\text{glass}} = 1.33/1.60 = 0.831$ and $C \approx 56.2^{\circ}$. The simple form $\sin C = 1/n$ is only valid when the rarer medium is air (or vacuum).\\
\textbf{4.} First find $r$: $\sin r = \sin 30.0^{\circ}/1.50 = 0.500/1.50 = 0.333$, so $r \approx 19.5^{\circ}$. Then $d = t\sin(i-r)/\cos r = \SI{4.0}{cm} \times \sin(10.5^{\circ})/\cos 19.5^{\circ} = \SI{4.0}{cm} \times 0.182/0.943 \approx \SI{0.77}{cm}$.\\
\textbf{5.} The ray enters the prism normally, meets the hypotenuse at $45^{\circ} > C \approx 42^{\circ}$, and is totally reflected with essentially no loss of light; a metal mirror absorbs a fraction of the light and its surface deteriorates with time. The prism also folds the light path, making the instrument compact.
\end{corrige}

\coursec{Section 5: Lenses, Prisms and Optical Instruments}

In the previous section we studied how light changes direction at a single plane interface through refraction, and how the refractive index governs this bending. We now extend these ideas to \emph{curved} refracting surfaces — lenses — and to combinations of plane surfaces — prisms. These elements allow us to control light in a precise way, forming real images on a screen or virtual images for direct viewing. From the simple magnifying glass to the compound microscope and the astronomical telescope, the principles are the same: refraction at spherical or plane surfaces, governed by Snell's law. We begin with the thin lens, the workhorse of geometric optics.

\courssub{Converging and Diverging Lenses}

A \cle{lens} is a piece of transparent material (usually glass or plastic) bounded by two spherical surfaces (or one spherical and one plane). Its action on light depends on its shape.

\begin{popfigure}
\begin{tikzpicture}[scale=1.2, >=stealth]
% Converging lens (double convex)
\draw[popblue, thick] (-0.1,1.5) .. controls (0.3,1.5) and (0.3,-1.5) .. (-0.1,-1.5);
\draw[popblue, thick] (0.1,1.5) .. controls (-0.3,1.5) and (-0.3,-1.5) .. (0.1,-1.5);
% Principal axis
\draw[popmuted, dashed] (-3,0) -- (3,0) node[right] {Principal axis};
% Optical centre
\fill[popdark] (0,0) circle (1.5pt) node[below left] {O};
% Focal points
\fill[poporange] (-2,0) circle (1.5pt) node[below] {F};
\fill[poporange] (2,0) circle (1.5pt) node[below] {F'};
% Parallel rays from left
\foreach \y in {1,0.5,-0.5,-1} {
    \draw[popgreen, ->] (-3,\y) -- (0,\y);
    \draw[popgreen, ->] (0,\y) -- (3,0);
}
\draw[popgreen, ->] (-3,0) -- (3,0);
% Labels
\node[popblue] at (0,1.8) {Converging lens};
\end{tikzpicture}
\legende{A double convex (converging) lens brings parallel rays to a real focus F' on the opposite side. The distance OF' is the focal length $f$.}
\end{popfigure}

\begin{popfigure}
\begin{tikzpicture}[scale=1.2, >=stealth]
% Diverging lens (double concave)
\draw[popblue, thick] (-0.1,1.5) .. controls (-0.3,1.5) and (-0.3,-1.5) .. (-0.1,-1.5);
\draw[popblue, thick] (0.1,1.5) .. controls (0.3,1.5) and (0.3,-1.5) .. (0.1,-1.5);
% Principal axis
\draw[popmuted, dashed] (-3,0) -- (3,0) node[right] {Principal axis};
% Optical centre
\fill[popdark] (0,0) circle (1.5pt) node[below left] {O};
% Focal points (virtual)
\fill[poporange] (-2,0) circle (1.5pt) node[below] {F};
\fill[poporange] (2,0) circle (1.5pt) node[below] {F'};
% Parallel rays from left
\foreach \y in {1,0.5,-0.5,-1} {
    \draw[popgreen, ->] (-3,\y) -- (0,\y);
    \draw[popgreen, dashed] (0,\y) -- (-2,0);
    \draw[popgreen, ->] (0,\y) -- (3,\y);
}
\draw[popgreen, ->] (-3,0) -- (3,0);
% Labels
\node[popblue] at (0,1.8) {Diverging lens};
\end{tikzpicture}
\legende{A double concave (diverging) lens spreads parallel rays as if they originated from a virtual focus F on the same side as the incident light. The focal length $f$ is negative.}
\end{popfigure}

\begin{definition}[Principal Focus of a Converging Lens]
The \cle{principal focus} $F'$ of a converging lens is the point on the principal axis where rays incident parallel to the axis converge after refraction through the lens. The distance from the optical centre $O$ to $F'$ is the \cle{focal length} $f$. For a diverging lens, the principal focus $F$ is the virtual point from which incident parallel rays appear to diverge; its focal length is taken as negative.
\end{definition}

\courssub{Power of a Lens}

The \cle{power} $P$ of a lens quantifies its ability to converge or diverge light. It is defined as the reciprocal of the focal length expressed in metres.

\begin{propriete}[Power of a Lens]
\[
P = \frac{1}{f}
\]
where $f$ is the focal length in metres (m). The unit of power is the \cle{dioptre} (symbol D). A converging lens has positive power ($f>0$); a diverging lens has negative power ($f<0$).
\end{propriete}

\begin{attention}[Units for Power]
A very common mistake is to use $f$ in centimetres when calculating $P$. Remember: \textbf{always convert $f$ to metres first}. For example, a lens of focal length $20\ \mathrm{cm}$ has $f = 0.20\ \mathrm{m}$, so $P = 1/0.20 = +5\ \mathrm{D}$, not $+0.05\ \mathrm{D}$.
\end{attention}

\begin{exemplebox}[Calculating Power]
A converging lens has a focal length of $25\ \mathrm{cm}$. A diverging lens has a focal length of $50\ \mathrm{cm}$. Find their powers.
\begin{align*}
\text{Converging: } f &= +0.25\ \mathrm{m} \quad\Rightarrow\quad P = \frac{1}{0.25} = +4.0\ \mathrm{D}. \\
\text{Diverging: } f &= -0.50\ \mathrm{m} \quad\Rightarrow\quad P = \frac{1}{-0.50} = -2.0\ \mathrm{D}.
\end{align*}
\end{exemplebox}

\courssub{Ray Diagrams for a Converging Lens}

To locate the image formed by a thin converging lens, we use three \cle{principal rays} (any two are sufficient):

\begin{enumerate}
\item A ray parallel to the principal axis refracts through the principal focus $F'$.
\item A ray passing through the optical centre $O$ continues undeviated (for a thin lens).
\item A ray passing through the first principal focus $F$ emerges parallel to the principal axis.
\end{enumerate}

\begin{methode}[Drawing Lens Ray Diagrams]
\begin{enumerate}
\item Draw the principal axis, the lens as a vertical line (or double arrow for converging), and mark $O$, $F$ and $F'$ at equal distances $f$ on either side.
\item Draw the object as an upright arrow perpendicular to the axis at distance $u$ from $O$.
\item From the tip of the object, draw Ray 1 (parallel to axis $\to$ through $F'$) and Ray 2 (through $O$ undeviated).
\item The intersection of the refracted rays (or their backward extensions) gives the tip of the image.
\item Measure $v$ (image distance) and image height $h'$.
\end{enumerate}
\end{methode}

The nature of the image depends on the object position relative to $F$ and $2F$ (where $2F$ is the point at distance $2f$ from $O$). We summarise the five standard cases:

\begin{center}
\begin{tabularx}{\linewidth}{|l|X|X|X|X|}\hline
\textbf{Object position} & \textbf{Image position} & \textbf{Nature} & \textbf{Size} & \textbf{Practical example} \\ \hline
Beyond $2F$ ($u>2f$) & Between $F'$ and $2F'$ & Real, inverted & Diminished & Camera \\ \hline
At $2F$ ($u=2f$) & At $2F'$ ($v=2f$) & Real, inverted & Same size & Symmetric projector \\ \hline
Between $F$ and $2F$ ($f<u<2f$) & Beyond $2F'$ ($v>2f$) & Real, inverted & Magnified & Projector, enlarger \\ \hline
At $F$ ($u=f$) & At infinity & Real, inverted & Highly magnified & Spotlight, collimator \\ \hline
Inside $F$ ($u<f$) & Same side as object ($v<0$) & Virtual, upright & Magnified & Magnifying glass \\ \hline
\end{tabularx}
\end{center}

\begin{popfigure}
\begin{tikzpicture}[scale=0.9, >=stealth]
% Three ray diagrams side by side
% Common settings
\def\lens{%
    \draw[popblue, thick] (0,1.2) .. controls (0.3,1.2) and (0.3,-1.2) .. (0,-1.2);
    \draw[popblue, thick] (0,1.2) .. controls (-0.3,1.2) and (-0.3,-1.2) .. (0,-1.2);
    \fill[popdark] (0,0) circle (1.5pt);
    \draw[popmuted, dashed] (-3,0) -- (4,0);
}
% Case 1: Object beyond 2F (diminished real)
\begin{scope}[xshift=-7cm]
\lens
\fill[poporange] (-2,0) circle (1.5pt) node[below] {F};
\fill[poporange] (2,0) circle (1.5pt) node[below] {F'};
\fill[poporange] (-4,0) circle (1.5pt) node[below] {2F};
\fill[poporange] (4,0) circle (1.5pt) node[below] {2F'};
% Object
\draw[popdark, thick] (-3.5,0) -- (-3.5,1.2) node[above] {Object};
% Ray 1: parallel then through F'
\draw[popgreen, ->] (-3.5,1.2) -- (0,1.2);
\draw[popgreen, ->] (0,1.2) -- (2,0);
% Ray 2: through O
\draw[popgreen, ->] (-3.5,1.2) -- (0,0) -- (1.5,-0.5);
% Image
\draw[popred, thick] (1.5,0) -- (1.5,-0.5) node[below] {Image};
\node at (0,1.6) {\textbf{(a) $u>2f$}};
\end{scope}

% Case 2: Object between F and 2F (magnified real)
\begin{scope}[xshift=0cm]
\lens
\fill[poporange] (-2,0) circle (1.5pt) node[below] {F};
\fill[poporange] (2,0) circle (1.5pt) node[below] {F'};
\fill[poporange] (-4,0) circle (1.5pt) node[below] {2F};
\fill[poporange] (4,0) circle (1.5pt) node[below] {2F'};
% Object
\draw[popdark, thick] (-3,0) -- (-3,1) node[above] {Object};
% Ray 1
\draw[popgreen, ->] (-3,1) -- (0,1);
\draw[popgreen, ->] (0,1) -- (2,0);
% Ray 2
\draw[popgreen, ->] (-3,1) -- (0,0) -- (4,-1.33);
% Image
\draw[popred, thick] (4,0) -- (4,-1.33) node[below] {Image};
\node at (0,1.6) {\textbf{(b) $f<u<2f$}};
\end{scope}

% Case 3: Object inside F (magnified virtual)
\begin{scope}[xshift=7cm]
\lens
\fill[poporange] (-2,0) circle (1.5pt) node[below] {F};
\fill[poporange] (2,0) circle (1.5pt) node[below] {F'};
% Object
\draw[popdark, thick] (-1,0) -- (-1,0.8) node[above] {Object};
% Ray 1: parallel then through F'
\draw[popgreen, ->] (-1,0.8) -- (0,0.8);
\draw[popgreen, ->] (0,0.8) -- (2,0);
% Ray 2: through O
\draw[popgreen, ->] (-1,0.8) -- (0,0) -- (-2,-1.6);
% Virtual image (backward extensions)
\draw[popred, dashed] (0,0.8) -- (-2,0);
\draw[popred, dashed] (0,0) -- (-2,-1.6);
\draw[popred, thick] (-2,0) -- (-2,-1.6) node[below] {Virtual Image};
\node at (0,1.6) {\textbf{(c) $u<f$}};
\end{scope}
\end{tikzpicture}
\legende{Ray diagrams for a converging lens: (a) object beyond $2F$ gives a diminished real image; (b) object between $F$ and $2F$ gives a magnified real image; (c) object inside $F$ gives a magnified virtual image.}
\end{popfigure}

\courssub{The Thin Lens Formula and Linear Magnification}

The geometric construction leads to an algebraic relation between object distance $u$, image distance $v$, and focal length $f$. We adopt the \cle{real-is-positive sign convention} (also called the Cartesian convention for lenses):

\begin{itemize}
\item Distances measured in the direction of incident light (left to right) are positive.
\item Distances measured opposite to the incident light are negative.
\item Object distance $u$ is positive for a real object (on the incident side).
\item Image distance $v$ is positive for a real image (on the opposite side), negative for a virtual image (on the same side).
\item Focal length $f$ is positive for converging lenses, negative for diverging lenses.
\item Heights measured upward from the axis are positive; downward are negative.
\end{itemize}

\begin{propriete}[Thin Lens Formula and Magnification]
For a thin lens:
\[
\frac{1}{f} = \frac{1}{u} + \frac{1}{v}
\]
The \cle{linear magnification} (transverse magnification) is
\[
m = \frac{h'}{h} = \frac{v}{u}
\]
where $h$ is object height and $h'$ is image height. If $m$ is negative, the image is inverted relative to the object; if positive, upright. $|m|>1$ means magnified, $|m|<1$ means diminished.
\end{propriete}

\begin{exoresolu}[Using the Lens Formula]
A converging lens of focal length $15\ \mathrm{cm}$ forms a real image of an object placed $45\ \mathrm{cm}$ from the lens. Find the image distance and the magnification.
\tcblower
Given: $f = +15\ \mathrm{cm}$, $u = +45\ \mathrm{cm}$ (real object).
\[
\frac{1}{v} = \frac{1}{f} - \frac{1}{u} = \frac{1}{15} - \frac{1}{45} = \frac{3-1}{45} = \frac{2}{45}
\]
\[
v = \frac{45}{2} = +22.5\ \mathrm{cm}
\]
The positive $v$ confirms a real image on the opposite side.
\[
m = \frac{v}{u} = \frac{22.5}{45} = +0.5
\]
Wait: $m = v/u = +0.5$? But the image is real and inverted, so $m$ should be negative. In the real-is-positive convention, $m = v/u$ gives a positive value for real images because both $u$ and $v$ are positive. However, the \emph{physical} magnification (including orientation) is $m = -v/u$ in many textbooks. The GCE syllabus uses $m = v/u$ with the understanding that a real image is inverted, so the sign of $m$ from the formula does not indicate orientation. \textbf{Always state the nature separately:} here the image is real, inverted, and diminished ($|m|=0.5$).
\end{exoresolu}

\begin{attention}[Sign of Magnification]
In the real-is-positive convention, $m = v/u$ is \textbf{always positive for real images} (since $u>0, v>0$) and \textbf{negative for virtual images} ($v<0$). This is opposite to the Cartesian sign convention used in some textbooks. For GCE, use $m = v/u$ and explicitly describe the image as inverted (real) or upright (virtual). Do not rely on the sign of $m$ alone to deduce orientation.
\end{attention}

\courssub{Measuring the Focal Length of a Converging Lens}

Three standard laboratory methods are required by the syllabus.

\begin{methode}[Distant Object Method]
\begin{enumerate}
\item Point the lens at a distant object (e.g., a tree or building at least $50\ \mathrm{m}$ away, effectively at infinity).
\item Move a screen behind the lens until a sharp image is formed.
\item Measure the distance from the lens to the screen; this is $f$ (since $u \approx \infty \Rightarrow 1/u \approx 0 \Rightarrow v = f$).
\end{enumerate}
\end{methode}

\begin{methode}[$u$–$v$ Method (No Parallax)]
\begin{enumerate}
\item Place the lens between an illuminated object (e.g., a cross-wire or a filament lamp) and a screen.
\item Adjust the lens position until a sharp image is formed on the screen. Record $u$ and $v$.
\item Repeat for at least five different object distances $u$ (ensuring $u > f$ for real images).
\item Plot a graph of $1/v$ against $1/u$ (or $v$ against $u$). The intercepts give $1/f$.
\end{enumerate}
\end{methode}

\begin{propriete}[Graphical Analysis for $u$–$v$ Method]
From $\frac{1}{f} = \frac{1}{u} + \frac{1}{v}$, we have $\frac{1}{v} = -\frac{1}{u} + \frac{1}{f}$.
A plot of $1/v$ (vertical) vs $1/u$ (horizontal) is a straight line with:
\begin{itemize}
\item Slope $= -1$
\item $y$-intercept $= 1/f$
\item $x$-intercept $= 1/f$
\end{itemize}
Alternatively, a plot of $v$ vs $u$ is a rectangular hyperbola; the line $v=u$ cuts it at $u=v=2f$.
\end{propriete}

\begin{methode}[Displacement Method (Alternative)]
If the distance $D$ between object and screen is fixed and $D > 4f$, there are two lens positions that give a sharp image (one magnified, one diminished). If the distance between these two lens positions is $d$, then
\[
f = \frac{D^2 - d^2}{4D}
\]
This method is quick and accurate but requires $D > 4f$.
\end{methode}

\courssub{Refraction Through a Prism}

A prism is a transparent block with two plane refracting surfaces inclined at an angle $A$ (the \cle{angle of the prism} or \cle{refracting angle}). A ray passing through the prism is deviated by an angle $D$ (the \cle{angle of deviation}).

\begin{popfigure}
\begin{tikzpicture}[scale=1.3, >=stealth]
% Prism
\draw[popblue, thick] (0,0) -- (3,0) -- (1.5,2.6) -- cycle;
% Normal lines
\draw[popmuted, dashed] (1.5,0) -- (1.5,1.3);
\draw[popmuted, dashed] (2.25,1.3) -- (0.75,1.3); % approximate normal to second face
% Incident ray
\draw[popgreen, ->] (0.5,-1) -- (1.5,0) node[midway, below right] {$i$};
% Refracted ray inside
\draw[popgreen, ->] (1.5,0) -- (2.25,1.3) node[midway, right] {$r_1$};
% Emergent ray
\draw[popgreen, ->] (2.25,1.3) -- (3.5,1.8) node[midway, above right] {$e$};
% Deviation angle
\draw[poporange, ->] (3.5,1.8) -- (3.5,0.5) node[midway, right] {$D$};
\draw[poporange, ->] (0.5,-1) -- (3.5,-1) node[midway, below] {Original direction};
% Angle A
\draw[poppurple, <->] (1.5,0.3) arc (0:60:0.3) node[midway, right] {$A$};
% Labels
\node[popblue] at (1.5, -0.5) {Prism};
\node[popdark] at (1.5,0) [below left] {O};
\end{tikzpicture}
\legende{A ray passing through a prism. The angle of deviation $D = i + e - A$. At minimum deviation, the path is symmetric: $i = e$ and $r_1 = r_2 = A/2$.}
\end{popfigure}

For a given prism and wavelength, $D$ varies with the angle of incidence $i$. There is a unique \cle{minimum deviation} $D_{\min}$ when the ray passes symmetrically through the prism ($i = e$, $r_1 = r_2 = A/2$). At this point, the refractive index $n$ of the prism material (relative to air) is given by:

\begin{propriete}[Refractive Index at Minimum Deviation]
\[
n = \frac{\sin\left(\frac{A + D_{\min}}{2}\right)}{\sin\left(\frac{A}{2}\right)}
\]
This formula is examinable and must be derived from Snell's law and geometry ($A = r_1 + r_2$, $D = i + e - A$).
\end{propriete}

\begin{exoresolu}[Minimum Deviation Calculation]
A prism of angle $A = 60^\circ$ produces a minimum deviation $D_{\min} = 40^\circ$. Find the refractive index of the glass.
\tcblower
\[
n = \frac{\sin\left(\frac{60^\circ + 40^\circ}{2}\right)}{\sin\left(\frac{60^\circ}{2}\right)}
= \frac{\sin 50^\circ}{\sin 30^\circ}
= \frac{0.7660}{0.5}
= 1.53
\]
\end{exoresolu}

\courssub{The Compound Microscope}

The compound microscope uses two converging lenses: the \cle{objective} (short focal length $f_o$, high power) and the \cle{eyepiece} (longer focal length $f_e$, lower power). The objective forms a real, inverted, magnified image $I_1$ at its focal plane (or just beyond). This image acts as the object for the eyepiece, which functions as a simple magnifier to produce a final virtual image at the near point ($D = 25\ \mathrm{cm}$) or at infinity (normal adjustment).

\begin{popfigure}
\begin{tikzpicture}[scale=0.9, >=stealth]
% Objective lens
\draw[popblue, thick] (0,1.5) .. controls (0.3,1.5) and (0.3,-1.5) .. (0,-1.5);
\draw[popblue, thick] (0,1.5) .. controls (-0.3,1.5) and (-0.3,-1.5) .. (0,-1.5);
\fill[popdark] (0,0) circle (1.5pt) node[below] {$O_o$};
\fill[poporange] (-1.5,0) circle (1.5pt) node[below] {$F_o$};
\fill[poporange] (1.5,0) circle (1.5pt) node[below] {$F'_o$};
% Eyepiece lens
\draw[popblue, thick] (8,1.5) .. controls (8.3,1.5) and (8.3,-1.5) .. (8,-1.5);
\draw[popblue, thick] (8,1.5) .. controls (7.7,1.5) and (7.7,-1.5) .. (8,-1.5);
\fill[popdark] (8,0) circle (1.5pt) node[below] {$O_e$};
\fill[poporange] (6.5,0) circle (1.5pt) node[below] {$F_e$};
\fill[poporange] (9.5,0) circle (1.5pt) node[below] {$F'_e$};
% Object
\draw[popdark, thick] (-2.5,0) -- (-2.5,0.8) node[above] {Object};
% Rays from object through objective
\draw[popgreen, ->] (-2.5,0.8) -- (0,0.8);
\draw[popgreen, ->] (0,0.8) -- (1.5,0);
\draw[popgreen, ->] (-2.5,0.8) -- (0,0) -- (4,1.28);
% Intermediate image I1
\draw[popred, thick] (4,0) -- (4,1.28) node[above] {$I_1$};
% Rays from I1 through eyepiece (normal adjustment: I1 at Fe)
\draw[popgreen, ->] (4,1.28) -- (8,1.28);
\draw[popgreen, ->] (8,1.28) -- (10,1.28); % parallel emergence
\draw[popgreen, ->] (4,1.28) -- (8,0) -- (10,-1.28);
% Final virtual image at infinity (parallel rays)
\draw[popred, dashed] (10,1.28) -- (12,1.28);
\draw[popred, dashed] (10,-1.28) -- (12,-1.28);
\node[popred] at (11,0) {Final image at $\infty$};
% Length
\draw[popmuted, <->] (0,-2) -- (8,-2) node[midway, below] {Length $L \approx v_o + f_e$};
\end{tikzpicture}
\legende{Ray diagram of a compound microscope in normal adjustment (final image at infinity). The objective forms a real intermediate image $I_1$ at the focal point of the eyepiece $F_e$.}
\end{popfigure}

\begin{propriete}[Angular Magnification of Compound Microscope]
In \cle{normal adjustment} (final image at infinity, relaxed eye), the angular magnification $M$ is:
\[
M = \frac{v_o}{f_o} \times \frac{D}{f_e}
\]
where $v_o$ is the image distance for the objective (approximately the tube length $L$ minus $f_e$), $f_o$ and $f_e$ are the focal lengths of objective and eyepiece, and $D = 25\ \mathrm{cm}$ is the least distance of distinct vision. If the final image is at the near point, the eyepiece magnification becomes $(1 + D/f_e)$.
\end{propriete}

\courssub{The Astronomical Telescope}

The astronomical telescope also uses two converging lenses: the \cle{objective} (large focal length $f_o$, large aperture to gather light) and the \cle{eyepiece} (short focal length $f_e$). The objective forms a real, inverted, diminished image of a distant object at its focal plane. The eyepiece then magnifies this image.

\begin{popfigure}
\begin{tikzpicture}[scale=0.9, >=stealth]
% Objective lens
\draw[popblue, thick] (0,2) .. controls (0.3,2) and (0.3,-2) .. (0,-2);
\draw[popblue, thick] (0,2) .. controls (-0.3,2) and (-0.3,-2) .. (0,-2);
\fill[popdark] (0,0) circle (1.5pt) node[below] {$O_o$};
\fill[poporange] (-3,0) circle (1.5pt) node[below] {$F_o$};
\fill[poporange] (3,0) circle (1.5pt) node[below] {$F'_o$};
% Eyepiece lens
\draw[popblue, thick] (8,2) .. controls (8.3,2) and (8.3,-2) .. (8,-2);
\draw[popblue, thick] (8,2) .. controls (7.7,2) and (7.7,-2) .. (8,-2);
\fill[popdark] (8,0) circle (1.5pt) node[below] {$O_e$};
\fill[poporange] (5,0) circle (1.5pt) node[below] {$F_e$};
\fill[poporange] (11,0) circle (1.5pt) node[below] {$F'_e$};
% Distant object: parallel rays at angle alpha
\draw[popgreen, ->] (-4,1.5) -- (0,1.5);
\draw[popgreen, ->] (-4,0) -- (0,0);
\draw[popgreen, ->] (-4,-1.5) -- (0,-1.5);
% Rays through objective converge at Fo'
\draw[popgreen, ->] (0,1.5) -- (3,0);
\draw[popgreen, ->] (0,0) -- (3,0);
\draw[popgreen, ->] (0,-1.5) -- (3,0);
% Intermediate image at Fo' (which coincides with Fe)
\draw[popred, thick] (3,-0.5) -- (3,0.5) node[right] {$I_1$};
% Rays through eyepiece emerge parallel (normal adjustment)
\draw[popgreen, ->] (3,0.5) -- (8,0.5) -- (10,0.5);
\draw[popgreen, ->] (3,0) -- (8,0) -- (10,0);
\draw[popgreen, ->] (3,-0.5) -- (8,-0.5) -- (10,-0.5);
% Angles
\draw[poppurple, <->] (0,0.2) arc (0:26.5:0.2) node[midway, right] {$\alpha$};
\draw[poppurple, <->] (8,-0.2) arc (-180:-153.5:0.2) node[midway, left] {$\beta$};
% Length
\draw[popmuted, <->] (0,-2.5) -- (8,-2.5) node[midway, below] {Length $L = f_o + f_e$};
\end{tikzpicture}
\legende{Astronomical telescope in normal adjustment. The final image is at infinity. The angular magnification $M = f_o/f_e$.}
\end{popfigure}

\begin{propriete}[Astronomical Telescope in Normal Adjustment]
\begin{itemize}
\item The objective forms a real image at its focal plane $F'_o$.
\item The eyepiece is placed so that its focal plane $F_e$ coincides with $F'_o$ (so the final image is at infinity).
\item The distance between the lenses (length of telescope) is $L = f_o + f_e$.
\item The \cle{angular magnification} is
\[
M = \frac{\beta}{\alpha} = \frac{f_o}{f_e}
\]
where $\alpha$ is the angle subtended by the object at the objective (and at the unaided eye), and $\beta$ is the angle subtended by the final image at the eye.
\end{itemize}
\end{propriete}

\begin{exoresolu}[Telescope Design]
An astronomical telescope has an objective of focal length $120\ \mathrm{cm}$ and an eyepiece of focal length $5\ \mathrm{cm}$. Calculate (a) the angular magnification in normal adjustment, (b) the length of the telescope, (c) the magnification if the final image is formed at the near point ($D=25\ \mathrm{cm}$).
\tcblower
(a) Normal adjustment: $M = \dfrac{f_o}{f_e} = \dfrac{120}{5} = 24$.
(b) Length $L = f_o + f_e = 120 + 5 = 125\ \mathrm{cm}$.
(c) Final image at near point: The eyepiece acts as a magnifier with magnification $m_e = 1 + \dfrac{D}{f_e} = 1 + \dfrac{25}{5} = 6$. The objective still gives a linear magnification of $v_o/u_o \approx f_o/u_o$ (but angular magnification is defined differently). For a telescope with image at near point, the angular magnification is
\[
M = \frac{f_o}{f_e}\left(1 + \frac{f_e}{D}\right) = 24 \times \left(1 + \frac{5}{25}\right) = 24 \times 1.2 = 28.8.
\]
(Alternatively, $M = \frac{f_o}{f_e} + \frac{f_o}{D} = 24 + \frac{120}{25} = 24 + 4.8 = 28.8$.)
\end{exoresolu}

\courssub{Summary of Key Formulae}

\begin{aretenir}[Lenses, Prisms and Optical Instruments]
\begin{itemize}
\item \textbf{Thin lens formula:} $\dfrac{1}{f} = \dfrac{1}{u} + \dfrac{1}{v}$ (real-is-positive convention).
\item \textbf{Linear magnification:} $m = \dfrac{v}{u}$ (sign indicates virtual/real, not orientation; state orientation separately).
\item \textbf{Power of a lens:} $P = \dfrac{1}{f(\mathrm{m})}$ in dioptres (D). Converging: $P>0$; Diverging: $P<0$.
\item \textbf{Prism at minimum deviation:} $n = \dfrac{\sin\frac{A+D_{\min}}{2}}{\sin\frac{A}{2}}$.
\item \textbf{Compound microscope (normal adjustment):} $M = \dfrac{v_o}{f_o} \times \dfrac{D}{f_e}$; length $L \approx v_o + f_e$.
\item \textbf{Astronomical telescope (normal adjustment):} $M = \dfrac{f_o}{f_e}$; length $L = f_o + f_e$.
\item \textbf{Astronomical telescope (image at near point):} $M = \dfrac{f_o}{f_e}\left(1 + \dfrac{f_e}{D}\right) = \dfrac{f_o}{f_e} + \dfrac{f_o}{D}$.
\end{itemize}
\end{aretenir}

\begin{exercice}[Practice Problems]
\begin{enumerate}
\item A converging lens of power $+6.0\ \mathrm{D}$ is used as a magnifying glass. An object is placed $10\ \mathrm{cm}$ from the lens. Find the image distance, magnification, and describe the image.
\item A prism of angle $60^\circ$ and refractive index $1.5$ is used in a spectrometer. Calculate the angle of minimum deviation.
\item A compound microscope has an objective of focal length $0.5\ \mathrm{cm}$ and an eyepiece of focal length $2.5\ \mathrm{cm}$. The tube length (distance between objective image and eyepiece) is $16\ \mathrm{cm}$. Calculate the angular magnification in normal adjustment.
\item An astronomical telescope has an objective of focal length $100\ \mathrm{cm}$ and an eyepiece of focal length $4\ \mathrm{cm}$. What is the length of the telescope in normal adjustment? If the final image is formed at the near point ($25\ \mathrm{cm}$), what is the magnification?
\end{enumerate}
\end{exercice}

\begin{corrige}[Exercise 1]
$P = +6.0\ \mathrm{D} \Rightarrow f = 1/6 = 0.1667\ \mathrm{m} = 16.67\ \mathrm{cm}$.
$u = +10\ \mathrm{cm}$.
$\frac{1}{v} = \frac{1}{f} - \frac{1}{u} = \frac{1}{16.67} - \frac{1}{10} = 0.06 - 0.1 = -0.04\ \mathrm{cm}^{-1}$.
$v = -25\ \mathrm{cm}$. Negative $\Rightarrow$ virtual image on same side as object.
$m = v/u = -25/10 = -2.5$. Magnified ($|m|>1$), upright (virtual).
\end{corrige}

\begin{corrige}[Exercise 2]
$n = 1.5$, $A = 60^\circ$.
$1.5 = \frac{\sin\frac{60+D_{\min}}{2}}{\sin 30^\circ} = 2 \sin\frac{60+D_{\min}}{2}$.
$\sin\frac{60+D_{\min}}{2} = 0.75$.
$\frac{60+D_{\min}}{2} = \sin^{-1}(0.75) \approx 48.59^\circ$.
$60 + D_{\min} = 97.18^\circ \Rightarrow D_{\min} \approx 37.2^\circ$.
\end{corrige}

\begin{corrige}[Exercise 3]
Normal adjustment: $v_o = \text{tube length} = 16\ \mathrm{cm}$ (since intermediate image at $F_e$).
$M = \frac{v_o}{f_o} \times \frac{D}{f_e} = \frac{16}{0.5} \times \frac{25}{2.5} = 32 \times 10 = 320$.
\end{corrige}

\begin{corrige}[Exercise 4]
Normal adjustment: $L = f_o + f_e = 100 + 4 = 104\ \mathrm{cm}$.
$M_{\text{normal}} = f_o/f_e = 100/4 = 25$.
Image at near point: $M = \frac{f_o}{f_e} + \frac{f_o}{D} = 25 + \frac{100}{25} = 25 + 4 = 29$.
\end{corrige}

\coursec{Section 6: Energy of Waves in Free Space and the Inverse Square Law}

In Section 5 we studied how lenses, prisms and optical instruments redirect light to form useful images. But whether light passes through a telescope objective or travels freely from a star to your eye, one fundamental question remains: \emph{how much energy does the wave actually deliver, and how does that amount change with distance from the source?} Anyone who has walked away from a loudspeaker at a football match in Yaound\'e, or noticed that a kerosene lamp looks dimmer from across the courtyard, has already met the answer qualitatively: the effect of a wave \emph{weakens with distance}. In this section we make that idea precise, deriving the celebrated \cle{inverse square law} and connecting it to the amplitude of the wave.

\courssub{Energy Transport by Waves and the Concept of Intensity}

\begin{experience}[A simple observation]
Stand a few metres from a radio playing at constant volume. The sound is loud. Now walk to the far end of the school field: the same radio, producing the same power of sound, is now barely audible. Nothing about the source has changed --- yet what reaches your ear has changed dramatically. The energy emitted by the speaker has been \emph{spread out} over a larger region of space. Understanding this spreading is the goal of this section.
\end{experience}

A wave is, at its heart, a mechanism for \textbf{transporting energy without transporting matter}. The ripples on a pond carry energy outward from the splash, sound carries energy from the loudspeaker to your eardrum, and light carries energy from the Sun to the Earth across about \SI{1.5e8}{km} of empty space. To compare sources and detectors quantitatively, we need a quantity that describes \emph{how concentrated} this energy flow is at a given point. That quantity is the \cle{intensity}.

\begin{definition}[Intensity of a wave]
The \cle{intensity} $I$ of a wave at a point is the \textbf{power received per unit area} measured perpendicular to the direction of propagation of the wave:
\[
I = \frac{P}{A}
\]
where $P$ is the power (energy per unit time, in watts) crossing a surface of area $A$ placed at right angles to the wave. The SI unit of intensity is the \textbf{watt per square metre}, $\mathrm{W\,m^{-2}}$.
\end{definition}

\begin{attention}[The area must be perpendicular]
The area $A$ in $I = P/A$ is the area of a surface held \emph{perpendicular} to the direction in which the wave travels. If you tilt a solar panel away from the Sun, it intercepts less power not because the intensity of sunlight has changed, but because the \emph{effective} area it presents to the beam has decreased.
\end{attention}

\begin{exemplebox}[Intensity of sunlight on a small panel]
A solar panel of area $A = \SI{2.0}{m^2}$ is placed perpendicular to sunlight and receives power at a rate of $P = \SI{800}{W}$. The intensity of the sunlight at the panel is
\[
I = \frac{P}{A} = \frac{800}{2.0} = \SI{400}{W.m^{-2}}.
\]
\end{exemplebox}

\courssub{A Point Source Radiating Uniformly in Free Space}

Consider a small source --- a lamp filament, a star, a compact loudspeaker --- that radiates energy \textbf{uniformly in all directions} into a homogeneous medium with \textbf{no absorption}. Such an idealised source is called an \emph{isotropic point source}, and the surrounding medium with no losses is called \cle{free space}.

Because the radiation is uniform in all directions, the wavefronts are \textbf{spheres} centred on the source. As the wave travels outward, each spherical wavefront expands: at distance $r$ from the source, the energy is spread over the surface of a sphere of radius $r$, whose area is
\[
A = 4\pi r^2.
\]

\begin{propriete}[Conservation of energy for waves in free space]
For a point source radiating in free space (no absorption, no reflection, no obstacles), energy is conserved: \textbf{the same total power $P$ crosses every sphere centred on the source}, whatever its radius. No energy is created or lost as the wavefront expands; it is simply distributed over a larger and larger area.
\end{propriete}

This conservation statement is the physical heart of everything that follows. Imagine counting the energy crossing a small sphere around the source each second; since nothing absorbs the wave, exactly the same amount of energy must cross a larger sphere each second --- otherwise energy would be piling up in the space between, which does not happen for a steady source.

\begin{popfigure}
\begin{center}
\begin{tikzpicture}[scale=1.0]
  % source
  \fill[poporange] (0,0) circle (0.09);
  \node[font=\small, align=center] at (0,-0.45) {point source, power $P$};
  % inner sphere (circle, radius r)
  \draw[popblue, thick] (0,0) circle (1.2);
  % outer sphere (circle, radius 2r)
  \draw[poppurple, thick] (0,0) circle (2.4);
  % radii
  \draw[->, popblue, thick] (0,0) -- (0.85,0.85) node[midway, above left, font=\small] {$r$};
  \draw[->, poppurple, thick] (0,0) -- (-1.70,-1.70) node[midway, below left, font=\small] {$2r$};
  % rays
  \draw[->, popgold, thick] (0,0) -- (2.9,0.6);
  \draw[->, popgold, thick] (0,0) -- (-0.6,2.9);
  \draw[->, popgold, thick] (0,0) -- (0.6,-2.9);
  \draw[->, popgold, thick] (0,0) -- (-2.9,-0.6);
  % labels of areas
  \node[popblue, font=\small, align=center] at (1.9,1.7) {area $4\pi r^{2}$};
  \node[poppurple, font=\small, align=center] at (-2.7,1.5) {area $4\pi (2r)^{2}$\\ $= 4 \times 4\pi r^{2}$};
  % intensity labels
  \node[popblue, font=\small] at (0.0,-1.55) {$I_1 = \dfrac{P}{4\pi r^{2}}$};
  \node[poppurple, font=\small] at (0.0,-2.95) {$I_2 = \dfrac{P}{4\pi (2r)^{2}} = \dfrac{I_1}{4}$};
\end{tikzpicture}
\legende{A point source radiating uniformly. The same power $P$ crosses the sphere of radius $r$ and the sphere of radius $2r$, but on the larger sphere it is spread over four times the area, so the intensity is four times smaller.}
\end{center}
\end{popfigure}

\courssub{Derivation of the Inverse Square Law}

The derivation is short and \textbf{examinable} --- you should be able to reproduce it in three lines.

\begin{enumerate}
\item The source emits power $P$ uniformly in all directions.
\item At distance $r$, this power is spread uniformly over a sphere of area $A = 4\pi r^{2}$ (conservation of energy: all of $P$ crosses the sphere).
\item Applying the definition of intensity:
\[
I = \frac{P}{A} = \frac{P}{4\pi r^{2}}.
\]
\end{enumerate}

\begin{propriete}[Inverse square law]
For a point source radiating uniformly in free space, the intensity at distance $r$ from the source is
\[
\boxed{\,I = \frac{P}{4\pi r^{2}}\,} \qquad\text{i.e.}\qquad I \propto \frac{1}{r^{2}}.
\]
Doubling the distance divides the intensity by four; tripling it divides the intensity by nine. (The derivation above is examinable.)
\end{propriete}

A very useful consequence follows when we compare the intensities $I_1$ and $I_2$ at two distances $r_1$ and $r_2$ from the \emph{same} source. Since $I_1 = P/(4\pi r_1^{2})$ and $I_2 = P/(4\pi r_2^{2})$, multiplying each by the square of its distance eliminates $P$ and $4\pi$:
\[
I_1 r_1^{2} = \frac{P}{4\pi} = I_2 r_2^{2}.
\]

\begin{aretenir}[The two-distance relation]
For the same point source measured at two different distances:
\[
I_1\, r_1^{2} = I_2\, r_2^{2}.
\]
This relation lets you solve problems without knowing the power $P$ of the source at all.
\end{aretenir}

\begin{popfigure}
\begin{center}
\begin{tikzpicture}
\begin{axis}[
    width=0.78\linewidth,
    height=6.2cm,
    xlabel={distance $r$ (m)},
    ylabel={intensity $I$ ($\mathrm{W\,m^{-2}}$)},
    xmin=0, xmax=5.2,
    ymin=0, ymax=105,
    xtick={1,2,3,4,5},
    ytick={0,25,50,75,100},
    grid=both,
    grid style={popline, very thin},
    axis line style={popink},
    label style={font=\small},
    tick label style={font=\small},
    legend style={font=\small, at={(0.97,0.95)}, anchor=north east},
]
\addplot[domain=0.32:5.2, samples=120, thick, popblue] {100/x^2};
\addlegendentry{$I = \dfrac{P}{4\pi r^{2}}$}
% marked points (1,100) and (2,25)
\addplot[only marks, mark=*, mark size=2.2pt, poporange] coordinates {(1,100) (2,25)};
\node[font=\small, popdark, anchor=west] at (axis cs:1.08,100) {$(r,\; I)$};
\node[font=\small, popdark, anchor=west] at (axis cs:2.08,25) {$(2r,\; I/4)$};
% dashed guide lines
\draw[dashed, popmuted] (axis cs:1,0) -- (axis cs:1,100);
\draw[dashed, popmuted] (axis cs:2,0) -- (axis cs:2,25);
\draw[dashed, popmuted] (axis cs:0,25) -- (axis cs:2,25);
\end{axis}
\end{tikzpicture}
\legende{Intensity versus distance for a point source (here $P/(4\pi) = \SI{100}{W}$). When the distance doubles from $r$ to $2r$, the intensity falls to one quarter of its value.}
\end{center}
\end{popfigure}

\courssub{Intensity and Amplitude}

Intensity describes the \emph{energy flow}; amplitude describes the \emph{size of the oscillation}. The two are linked by a general result valid for all types of wave (sound, light, waves on strings):

\begin{propriete}[Intensity is proportional to amplitude squared]
For any wave, the intensity is proportional to the square of the amplitude $a$ of the oscillation:
\[
I \propto a^{2}.
\]
Physically, the energy carried by a wave depends on how \emph{strongly} the medium (or the electromagnetic field) oscillates, and this dependence is quadratic --- just as the energy stored in a stretched spring is $\tfrac12 k x^{2}$, proportional to the square of the displacement.
\end{propriete}

Combining $I \propto a^{2}$ with the inverse square law $I \propto 1/r^{2}$ for a point source gives a result that surprises many students:
\[
a^{2} \propto \frac{1}{r^{2}} \quad\Longrightarrow\quad \boxed{\,a \propto \frac{1}{r}\,}.
\]

\begin{attention}[Do not confuse intensity and amplitude]
For a point source in free space:
\begin{itemize}
\item \textbf{Intensity} falls off as $I \propto \dfrac{1}{r^{2}}$ (inverse \emph{square} law);
\item \textbf{Amplitude} falls off as $a \propto \dfrac{1}{r}$ (simple inverse law).
\end{itemize}
A classic GCE trap asks: ``the distance from a point source is doubled; what happens to the amplitude?'' The answer is that the amplitude is \emph{halved}, while the \emph{intensity} is divided by four. Mixing these two up is one of the most common errors in this topic.
\end{attention}

\courssub{Worked Examples and the Two-Distance Method}

\begin{methode}[Solving two-distance problems with $I_1 r_1^{2} = I_2 r_2^{2}$]
\begin{enumerate}
\item \textbf{Check the assumptions}: the source must behave as a point source radiating uniformly, with negligible absorption between source and detector.
\item \textbf{Identify} the two distances $r_1$, $r_2$ and the known intensity (or the unknown one).
\item \textbf{Write} $I_1 r_1^{2} = I_2 r_2^{2}$ and rearrange for the unknown:
\[
I_2 = I_1 \left(\frac{r_1}{r_2}\right)^{2}.
\]
\item \textbf{Compute}, keeping units consistent (both distances in the same unit) and quoting the intensity in $\mathrm{W\,m^{-2}}$ with sensible significant figures.
\item \textbf{Sanity check}: if $r_2 > r_1$, then $I_2$ must be \emph{smaller} than $I_1$.
\end{enumerate}
\end{methode}

\begin{exoresolu}[Brightness of a lamp at two distances]
A small lamp radiates uniformly in all directions. A light meter measures an intensity $I_1 = \SI{5.0}{W.m^{-2}}$ at a distance $r_1 = \SI{2.0}{m}$ from the lamp. (a) Calculate the intensity at $r_2 = \SI{6.0}{m}$. (b) At what distance would the intensity be $I_3 = \SI{0.20}{W.m^{-2}}$? (c) Deduce the total power $P$ emitted by the lamp as light.
\tcblower
\textbf{(a)} Treating the lamp as a point source in free space, apply the two-distance relation:
\[
I_2 = I_1\left(\frac{r_1}{r_2}\right)^{2} = 5.0 \times \left(\frac{2.0}{6.0}\right)^{2} = 5.0 \times \frac{1}{9} \approx \SI{0.56}{W.m^{-2}}.
\]
The distance was tripled, so the intensity is divided by $3^{2} = 9$. \checkmark

\textbf{(b)} Rearranging $I_1 r_1^{2} = I_3 r_3^{2}$:
\[
r_3 = r_1\sqrt{\frac{I_1}{I_3}} = 2.0 \times \sqrt{\frac{5.0}{0.20}} = 2.0 \times \sqrt{25} = \SI{10}{m}.
\]

\textbf{(c)} Using $I = P/(4\pi r^{2})$ at the first position:
\[
P = 4\pi r_1^{2} I_1 = 4\pi \times (2.0)^{2} \times 5.0 \approx \SI{2.5e2}{W}.
\]
Note that this is the power carried by the light alone; a real lamp also dissipates heat.
\end{exoresolu}

\begin{exoresolu}[Solar radiation at Earth and at Mars]
The Sun radiates a total power $P = \SI{3.9e26}{W}$ uniformly in all directions. The mean distance of the Earth from the Sun is $r_E = \SI{1.5e11}{m}$ and that of Mars is $r_M = \SI{2.3e11}{m}$. (a) Calculate the solar intensity at the Earth. (b) Calculate the solar intensity at Mars using the two-distance relation.
\tcblower
\textbf{(a)} Direct application of the inverse square law:
\[
I_E = \frac{P}{4\pi r_E^{2}} = \frac{3.9\times10^{26}}{4\pi \times (1.5\times10^{11})^{2}}
= \frac{3.9\times10^{26}}{2.83\times10^{23}} \approx \SI{1.4e3}{W.m^{-2}}.
\]
This value, about \SI{1.4}{kW.m^{-2}}, is known as the \emph{solar constant} --- the reason solar panels are so effective in sun-rich regions of Cameroon.

\textbf{(b)} Using $I_M r_M^{2} = I_E r_E^{2}$:
\[
I_M = I_E\left(\frac{r_E}{r_M}\right)^{2} = 1.4\times10^{3}\times\left(\frac{1.5}{2.3}\right)^{2}
\approx 1.4\times10^{3}\times 0.43 \approx \SI{6.0e2}{W.m^{-2}}.
\]
Mars receives less than half the solar intensity available at Earth, which is why solar-powered missions to Mars need much larger panels.
\end{exoresolu}

\courssub{Everyday and Technological Applications}

\begin{itemize}
\item \textbf{Brightness of a lamp.} A lamp that appears comfortably bright at \SI{1}{m} delivers only a quarter of that intensity at \SI{2}{m}. This is why reading lamps must be placed close to the page.
\item \textbf{Loudness of a speaker.} The sound intensity from a loudspeaker at an open-air rally falls roughly as $1/r^{2}$ (while the ground and buildings are far away). Sound engineers position multiple speakers along a field precisely because a single speaker cannot deliver adequate intensity to distant listeners.
\item \textbf{Radio and television signals.} The intensity of the signal from a transmitter mast decreases as $1/r^{2}$, which is why reception is poor in villages far from the mast and why relay stations are needed to cover the national territory.
\item \textbf{Astronomy.} The apparent brightness of a star is the intensity of its light at the Earth. Comparing the apparent brightness of two stars of known similar power allows astronomers to compare their distances using $I_1 r_1^{2} = I_2 r_2^{2}$ --- a cornerstone of distance measurement in the Universe.
\end{itemize}

\begin{savaistu}[Why distant stars are invisible to the naked eye]
The human eye detects light only above a certain threshold intensity. A star may emit as much power as the Sun, but if it is a million times farther away, the inverse square law reduces its intensity by a factor of $10^{12}$ --- far below what the eye can detect. Telescopes help not by ``magnifying brightness'' but by collecting light over a large objective area $A$, gathering a power $P = IA$ large enough to register.
\end{savaistu}

\courssub{Limitations of the Inverse Square Law}

The law $I = P/(4\pi r^{2})$ rests on idealised assumptions. Knowing its limits is excellent GCE extension material.

\begin{attention}[When the inverse square law fails or needs care]
\begin{itemize}
\item \textbf{Real sources are not perfect points.} Very close to an extended source (a long fluorescent tube, a large loudspeaker), the $1/r^{2}$ law does not apply; it becomes accurate only at distances large compared with the size of the source.
\item \textbf{Sources are not always isotropic.} A torch with a reflector or a radio antenna radiates preferentially in certain directions, so the power is not spread uniformly over $4\pi r^{2}$.
\item \textbf{Absorption and scattering.} Fog, dust, walls and the atmosphere absorb or scatter energy, so the intensity falls \emph{faster} than $1/r^{2}$. Sound in a hall is also reflected, which can make the intensity fall \emph{slower} than predicted.
\item \textbf{Beams.} A laser produces an approximately parallel beam: the energy is not spread over expanding spheres, and the intensity stays nearly constant over long distances.
\end{itemize}
\end{attention}

\begin{methode}[Exam tip (GCE): state your assumptions]
Whenever you use $I = P/(4\pi r^{2})$ or $I_1 r_1^{2} = I_2 r_2^{2}$ in an examination, begin by writing the assumptions explicitly: \emph{``treating the source as a point source radiating uniformly in all directions, with no absorption between source and detector.''} Examiners award marks for this statement, and it protects you when the question later asks why real measurements deviate from the prediction.
\end{methode}

\begin{exercice}[Loudspeaker at a school rally]
A loudspeaker radiates sound uniformly in all directions with a power of \SI{30}{W}. (a) Calculate the sound intensity at distances of \SI{5.0}{m} and \SI{20}{m} from the speaker. (b) A student finds the sound ``just comfortable'' at \SI{5.0}{m}. At what distance would the intensity be one hundredth of that value? (c) Give two reasons why the measured intensity at \SI{20}{m} in a real schoolyard might differ from your calculated value.
\end{exercice}

\begin{corrige}[Exercise 1]
\textbf{(a)} At $r_1 = \SI{5.0}{m}$:
\[
I_1 = \frac{P}{4\pi r_1^{2}} = \frac{30}{4\pi \times 25} \approx \SI{9.5e-2}{W.m^{-2}}.
\]
At $r_2 = \SI{20}{m}$ the distance is four times greater, so the intensity is $4^{2} = 16$ times smaller:
\[
I_2 = \frac{I_1}{16} \approx \SI{6.0e-3}{W.m^{-2}}.
\]
\textbf{(b)} We require $I_3 = I_1/100$. Since $I \propto 1/r^{2}$, the distance must be $\sqrt{100} = 10$ times larger: $r_3 = 10 \times 5.0 = \SI{50}{m}$.
\textbf{(c)} Any two of: absorption of sound by the air; reflection from the ground and nearby buildings adding to the direct sound; the speaker not radiating uniformly (it is directional); wind and obstacles scattering the sound.
\end{corrige}

\begin{aretenir}[Key points of Section 6]
\begin{itemize}
\item Intensity is power per unit perpendicular area: $I = P/A$, in $\mathrm{W\,m^{-2}}$.
\item A point source radiating uniformly in free space spreads its power over spheres of area $4\pi r^{2}$; conservation of energy then gives the \cle{inverse square law} $I = \dfrac{P}{4\pi r^{2}}$.
\item For the same source at two distances: $I_1 r_1^{2} = I_2 r_2^{2}$.
\item Intensity is proportional to the square of the amplitude, $I \propto a^{2}$; hence for a point source the \emph{amplitude} falls as $a \propto 1/r$ while the \emph{intensity} falls as $1/r^{2}$.
\item The law assumes a point source, uniform radiation and no absorption --- state these assumptions in the examination.
\end{itemize}
\end{aretenir}

\newpage

\coursec{Integration activity}

\courssub{Aims of the activity}

This integration activity closes the chapter by asking you to act as a \cle{physics consultant} for a real school. By the end of the activity, you should be able to:

\begin{itemize}
\item apply \cle{Snell's law} and the condition for \cle{total internal reflection} to compute a critical angle and an acceptance angle for an optical fibre;
\item use a \cle{diffraction grating} relation $d\sin\theta = n\lambda$ to identify the wavelength of a laser and confirm that it is monochromatic;
\item use the \cle{thin-lens equation} and \cle{linear magnification} to select a projector lens;
\item apply the \cle{inverse square law} $I = \dfrac{P}{4\pi r^{2}}$ to design a lighting plan;
\item organise a group investigation, present numerical results with correct SI units and sensible significant figures, and justify every physical principle invoked.
\end{itemize}

\begin{attention}[Skills being assessed]
This activity mobilises Lessons 21, 24, 28, 29, 32 and 33 simultaneously, exactly as the GCE Advanced Level problem questions do. Marks are awarded not only for correct numbers but for \emph{stating the law used}, \emph{showing the substitution}, and \emph{giving the unit}.
\end{attention}

\begin{methode}[Approach to multi-part physics problems]
\begin{enumerate}
\item \textbf{Identify the physics:} For each part, name the relevant law or principle (e.g. Snell's law, grating equation, thin-lens formula, inverse square law).
\item \textbf{Draw a diagram:} Sketch the geometry (fibre cross-section, grating, lens, room plan). Label known quantities and unknowns.
\item \textbf{Write the equation:} State the formula in symbols before substituting numbers.
\item \textbf{Substitute with units:} Use SI units throughout; carry units through the calculation.
\item \textbf{Calculate and round:} Give the final answer to 2 or 3 significant figures, matching the data precision, with the correct unit.
\item \textbf{Interpret:} Explain what the number means in the context of the problem and make a clear recommendation.
\end{enumerate}
\end{methode}

\courssub{Situation}

\textbf{The Bamenda College of Science project.}\par
The Bamenda College of Science, a secondary school on the hills of the North-West Region, has received a donation of equipment and wants to modernise two facilities. The Parent-Teacher Association has hired your Upper Sixth physics class as consultants before spending any money.

\textbf{Part A — Connecting the computer lab.}\par
The administration block, where the internet router is installed, is separated from the new computer laboratory by a courtyard. A technician proposes to run a \cle{step-index optical fibre} between the two buildings. The fibre core is made of glass of refractive index $n_{1} = 1.48$, surrounded by a cladding of refractive index $n_{2} = 1.46$. The fibre will be bent around corners along its path, so light must be guided by total internal reflection at the core--cladding boundary. To test the installed fibre, the technician shines light from a small laser into one end. The laser's wavelength is unknown: its label has been scratched off, and the school suspects it may not even be monochromatic.

\begin{popfigure}
\begin{tikzpicture}[scale=1.2]
% Core
\draw[popblue, thick, fill=popblueL] (-2, -0.5) rectangle (2, 0.5);
\node[popdark] at (0, 0) {Core $n_1=1.48$};
% Cladding
\draw[poporange, thick, fill=poporange!20] (-2, -1) rectangle (-2, -0.5);
\draw[poporange, thick, fill=poporange!20] (-2, 0.5) rectangle (2, 1);
\draw[poporange, thick, fill=poporange!20] (-2, -1) rectangle (2, -0.5);
\draw[poporange, thick, fill=poporange!20] (-2, 0.5) rectangle (2, 1);
\node[popdark] at (0, -0.75) {Cladding $n_2=1.46$};
\node[popdark] at (0, 0.75) {Cladding $n_2=1.46$};
% Ray
\draw[popred, very thick, ->] (-1.5, 0) -- (-0.5, 0.4);
\draw[popred, very thick, ->] (-0.5, 0.4) -- (0.5, -0.4);
\draw[popred, very thick, ->] (0.5, -0.4) -- (1.5, 0.4);
% Normal lines
\draw[dashed, thin, popmuted] (-0.5, 0.5) -- (-0.5, -0.5);
\draw[dashed, thin, popmuted] (0.5, 0.5) -- (0.5, -0.5);
% Angles
\draw[popred, <->] (-0.5, 0.1) arc (90:135:0.4) node[midway, left, popred] {$\theta > \theta_c$};
\draw[popred, <->] (0.5, -0.1) arc (-90:-45:0.4) node[midway, right, popred] {$\theta > \theta_c$};
\end{tikzpicture}
\legende{Step-index fibre: light guided by total internal reflection at core-cladding boundary.}
\end{popfigure}

\textbf{Part B — Testing the laser.}\par
The physics department owns a \cle{transmission grating} ruled with $600$ lines per millimetre and a screen. When the laser beam is shone at normal incidence onto the grating, a single sharp bright spot appears on each side of the central maximum at first order, measured at an angle $\theta_{1} = 23.1^{\circ}$ from the straight-through direction. No coloured spreading of the spot is observed.

\begin{popfigure}
\begin{tikzpicture}[scale=1.0]
% Grating
\draw[popdark, very thick] (0, -1.5) -- (0, 1.5);
\node[popdark, left] at (0, 0) {Grating};
% Incident ray
\draw[popblue, thick, ->] (-2, 0) -- (0, 0) node[midway, above] {Laser};
% Central max
\draw[popblue, thick, ->] (0, 0) -- (2, 0) node[midway, above] {$n=0$};
% First order
\draw[popred, thick, ->] (0, 0) -- (2, 0.85) node[midway, above right] {$n=1,\ \theta_1=23.1^\circ$};
\draw[popred, thick, ->] (0, 0) -- (2, -0.85) node[midway, below right] {$n=-1$};
% Angle arc
\draw[popred, <->] (0.3, 0) arc (0:23.1:0.3) node[midway, right=2pt] {$\theta_1$};
% Screen
\draw[popdark, ultra thick] (3, -1.5) -- (3, 1.5);
\node[popdark, right] at (3, 0) {Screen};
% Spots
\fill[popblue] (3, 0) circle (3pt);
\fill[popred] (3, 1.27) circle (3pt);
\fill[popred] (3, -1.27) circle (3pt);
\end{tikzpicture}
\legende{Diffraction grating test: single sharp first-order spots indicate monochromatic light.}
\end{popfigure}

\textbf{Part C — The projector.}\par
For lessons, the school wants a projector that forms a real image of a $5.0\ \mathrm{cm}$ high slide onto a screen placed $4.0\ \mathrm{m}$ from the slide. The image must be $2.0\ \mathrm{m}$ high so that students at the back of the classroom can read it. A single thin converging lens is to be used.

\begin{popfigure}
\begin{tikzpicture}[scale=0.6]
% Lens
\draw[popblue, thick] (0, -1) -- (0, 1);
\draw[popblue, thick] (-0.1, 1) -- (0.1, 1);
\draw[popblue, thick] (-0.1, -1) -- (0.1, -1);
\node[popdark, above] at (0, 1.2) {Lens};
% Object (slide)
\draw[poporange, thick] (-2.5, 0) -- (-2.5, 0.5);
\node[popdark, left] at (-2.5, 0.25) {Slide $h_o=5.0\ \mathrm{cm}$};
\draw[<->, popmuted] (-2.5, -0.5) -- (0, -0.5) node[midway, below] {$u$};
% Image
\draw[popgreen, thick] (3.9, 0) -- (3.9, -2.0);
\node[popdark, right] at (3.9, -1.0) {Image $h_i=2.0\ \mathrm{m}$};
\draw[<->, popmuted] (0, -1.5) -- (3.9, -1.5) node[midway, below] {$v$};
% Rays
\draw[poporange, thick, ->] (-2.5, 0.5) -- (0, 0.5) -- (3.9, -2.0);
\draw[poporange, thick, ->] (-2.5, 0.5) -- (0, 0) -- (3.9, -2.0);
\draw[poporange, thick, ->] (-2.5, 0) -- (0, 0) -- (3.9, 0);
% Total distance
\draw[<->, poppurple] (-2.5, 1.5) -- (3.9, 1.5) node[midway, above] {$u+v = 4.0\ \mathrm{m}$};
\end{tikzpicture}
\legende{Projector geometry: real, inverted, magnified image formed by a converging lens.}
\end{popfigure}

\textbf{Part D — Lighting the classroom.}\par
The classroom measures $6\ \mathrm{m} \times 8\ \mathrm{m}$, with desks $3.0\ \mathrm{m}$ below the ceiling. The school will install identical lamps fixed to the ceiling, each treated as a \cle{point source} radiating uniformly in all directions, of total luminous power $P = 40\ \mathrm{W}$. For comfortable reading, the intensity received at desk level directly under a lamp must be at least $I_{\min} = 0.30\ \mathrm{W\,m^{-2}}$. Lamps cost money to buy and to run, so the school wants the \emph{smallest} number of lamps that satisfies the requirement, spread evenly over the ceiling.

\begin{popfigure}
\begin{tikzpicture}[scale=0.4]
% Ceiling outline
\draw[popdark, thick] (0, 0) rectangle (8, 6);
% Lamps (4)
\foreach \x/\y in {1.5/2, 1.5/4, 6.5/2, 6.5/4} {
    \fill[popgold] (\x, \y) circle (0.2);
    \draw[popdark] (\x, \y) circle (0.2);
    \node[popdark, below] at (\x, \y-0.3) {Lamp};
}
% Desk level points (worst case: center of room)
\fill[popblue] (4, 3) circle (0.15);
\node[popdark, above] at (4, 3.3) {Worst-lit desk};
% Distances
\draw[dashed, popmuted] (1.5, 2) -- (4, 3);
\draw[dashed, popmuted] (1.5, 4) -- (4, 3);
\draw[dashed, popmuted] (6.5, 2) -- (4, 3);
\draw[dashed, popmuted] (6.5, 4) -- (4, 3);
\node[popmuted, rotate=30] at (2.5, 2.2) {$r_{\text{horiz}}=2.5\ \mathrm{m}$};
\node[popmuted, rotate=-30] at (2.5, 3.8) {$r_{\text{horiz}}=2.5\ \mathrm{m}$};
% Dimensions
\draw[<->, poppurple] (0, -0.5) -- (8, -0.5) node[midway, below] {$8\ \mathrm{m}$};
\draw[<->, poppurple] (-0.5, 0) -- (-0.5, 6) node[midway, left] {$6\ \mathrm{m}$};
\end{tikzpicture}
\legende{Proposed 4-lamp layout on $6\ \mathrm{m}\times 8\ \mathrm{m}$ ceiling; worst-lit point is the room centre.}
\end{popfigure}

\courssub{Tasks}

\begin{enumerate}
\item \textbf{Critical angle of the fibre.} State the condition for total internal reflection at the core--cladding boundary. Calculate the critical angle $\theta_{c}$ for light travelling in the core ($n_{1} = 1.48$) towards the cladding ($n_{2} = 1.46$). Explain, in two or three sentences, why light entering the fibre along its axis remains trapped even when the fibre is bent.

\item \textbf{Acceptance angle.} Show that a ray entering the flat end of the fibre from air ($n_{\text{air}} = 1.00$) is guided only if its angle of incidence $i$ on the end face satisfies $\sin i \leq \sqrt{n_{1}^{2}-n_{2}^{2}}$. Compute the maximum acceptance angle $i_{\max}$. Why does this result matter to the technician aligning the test laser?

\item \textbf{Is the laser monochromatic?} Calculate the grating spacing $d$ for $600$ lines per millimetre. Using the first-order maximum at $\theta_{1} = 23.1^{\circ}$, determine the laser wavelength $\lambda$ and state the colour of the light. Explain why the observation of a \emph{single sharp spot} at first order, with no spreading, confirms that the laser is monochromatic. What would be seen if the source contained two different wavelengths?

\item \textbf{Choosing the projector lens.} Using the magnification $m = \dfrac{v}{u}$ and the thin-lens relation $\dfrac{1}{f} = \dfrac{1}{u} + \dfrac{1}{v}$ (real-is-positive convention), with $u + v = 4.0\ \mathrm{m}$ and an image height of $2.0\ \mathrm{m}$ for a $5.0\ \mathrm{cm}$ slide, find the object distance $u$, the image distance $v$, the magnification $m$, and the required focal length $f$ of the converging lens. Express the power of the lens in dioptres.

\item \textbf{How many lamps?} Write down the inverse square law for the intensity at distance $r$ from a point source of power $P$, and state the assumption about the medium. Calculate the intensity at a desk $3.0\ \mathrm{m}$ directly below one $40\ \mathrm{W}$ lamp. Deduce the minimum number of lamps needed so that the combined intensity at a desk midway between neighbouring lamps is at least $0.30\ \mathrm{W\,m^{-2}}$. (Hint: for a desk directly under one lamp, contributions from the other lamps are smaller because they are further away; a simple acceptable estimate is to require $N \times I_{\text{one lamp}} \geq I_{\min}$ at the worst-lit desk.) Propose a layout for the $6\ \mathrm{m} \times 8\ \mathrm{m}$ ceiling.

\item \textbf{Report and presentation.} Prepare a five-minute group presentation: one slide per part (A--D), each showing the law used, the calculation, the result with unit, and one sentence of recommendation to the school.
\end{enumerate}

\courssub{Model answer}

\begin{exoresolu}[Task 1 — Critical angle of the fibre]
\textbf{Condition for total internal reflection (TIR):} Light must travel from a medium of higher refractive index to one of lower refractive index ($n_1 > n_2$), and the angle of incidence $\theta$ at the boundary must exceed the critical angle $\theta_c$ defined by Snell's law with a refracted angle of $90^\circ$:
\[
n_1 \sin\theta_c = n_2 \sin 90^\circ \quad\Longrightarrow\quad \sin\theta_c = \frac{n_2}{n_1}.
\]
\textbf{Calculation:}
\[
\sin\theta_c = \frac{1.46}{1.48} = 0.986486\ldots
\]
\[
\theta_c = \arcsin(0.986486\ldots) = 80.6^\circ \quad (\text{to 3 significant figures}).
\]
\[\boxed{\theta_c = 80.6^\circ}\]
\textbf{Why light stays trapped:} A ray propagating near the fibre axis strikes the core-cladding interface at a grazing angle (measured from the normal) greater than $80.6^\circ$. Since $n_1 > n_2$ and $\theta > \theta_c$, Snell's law cannot be satisfied for a transmitted ray; all the light is reflected back into the core. This process repeats at every bounce, guiding the light along the fibre even around gentle bends, provided the bend radius is large enough that the local incidence angle remains above $\theta_c$.
\end{exoresolu}

\begin{exoresolu}[Task 2 — Acceptance angle]
\textbf{Derivation:} Let a ray enter the flat end face from air ($n_0=1.00$) at incidence $i$, refracting to angle $r$ inside the core: $\sin i = n_1 \sin r$ (Snell's law). The ray then meets the core-cladding boundary at an angle of incidence $\theta = 90^\circ - r$. For guidance we need $\theta \ge \theta_c$, i.e. $\sin\theta \ge \sin\theta_c = n_2/n_1$. Since $\sin\theta = \cos r$, the limiting case gives
\[
\cos r = \frac{n_2}{n_1} \quad\Longrightarrow\quad \sin r = \sqrt{1 - \frac{n_2^2}{n_1^2}}.
\]
Substituting into Snell's law at the entrance face:
\[
\sin i_{\max} = n_1 \sin r = n_1 \sqrt{1 - \frac{n_2^2}{n_1^2}} = \sqrt{n_1^2 - n_2^2}.
\]
\textbf{Numerical value:}
\[
\sin i_{\max} = \sqrt{1.48^2 - 1.46^2} = \sqrt{2.1904 - 2.1316} = \sqrt{0.0588} = 0.2425\ldots
\]
\[
i_{\max} = \arcsin(0.2425\ldots) = 14.0^\circ \quad (\text{to 3 significant figures}).
\]
\[\boxed{i_{\max} = 14.0^\circ}\]
\textbf{Significance:} Any ray entering within a cone of half-angle $14.0^\circ$ about the fibre axis will be guided. The technician must therefore align the test laser carefully with the fibre axis; a beam entering at an angle larger than $14.0^\circ$ would strike the core-cladding boundary below the critical angle and leak into the cladding, giving a false impression of fibre quality.
\end{exoresolu}

\begin{exoresolu}[Task 3 — Testing the laser with the grating]
\textbf{Grating spacing:} $600$ lines per millimetre means
\[
d = \frac{1\ \mathrm{mm}}{600} = \frac{10^{-3}\ \mathrm{m}}{600} = 1.667\times 10^{-6}\ \mathrm{m}.
\]
\textbf{Wavelength from first-order maximum:} For normal incidence the grating equation is $d\sin\theta = n\lambda$. At first order ($n=1$) with $\theta_1 = 23.1^\circ$:
\[
\lambda = d\sin\theta_1 = (1.667\times 10^{-6}\ \mathrm{m}) \times \sin 23.1^\circ.
\]
$\sin 23.1^\circ = 0.3923\ldots$
\[
\lambda = 1.667\times 10^{-6} \times 0.3923 = 6.54\times 10^{-7}\ \mathrm{m} = 654\ \mathrm{nm}.
\]
\[\boxed{\lambda = 654\ \mathrm{nm}}\]
This wavelength corresponds to \textbf{red} light.
\textbf{Monochromaticity check:} The grating equation shows that the diffraction angle $\theta$ depends on wavelength ($\sin\theta \propto \lambda$). If the source contained two or more distinct wavelengths, each would produce its own first-order maximum at a different angle, resulting in several separated spots (or a continuous coloured band if the source had a broad spectrum). The observation of a \emph{single sharp spot} on each side of the central maximum, with no spreading or colour fringing, confirms that only one wavelength is present: the laser is monochromatic. This is essential for fibre testing because a non-monochromatic source would suffer material dispersion (different wavelengths travel at different speeds), broadening pulses and limiting bandwidth.
\end{exoresolu}

\begin{exoresolu}[Task 4 — Choosing the projector lens]
\textbf{Magnification:}
\[
m = \frac{\text{image height}}{\text{object height}} = \frac{2.0\ \mathrm{m}}{5.0\ \mathrm{cm}} = \frac{200\ \mathrm{cm}}{5.0\ \mathrm{cm}} = 40.
\]
\[\boxed{m = 40}\]
\textbf{Object and image distances:} For a real image formed by a single converging lens, $m = v/u$ (real-is-positive convention), so $v = 40\,u$. The slide-to-screen distance is fixed:
\[
u + v = 4.0\ \mathrm{m} \quad\Longrightarrow\quad u + 40u = 41u = 4.0\ \mathrm{m}.
\]
\[
u = \frac{4.0}{41} = 0.09756\ \mathrm{m} = 9.76\ \mathrm{cm}, \qquad
v = 40u = 3.902\ \mathrm{m}.
\]
\[\boxed{u = 9.76\ \mathrm{cm}}, \quad \boxed{v = 3.90\ \mathrm{m}}\]
\textbf{Focal length and power:} Thin-lens equation:
\[
\frac{1}{f} = \frac{1}{u} + \frac{1}{v} = \frac{1}{0.09756} + \frac{1}{3.902} = 10.250 + 0.256 = 10.506\ \mathrm{m^{-1}}.
\]
\[
f = \frac{1}{10.506} = 0.09518\ \mathrm{m} = 9.52\ \mathrm{cm}.
\]
\[\boxed{f = 9.52\ \mathrm{cm}}\]
Lens power:
\[
F = \frac{1}{f} = 10.5\ \mathrm{D} \quad (\text{dioptres}).
\]
\[\boxed{F = 10.5\ \mathrm{D}}\]
\textbf{Recommendation:} Purchase a converging lens of focal length approximately $9.5\ \mathrm{cm}$ (power $+10.5\ \mathrm{D}$). Mount the slide just outside the focal point, $9.8\ \mathrm{cm}$ from the lens, inverted so that the projected image on the screen is upright for the audience.
\end{exoresolu}

\begin{exoresolu}[Task 5 — How many lamps?]
\textbf{Inverse square law:} For a point source radiating total power $P$ uniformly in all directions in free space (no absorption, no reflections), energy conservation over a sphere of radius $r$ gives the intensity (power per unit area):
\[
I = \frac{P}{4\pi r^2}.
\]
\textbf{Intensity directly below one lamp:} Desk height $r = 3.0\ \mathrm{m}$, $P = 40\ \mathrm{W}$.
\[
I_1 = \frac{40}{4\pi (3.0)^2} = \frac{40}{36\pi} = \frac{10}{9\pi} = 0.3537\ \mathrm{W\,m^{-2}}.
\]
\[\boxed{I_{1} = 0.354\ \mathrm{W\,m^{-2}}}\]
One lamp alone exceeds the minimum $0.30\ \mathrm{W\,m^{-2}}$ at the point directly beneath it. However, the requirement must be met \emph{everywhere} on the desks, including the worst-lit points midway between lamps and in the corners.
\textbf{Estimating the number of lamps:} Consider a symmetric layout. With a single central lamp, the corners are at horizontal distance $\sqrt{3^2+4^2}=5.0\ \mathrm{m}$, slant range $r = \sqrt{3.0^2+5.0^2} = 5.83\ \mathrm{m}$, giving $I_{\text{corner}} = 40/(4\pi\times 5.83^2) = 0.094\ \mathrm{W\,m^{-2}} < 0.30\ \mathrm{W\,m^{-2}}$ — insufficient.
\textbf{Four-lamp layout:} Place one lamp at the centre of each quarter of the ceiling ($6\ \mathrm{m}\times 8\ \mathrm{m}$ $\rightarrow$ quarters of $3\ \mathrm{m}\times 4\ \mathrm{m}$). Lamp coordinates (from one corner): $(1.5, 2)$, $(1.5, 6)$, $(4.5, 2)$, $(4.5, 6)$ metres. The worst-lit point is the room centre $(3, 4)$. Horizontal distance from each lamp to centre:
\[
r_{\text{horiz}} = \sqrt{(3-1.5)^2 + (4-2)^2} = \sqrt{1.5^2 + 2^2} = 2.5\ \mathrm{m}.
\]
Slant range to desk level ($3.0\ \mathrm{m}$ below ceiling):
\[
r = \sqrt{3.0^2 + 2.5^2} = \sqrt{15.25} = 3.905\ \mathrm{m}.
\]
Intensity from one lamp at this point:
\[
I_{\text{each}} = \frac{40}{4\pi (3.905)^2} = \frac{40}{4\pi \times 15.25} = 0.208\ \mathrm{W\,m^{-2}}.
\]
Total intensity from four lamps (additive, since intensities are scalar):
\[
I_{\text{total}} = 4 \times 0.208 = 0.832\ \mathrm{W\,m^{-2}} \ge 0.30\ \mathrm{W\,m^{-2}}.
\]
\[\boxed{\text{Install 4 lamps of } 40\ \mathrm{W} \text{, evenly spaced in a } 2\times 2 \text{ grid on the ceiling.}}\]
This satisfies the requirement with a comfortable safety margin and uses the minimum number of lamps (3 lamps cannot cover the corners adequately).
\end{exoresolu}

\begin{exoresolu}[Task 6 — Report and presentation structure]
\textbf{Slide 1 (Part A — Fibre link):} Law: Snell's law and TIR condition $\sin\theta_c = n_2/n_1$. Calculation: $\theta_c = 80.6^\circ$. Recommendation: The fibre will guide light reliably around gentle bends; ensure bend radius $> 5\ \mathrm{cm}$ to keep incidence angles above $\theta_c$.
\textbf{Slide 2 (Part B — Laser test):} Law: Grating equation $d\sin\theta = n\lambda$. Calculation: $d=1.67\ \mu\mathrm{m}$, $\lambda=654\ \mathrm{nm}$ (red). Recommendation: Laser is monochromatic and suitable for fibre testing; label it ``654 nm'' for future use.
\textbf{Slide 3 (Part C — Projector):} Laws: $m=v/u$, $1/f=1/u+1/v$. Calculation: $u=9.8\ \mathrm{cm}$, $v=3.90\ \mathrm{m}$, $f=9.5\ \mathrm{cm}$ ($+10.5\ \mathrm{D}$). Recommendation: Buy a $+10.5\ \mathrm{D}$ converging lens; mount slide $9.8\ \mathrm{cm}$ from lens.
\textbf{Slide 4 (Part D — Lighting):} Law: Inverse square law $I=P/(4\pi r^2)$. Calculation: $I_1=0.354\ \mathrm{W\,m^{-2}}$ at $3.0\ \mathrm{m}$; 4 lamps give $0.83\ \mathrm{W\,m^{-2}}$ at worst point. Recommendation: Install four $40\ \mathrm{W}$ lamps in a $2\times 2$ grid on the ceiling.
\end{exoresolu}

\begin{savaistu}[Optical fibres in Cameroon]
Cameroon's national backbone network (CAMTEL) relies heavily on step-index and graded-index optical fibres to connect major cities like Douala, Yaoundé, Bafoussam and Bamenda. The principle of total internal reflection demonstrated in this activity allows data to travel hundreds of kilometres with minimal loss. Modern fibres have a numerical aperture $\mathrm{NA} = \sqrt{n_1^2-n_2^2}$ around $0.2$ (similar to our $0.24$), enabling efficient coupling with laser diodes. The $654\ \mathrm{nm}$ red laser you identified is visible; telecom fibres actually use infrared lasers at $1310\ \mathrm{nm}$ or $1550\ \mathrm{nm}$ where glass attenuation is lowest ($<0.2\ \mathrm{dB/km}$).
\end{savaistu}

\begin{aretenir}[What this activity demonstrates]
One authentic project mobilised the whole chapter: total internal reflection and the critical angle (fibre), Snell's law (acceptance cone), the grating equation $d\sin\theta = n\lambda$ (laser test), the thin-lens equation and magnification (projector), and the inverse square law $I = P/(4\pi r^2)$ (lighting). In the examination, always state the law, substitute with units, and conclude with a physically sensible recommendation.
\end{aretenir}

\newpage

\coursec{Methods and worked examples}

\coursec{Section 1: Measuring the Speed of Sound and the Nature of Electromagnetic Waves}

\begin{experience}[Measuring the Speed of Sound in Air -- Resonance Tube Method]
\textbf{Objective:} To determine the speed of sound in air at room temperature using a resonance tube and a tuning fork.

\textbf{Apparatus:} Tuning fork of known frequency $f$ (e.g., $512\ \mathrm{Hz}$), resonance tube (graduated glass tube $\approx 1\ \mathrm{m}$ long) connected to a water reservoir, thermometer, rubber pad, metre rule.

\textbf{Theory:} A stationary wave is formed in the air column by the superposition of the incident wave from the fork and the wave reflected at the water surface. At the water surface (closed end) there is a \textbf{displacement node} (pressure antinode). At the open end there is a \textbf{displacement antinode} (pressure node), but the antinode lies slightly outside the tube -- this is the \textbf{end correction} $e$.

For the first resonance (fundamental): $L_1 + e = \frac{\lambda}{4}$.
For the second resonance (first overtone): $L_2 + e = \frac{3\lambda}{4}$.
Subtracting eliminates $e$: $L_2 - L_1 = \frac{\lambda}{2} \Rightarrow \lambda = 2(L_2 - L_1)$.
Speed of sound: $v = f\lambda = 2f(L_2 - L_1)$.

\textbf{Procedure:}
\begin{enumerate}
\item Fill the reservoir with water and raise it so the water level is near the top of the tube.
\item Strike the tuning fork gently on the rubber pad (do not hit hard -- harmonics distort the note). Hold it horizontally just above the open end of the tube.
\item Lower the reservoir slowly (or raise the tube) while listening for a sudden increase in loudness -- the \textbf{first resonance}. Record the length $L_1$ of the air column (water level to top of tube).
\item Continue lowering the reservoir until the next loudness maximum -- the \textbf{second resonance}. Record $L_2$.
\item Repeat with a different tuning fork if available. Measure room temperature $T$.
\end{enumerate}

\textbf{Results Table:}
\begin{center}
\begin{tabular}{|c|c|c|c|c|c|}\hline
Fork $f$ (\si{Hz}) & $L_1$ (\si{cm}) & $L_2$ (\si{cm}) & $\Delta L = L_2-L_1$ (\si{cm}) & $\lambda = 2\Delta L$ (\si{m}) & $v = f\lambda$ (\si{m.s^{-1}}) \\ \hline
512 & 15.2 & 48.5 & 33.3 & 0.666 & 341 \\ \hline
\end{tabular}
\end{center}

\textbf{Analysis:} Calculate $v$ for each fork. Average the values. Compare with theoretical $v_T = 331\sqrt{1 + T/273}\ \mathrm{m\,s^{-1}}$. Determine end correction $e = \lambda/4 - L_1$.

\textbf{Precautions / Sources of Error:}
\begin{itemize}
\item Fork must be struck gently on rubber (not bench) to avoid overtones.
\item Tube must be vertical; water surface horizontal.
\item Temperature of air column may differ from room thermometer (evaporative cooling at water surface).
\item End correction assumes flanged open end; unflanged correction is $\approx 0.6r$ ($r$ = tube radius).
\end{itemize}
\end{experience}

\begin{definition}[Mechanical vs Electromagnetic Waves]
A \textbf{mechanical wave} (e.g., sound) requires a material medium for propagation; energy is transferred by oscillations of particles. An \textbf{electromagnetic (EM) wave} consists of oscillating electric ($\vec{E}$) and magnetic ($\vec{B}$) fields perpendicular to each other and to the direction of propagation; it requires \textbf{no medium} and travels through vacuum at speed $c = 3.00 \times 10^8\ \mathrm{m\,s^{-1}}$.
\end{definition}

\begin{propriete}[Properties of Electromagnetic Waves in Vacuum]
An electromagnetic wave propagating in vacuum has the following properties (proof not examinable):
\begin{itemize}
\item It is a \cle{transverse} wave: the electric field $\vec{E}$ and the magnetic field $\vec{B}$ are both perpendicular to the direction of propagation.
\item The vectors $\vec{E}$, $\vec{B}$ and the propagation vector form a \cle{right-handed triad}: $\vec{E}$, $\vec{B}$ and the direction of propagation are mutually perpendicular.
\item The fields $\vec{E}$ and $\vec{B}$ oscillate \cle{in phase} with each other.
\item The magnitudes of the fields are related by
\[ \frac{E}{B} = c, \]
where $c = 3.00\times 10^{8}\ \mathrm{m\,s^{-1}}$ is the speed of light in vacuum.
\item All electromagnetic waves travel in vacuum with the same speed $c$, whatever their frequency or wavelength, with
\[ c = \lambda\, f, \]
where $\lambda$ is the wavelength (in $\mathrm{m}$) and $f$ the frequency (in $\mathrm{Hz}$).
\item Electromagnetic waves transport \cle{energy} and can propagate through vacuum, unlike mechanical waves which require a material medium.
\end{itemize}
\end{propriete}

\begin{aretenir}[Key Points -- Section 1]
\begin{itemize}
\item Speed of sound in air measured by resonance tube: $v = 2f(L_2 - L_1)$ (end correction eliminated).
\item EM waves are transverse, travel in vacuum at $c$, obey $c = f\lambda$.
\item $\vec{E}$, $\vec{B}$, $\vec{k}$ mutually perpendicular; $E_0 = cB_0$.
\item Intensity $I \propto E_0^2$; radiation pressure exists.
\end{itemize}
\end{aretenir}

\coursec{Section 2: The Electromagnetic Spectrum, X-rays and Polarisation}

\begin{definition}[The Electromagnetic Spectrum]
The complete range of EM waves classified by frequency $f$ or wavelength $\lambda$ in vacuum. All travel at $c$. Photon energy $E = hf = hc/\lambda$.
\end{definition}

\begin{center}
\begin{tabularx}{\linewidth}{|l|c|c|c|X|}\hline
\rowcolor{popblueL}
\textbf{Region} & \textbf{$\lambda$ range (vacuum)} & \textbf{$f$ range} & \textbf{Photon Energy} & \textbf{Production / Detection / Uses} \\ \hline
Radio & $> 1\ \mathrm{m}$ & $< 300\ \mathrm{MHz}$ & $< 1.2\ \mu\mathrm{eV}$ & AC circuits, aerials / aerials, resonant circuits / Broadcasting, communication, MRI \\ \hline
Microwave & $1\ \mathrm{mm} - 1\ \mathrm{m}$ & $300\ \mathrm{MHz} - 300\ \mathrm{GHz}$ & $1.2\ \mu\mathrm{eV} - 1.2\ \mathrm{meV}$ & Klystron, magnetron, Gunn diode / Horn antenna, bolometer / Radar, satellite TV, microwave ovens, mobile phones \\ \hline
Infrared (IR) & $700\ \mathrm{nm} - 1\ \mathrm{mm}$ & $300\ \mathrm{GHz} - 430\ \mathrm{THz}$ & $1.2\ \mathrm{meV} - 1.7\ \mathrm{eV}$ & Hot bodies, LEDs / Thermopile, photodiode, photographic plate / Thermal imaging, remote controls, fibre optics, night vision \\ \hline
Visible & $400 - 700\ \mathrm{nm}$ & $430 - 750\ \mathrm{THz}$ & $1.7 - 3.1\ \mathrm{eV}$ & Hot bodies, lasers, LEDs / Eye, photographic film, CCD / Vision, photography, optical fibres, spectroscopy \\ \hline
Ultraviolet (UV) & $10 - 400\ \mathrm{nm}$ & $750\ \mathrm{THz} - 30\ \mathrm{PHz}$ & $3.1\ \mathrm{eV} - 124\ \mathrm{eV}$ & Mercury vapour lamp, synchrotron / Photographic plate, fluorescent screen, photodiode / Sterilisation, fluorescence, vitamin D synthesis, forgery detection \\ \hline
X-rays & $0.01 - 10\ \mathrm{nm}$ & $30\ \mathrm{PHz} - 30\ \mathrm{EHz}$ & $124\ \mathrm{eV} - 124\ \mathrm{keV}$ & X-ray tube (bremsstrahlung + characteristic), synchrotron / Photographic film, Geiger counter, scintillator / Medical imaging, crystallography, security scanning, radiotherapy \\ \hline
$\gamma$-rays & $< 0.01\ \mathrm{nm}$ & $> 30\ \mathrm{EHz}$ & $> 124\ \mathrm{keV}$ & Radioactive nuclei, nuclear reactions / Geiger counter, scintillator, semiconductor detector / Radiotherapy, sterilisation, nuclear medicine, astronomy \\ \hline
\end{tabularx}
\end{center}

\begin{savaistu}[Cameroon Context]
\begin{itemize}
\item \textbf{Radio:} CRTV broadcasts on FM ($88-108\ \mathrm{MHz}$, $\lambda \approx 3\ \mathrm{m}$) and AM ($531-1602\ \mathrm{kHz}$, $\lambda \approx 200-500\ \mathrm{m}$). Community radios use low-power FM transmitters.
\item \textbf{Microwave:} MTN, Orange, Camtel microwave links (backbone) and 3G/4G/5G base stations use cm-wave bands. VSAT dishes (Ku-band $\approx 12-18\ \mathrm{GHz}$) for internet in remote areas.
\item \textbf{X-rays:} Hospital radiography departments (Yaoundé, Douala, regional hospitals) use Coolidge tubes ($30-150\ \mathrm{kV}$ for diagnostics). Airport baggage scanners use dual-energy X-rays.
\item \textbf{Optical Fibre:} Camtel's national backbone uses IR lasers ($1550\ \mathrm{nm}$) in single-mode fibre for low-loss long-distance transmission.
\end{itemize}
\end{savaistu}

\courssub{X-ray Production and Properties}

\begin{experience}[The Coolidge X-ray Tube]
\textbf{Principle:} Thermionic emission from a heated cathode (filament), acceleration by high voltage $V$, sudden deceleration at a heavy metal target (anode) $\rightarrow$ X-rays.

\textbf{Construction:}
\begin{itemize}
\item \textbf{Cathode:} Tungsten filament (heated by low voltage $V_f \approx 10\ \mathrm{V}$, current $I_f \approx 3-5\ \mathrm{A}$) surrounded by a focusing cup (negative bias) to focus electrons onto a small spot (focal spot) on the anode.
\item \textbf{Anode (Target):} Tungsten (high $Z=74$, high melting point $3422^\circ\mathrm{C}$) embedded in copper stem for heat conduction. \textbf{Rotating anode} (up to $10\,000\ \mathrm{rpm}$) spreads heat over a track $\rightarrow$ higher power loading.
\item \textbf{Envelope:} Pyrex glass or metal/ceramic, high vacuum ($\approx 10^{-6}\ \mathrm{Pa}$).
\item \textbf{High Voltage:} $20-150\ \mathrm{kV}$ (diagnostic), up to $300\ \mathrm{kV}$ (therapy). Rectified and smoothed (constant potential) for stable spectrum.
\end{itemize}

\textbf{Two Mechanisms of X-ray Production:}
\begin{enumerate}
\item \textbf{Bremsstrahlung (Continuous Spectrum):} Incident electron decelerated in Coulomb field of target nucleus. Photon energy $hf = eV - \text{residual KE}$. Maximum photon energy (minimum wavelength) when electron loses all KE in one collision:
\[
eV = \frac{hc}{\lambda_{\min}} \quad \Rightarrow \quad \lambda_{\min} = \frac{hc}{eV} \approx \frac{1240\ \mathrm{eV\,nm}}{V(\mathrm{V})} \quad \text{(Duane-Hunt Law)}.
\]
Intensity $\propto Z (V - V_{\min})^2$ approximately. Spectrum is continuous up to $\lambda_{\min}$.
\item \textbf{Characteristic Lines (Line Spectrum):} Incident electron knocks out inner-shell electron (K, L, M...). Outer electron fills vacancy $\rightarrow$ photon of precise energy $E_K - E_L$ etc. (e.g., $K_\alpha$, $K_\beta$). Threshold voltage $V_K$ needed to eject K-shell electron. Independent of $V$ (once $V > V_K$). Superimposed on bremsstrahlung background.
\end{enumerate}

\textbf{Quality and Quantity Control:}
\begin{itemize}
\item \textbf{Tube Voltage $V$ (kVp):} Controls \textbf{penetrating power} (quality / hardness). Higher $V \rightarrow$ shorter $\lambda_{\min}$, more energetic photons, harder beam.
\item \textbf{Tube Current $I$ (mA):} Controls \textbf{intensity} (quantity). Number of photons $\propto I$. Exposure $\propto I \times t$ (mAs).
\item \textbf{Filtration:} Aluminium/Copper filters absorb low-energy (soft) photons that contribute to patient dose but not image $\rightarrow$ hardens beam.
\end{itemize}
\end{experience}

\begin{propriete}[X-ray Attenuation and Imaging]
Intensity after thickness $x$ of material: $I = I_0 e^{-\mu x}$ where $\mu$ = linear attenuation coefficient ($\mathrm{m^{-1}}$). Depends on $\lambda$ (or $E$) and $Z$ ($\mu \propto Z^3/E^3$ photoelectric effect dominant at diagnostic energies).
\textbf{Contrast:} Bone (high $Z$, high $\mu$) absorbs more $\rightarrow$ lighter on film (negative) / darker on digital detector. Soft tissue (low $Z$) transmits more.
\textbf{Dose:} Absorbed dose $D = \frac{\text{Energy deposited}}{\text{Mass}}$ (\si{Gy} = \si{J.kg^{-1}}). Equivalent dose $H = D \times w_R$ (\si{Sv}). For X-rays $w_R = 1$. \textbf{ALARA principle:} As Low As Reasonably Achievable.
\end{propriete}

\courssub{Plane Polarised Electromagnetic Waves}

\begin{definition}[Polarisation]
Unpolarised light: $\vec{E}$ oscillates in all directions perpendicular to $\vec{k}$, randomly changing on timescale $\sim 10^{-8}\ \mathrm{s}$.
\textbf{Plane (linearly) polarised light:} $\vec{E}$ oscillates in a \textbf{fixed plane} containing the propagation direction. The plane of polarisation is defined by the direction of $\vec{E}$.
\end{definition}

\begin{propriete}[Malus's Law]
If plane polarised light of intensity $I_0$ is incident on an analyser (Polaroid) whose transmission axis makes an angle $\theta$ with the polarisation direction of the incident light, the transmitted intensity is:
\[
I = I_0 \cos^2\theta.
\]
If unpolarised light of intensity $I_u$ is incident on a polariser, transmitted intensity is $I_0 = I_u/2$ (ideal polariser).
\end{propriete}

\begin{propriete}[Polarisation by Reflection -- Brewster's Law]
When unpolarised light is incident at the \textbf{Brewster angle} $i_p$ on a dielectric interface (air/glass, water/air), the reflected ray is \textbf{perfectly plane polarised} with $\vec{E}$ perpendicular to the plane of incidence (parallel to the interface). The reflected and refracted rays are perpendicular ($i_p + r = 90^\circ$).
\[
\tan i_p = n_2/n_1 \quad (\text{relative refractive index}).
\]
\end{propriete}

\begin{attention}[Common Mistakes -- Polarisation]
\begin{itemize}
\item \textbf{Malus's Law angle:} $\theta$ is the angle between the \textbf{polarisation direction of the incident light} and the \textbf{transmission axis of the analyser}. Not the angle between polariser and analyser unless the first polariser defines the polarisation direction.
\item \textbf{Brewster's angle:} Only the reflected ray is perfectly polarised (s-polarised). The refracted ray is partially polarised (p-component stronger).
\item \textbf{Intensity at Brewster:} The reflected intensity is \textbf{not} $I_0/2$. It is given by the Fresnel reflection coefficient $R_s$ for s-polarisation. At GCE level, if not given, assume the problem provides the reflected intensity or asks only for the polarisation state.
\item \textbf{Sound waves:} Cannot be polarised (longitudinal in fluids). Only transverse waves can be polarised.
\end{itemize}
\end{attention}

\begin{methode}[EM Wave Calculations and Spectrum Identification]
\begin{enumerate}
\item \textbf{Identify knowns:} $f$, $\lambda$, $c = 3.00 \times 10^8\ \mathrm{m\,s^{-1}}$, $h = 6.63 \times 10^{-34}\ \mathrm{J\,s}$, $e = 1.60 \times 10^{-19}\ \mathrm{C}$.
\item \textbf{Wave equation:} $c = f\lambda$. Convert $\lambda$ to metres, $f$ to Hz.
\item \textbf{Photon energy:} $E(\mathrm{J}) = hf = hc/\lambda$. $E(\mathrm{eV}) = \frac{1240}{\lambda(\mathrm{nm})}$ (using $hc \approx 1240\ \mathrm{eV\,nm}$).
\item \textbf{Spectrum location:} Compare $\lambda$ or $f$ with standard bands (see table above).
\item \textbf{Medium change:} Frequency $f$ constant. Speed $v = c/n$, wavelength $\lambda' = \lambda/n$.
\item \textbf{X-ray tube voltage:} Duane-Hunt law $\lambda_{\min} = hc/(eV) \Rightarrow V = 1240/\lambda_{\min}(\mathrm{nm})$ volts.
\end{enumerate}
\end{methode}

\begin{exoresolu}[EM Wave Properties and X-ray Threshold]
An electromagnetic wave has a wavelength of $2.5\ \mathrm{nm}$ in vacuum.
\begin{enumerate}
\item Calculate its frequency and photon energy in joules and electronvolts.
\item Identify the region of the electromagnetic spectrum to which it belongs.
\item This wave is incident on a metal target to produce X-rays via bremsstrahlung. What is the minimum accelerating voltage required across an X-ray tube to produce photons of this wavelength?
\end{enumerate}
\tcblower
\textbf{Detailed Solution}

\textbf{1. Frequency and Photon Energy}
$\lambda = 2.5\ \mathrm{nm} = 2.5 \times 10^{-9}\ \mathrm{m}$.
$f = \frac{c}{\lambda} = \frac{3.00 \times 10^8}{2.5 \times 10^{-9}} = 1.20 \times 10^{17}\ \mathrm{Hz}$.
$E = hf = (6.63 \times 10^{-34}) \times (1.20 \times 10^{17}) = 7.96 \times 10^{-17}\ \mathrm{J}$.
$E(\mathrm{eV}) = \frac{7.96 \times 10^{-17}}{1.60 \times 10^{-19}} = 497\ \mathrm{eV}$.
(Shortcut: $E(\mathrm{eV}) = \frac{1240}{2.5} = 496\ \mathrm{eV}$).

\textbf{2. Spectrum Region}
$2.5\ \mathrm{nm}$ lies between $0.01\ \mathrm{nm}$ and $10\ \mathrm{nm}$ $\rightarrow$ \textbf{X-ray region} (soft X-rays).

\textbf{3. Minimum Accelerating Voltage (Duane-Hunt Law)}
To produce a photon of wavelength $\lambda = 2.5\ \mathrm{nm}$, electron KE must be at least $hc/\lambda$.
$V_{\min} = \frac{hc}{e\lambda} = \frac{1240\ \mathrm{eV\,nm}}{2.5\ \mathrm{nm}} = 496\ \mathrm{V}$.
\textbf{Answer:} \SI{496}{V}.
\end{exoresolu}

\begin{methode}[Polarisation: Malus's Law and Brewster's Angle]
\begin{enumerate}
\item \textbf{Setup:} Unpolarised light ($I_u$) $\xrightarrow{\text{Polariser}}$ Plane polarised ($I_0 = I_u/2$) $\xrightarrow{\text{Analyser at }\theta}$ Transmitted $I = I_0\cos^2\theta$.
\item \textbf{Crossed Polaroids ($\theta=90^\circ$):} $I=0$. Insert third at $45^\circ$: $I = \frac{I_u}{2}\cos^2 45^\circ \cos^2 45^\circ = I_u/8$.
\item \textbf{Brewster's Law:} $\tan i_p = n$. At $i_p$, reflected ray $\perp$ refracted ray. Reflected ray perfectly plane polarised ($\vec{E} \parallel$ interface).
\item \textbf{Applications:} Polaroid sunglasses (block horizontal glare from road/water), LCD screens (twisted nematic), stress analysis (photoelasticity), 3D cinema (circular polarisation).
\end{enumerate}
\end{methode}

\begin{exoresolu}[Polarisation by Reflection and Malus's Law]
Unpolarised light of intensity $I_u = 100\ \mathrm{W\,m^{-2}}$ is incident on a glass block of refractive index $n = 1.52$ at the Brewster angle.
\begin{enumerate}
\item Calculate the Brewster angle $i_p$.
\item The reflected beam is perfectly plane polarised. If the reflectance for s-polarisation at this angle is $R_s = 0.15$, find the intensity of the reflected beam.
\item This reflected beam passes through a Polaroid analyser whose transmission axis makes an angle of $30^\circ$ with the plane of polarisation of the reflected beam. Calculate the transmitted intensity.
\end{enumerate}
\tcblower
\textbf{Detailed Solution}

\textbf{1. Brewster Angle}
$\tan i_p = n = 1.52 \Rightarrow i_p = \arctan(1.52) = \SI{56.7}{\ensuremath{{}^{\circ}}}$.

\textbf{2. Reflected Intensity}
Incident unpolarised light has equal s and p components: $I_s = I_p = I_u/2 = 50\ \mathrm{W\,m^{-2}}$.
At $i_p$, $R_p = 0$, $R_s = 0.15$ (given).
Reflected intensity $I_r = R_s \times I_s = 0.15 \times 50 = \SI{7.5}{W.m^{-2}}$.
(Note: The reflected beam is 100\% plane polarised with $\vec{E}$ perpendicular to plane of incidence.)

\textbf{3. Transmitted Intensity (Malus's Law)}
Analyser angle $\theta = 30^\circ$ relative to polarisation direction of reflected beam.
$I_t = I_r \cos^2\theta = 7.5 \times \cos^2 30^\circ = 7.5 \times (\frac{\sqrt{3}}{2})^2 = 7.5 \times 0.75 = \SI{5.625}{W.m^{-2}}$.
\end{exoresolu}

\begin{exercice}[Polarisation Practice]
\begin{enumerate}
\item Unpolarised light of intensity $I_0$ passes through two Polaroids. The transmission axis of the second makes an angle of $60^\circ$ with the first. Find the transmitted intensity.
\item Light is incident on a water surface ($n=4/3$) at Brewster's angle. Calculate the angle of refraction. Verify that reflected and refracted rays are perpendicular.
\item A polariser and analyser are crossed. A third Polaroid is inserted between them at $30^\circ$ to the polariser. What fraction of the incident unpolarised intensity is transmitted?
\end{enumerate}
\end{exercice}

\begin{corrige}[Polarisation Practice]
\begin{enumerate}
\item $I = \frac{I_0}{2}\cos^2 60^\circ = \frac{I_0}{2} \times \frac{1}{4} = I_0/8$.
\item $\tan i_p = 4/3 \Rightarrow i_p \approx 53.1^\circ$. $r = 90^\circ - i_p = 36.9^\circ$. Check Snell: $\sin i_p = 0.8$, $n\sin r = 1.333 \times 0.6 = 0.8$. Perpendicular: $i_p + r = 90^\circ$. Verified.
\item $I = \frac{I_0}{2}\cos^2 30^\circ \cos^2 60^\circ = \frac{I_0}{2} \times \frac{3}{4} \times \frac{1}{4} = \frac{3I_0}{32}$.
\end{enumerate}
\end{corrige}

\coursec{Section 3: Light Sources, Diffraction Gratings and Multiple-Slit Diffraction}

\begin{definition}[Light Sources]
\begin{itemize}
\item \textbf{Thermal sources:} Incandescent filament (black-body spectrum, continuous, low efficiency), Gas discharge tubes (line spectra -- Na, Hg, Ne).
\item \textbf{Luminescence:} Fluorescence/Phosphorescence (UV $\rightarrow$ visible).
\item \textbf{Laser (Light Amplification by Stimulated Emission of Radiation):} Population inversion, optical cavity. Properties: \textbf{Monochromatic} (narrow $\Delta\lambda$), \textbf{Coherent} (constant phase relationship), \textbf{Collimated} (low divergence), \textbf{High intensity}. He-Ne ($632.8\ \mathrm{nm}$), Semiconductor laser ($650-1550\ \mathrm{nm}$), Nd:YAG ($1064\ \mathrm{nm}$).
\end{itemize}
\end{definition}

\begin{propriete}[Transmission Diffraction Grating -- Normal Incidence]
A grating consists of $N$ equally spaced parallel slits per metre. Grating spacing $d = 1/N$.
\textbf{Grating Equation (Principal Maxima):}
\[
d\sin\theta = n\lambda \quad (n = 0, \pm 1, \pm 2, \dots)
\]
where $\theta$ is the angle of diffraction (deviation from normal).
\textbf{Maximum Order:} $n_{\max} = \lfloor d/\lambda \rfloor = \lfloor 1/(N\lambda) \rfloor$ (since $\sin\theta \le 1$).
\textbf{Angular Dispersion:} $\frac{d\theta}{d\lambda} = \frac{n}{d\cos\theta}$. Higher order $n$, smaller $d$ $\rightarrow$ greater dispersion.
\textbf{Intensity Envelope:} Modulated by single-slit diffraction pattern ($a$ = slit width). Minima at $a\sin\theta = m\lambda$ ($m \ne 0$).
\textbf{Missing Orders:} Occur when $d\sin\theta = n\lambda$ coincides with $a\sin\theta = m\lambda \Rightarrow n/m = d/a$.
\textbf{Resolving Power (Rayleigh Criterion):} $R = \frac{\lambda}{\Delta\lambda} = n N_{\text{total}}$ where $N_{\text{total}}$ = total illuminated lines.
\end{propriete}

\begin{experience}[Spectrometer and Diffraction Grating -- Normal Incidence]
\textbf{Apparatus:} Spectrometer (collimator, telescope, rotating table), diffraction grating ($N$ lines/mm), mercury or sodium lamp.

\textbf{Adjustments (Critical for GCE Practical):}
\begin{enumerate}
\item \textbf{Eyepiece:} Focus on cross-wires (relaxed eye).
\item \textbf{Collimator:} Focus on slit $\rightarrow$ parallel beam (telescope focused at infinity, slit at collimator focus).
\item \textbf{Grating Normal Incidence:} (i) Rotate table until slit image reflected from grating face coincides with cross-wires $\rightarrow$ grating perpendicular to collimator axis. (ii) Rotate table by exactly $90^\circ$ $\rightarrow$ grating normal to collimator beam (normal incidence). Clamp table.
\item \textbf{Observation:} Rotate telescope to left/right to find spectra orders $n=1,2,\dots$. Measure $\theta_n$ for each spectral line.
\end{enumerate}

\textbf{Calculations:}
$\lambda = \frac{d\sin\theta_n}{n}$. Average values for different $n$.
Resolving power check: Can the grating resolve Na D-lines ($589.0, 589.6\ \mathrm{nm}$) in order $n$? Need $N_{\text{total}} > \lambda/(n\Delta\lambda)$.
\end{experience}

\begin{methode}[Diffraction Grating: Normal Incidence and Resolving Power]
\begin{enumerate}
\item \textbf{Grating spacing:} $d = 1/N$ ($N$ in lines per metre).
\item \textbf{Angles:} $\sin\theta_n = n\lambda/d$. Calculate $\theta_n$ for each order $n$ up to $n_{\max}$.
\item \textbf{Angular separation of two close $\lambda$'s:} Compute $\theta$ for each $\lambda$ exactly using $\sin\theta = n\lambda/d$, then $\Delta\theta = |\theta_2 - \theta_1|$. (Small angle approx: $\Delta\theta \approx \frac{n\Delta\lambda}{d\cos\theta}$).
\item \textbf{Resolving Power:} Required $R = \lambda/\Delta\lambda$. Minimum illuminated lines $N_{\text{total}} = R/n$. Beam width $w = N_{\text{total}} \times d$.
\item \textbf{Missing Orders:} Check if $d/a$ is integer ratio. If $d/a = 2$, orders $n=2,4,6\dots$ missing.
\end{enumerate}
\end{methode}

\begin{exoresolu}[Diffraction Grating Analysis]
A diffraction grating has $3000$ lines per centimetre. It is illuminated normally by a sodium lamp emitting two close wavelengths $\lambda_1 = 589.0\ \mathrm{nm}$ and $\lambda_2 = 589.6\ \mathrm{nm}$.
\begin{enumerate}
\item Calculate the grating spacing $d$.
\item Determine the highest order spectrum visible.
\item Calculate the angular separation between the two sodium lines in the second order spectrum.
\item What is the minimum number of lines on the grating that must be illuminated to resolve these two lines in the second order?
\end{enumerate}
\tcblower
\textbf{Detailed Solution}

\textbf{1. Grating Spacing $d$}
$N = 3000\ \mathrm{lines/cm} = 3.00 \times 10^5\ \mathrm{lines/m}$.
$d = \frac{1}{N} = \frac{1}{3.00 \times 10^5} = 3.333 \times 10^{-6}\ \mathrm{m} = \SI{3.33}{\mu m}$.

\textbf{2. Highest Order $n_{\max}$}
Condition: $n\lambda \le d$. Use longer $\lambda_2 = 589.6\ \mathrm{nm} = 5.896 \times 10^{-7}\ \mathrm{m}$.
$n_{\max} = \left\lfloor \frac{3.333 \times 10^{-6}}{5.896 \times 10^{-7}} \right\rfloor = \lfloor 5.65 \rfloor = 5$.
\textbf{Highest order visible: 5th order.}

\textbf{3. Angular Separation in 2nd Order ($n=2$)}
$\sin\theta_1 = \frac{2 \times 589.0 \times 10^{-9}}{3.333 \times 10^{-6}} = 0.3534 \Rightarrow \theta_1 = 20.68^\circ$.
$\sin\theta_2 = \frac{2 \times 589.6 \times 10^{-9}}{3.333 \times 10^{-6}} = 0.3538 \Rightarrow \theta_2 = 20.71^\circ$.
$\Delta\theta = \theta_2 - \theta_1 = \SI{0.03}{\ensuremath{{}^{\circ}}}$ (or $5.2 \times 10^{-4}\ \mathrm{rad}$).

\textbf{4. Minimum Illuminated Lines for Resolution}
Mean $\lambda = 589.3\ \mathrm{nm}$, $\Delta\lambda = 0.6\ \mathrm{nm}$.
$R = \frac{\lambda}{\Delta\lambda} = \frac{589.3}{0.6} = 982$.
$R = n N_{\text{total}} \Rightarrow N_{\text{total}} = \frac{R}{n} = \frac{982}{2} = 491$.
\textbf{Minimum lines: 491.} Beam width $= 491 \times 3.333\ \mu\mathrm{m} \approx \SI{1.64}{mm}$.
\end{exoresolu}

\begin{attention}[Grating Common Mistakes]
\begin{itemize}
\item \textbf{Units:} $N$ often given in lines/mm or lines/cm. Convert to lines/m for $d$ in metres.
\item \textbf{Order $n=0$:} Central white maximum (undeviated). Not used for $\lambda$ measurement.
\item \textbf{$n_{\max}$:} Always use the \textbf{longest} wavelength present to find the absolute maximum order visible.
\item \textbf{Resolving Power:} $N_{\text{total}}$ is the number of lines \textbf{illuminated}, not total lines on the grating. If beam width $w$ is given, $N_{\text{total}} = w/d$.
\item \textbf{Missing Orders:} Don't confuse single-slit minima ($a\sin\theta=m\lambda$) with grating maxima ($d\sin\theta=n\lambda$).
\end{itemize}
\end{attention}

\coursec{Section 4: Reflection, Refraction and Refractive Index}

\begin{propriete}[Laws of Reflection and Refraction]
\textbf{Reflection (Plane Interface):}
\begin{enumerate}
\item Incident ray, reflected ray, and normal lie in the same plane.
\item Angle of incidence $i$ = Angle of reflection $r$ ($i = r$).
\end{enumerate}
\textbf{Refraction (Snell's Law):}
\begin{enumerate}
\item Incident ray, refracted ray, and normal lie in the same plane.
\item $n_1 \sin i = n_2 \sin r$ where $n_1, n_2$ are absolute refractive indices ($n = c/v = \lambda_0/\lambda$).
\end{enumerate}
\textbf{Refractive Index:} $n = \frac{c}{v} = \frac{\lambda_0}{\lambda} = \frac{\sin i}{\sin r}$ (relative to vacuum/air). $n \ge 1$.
\end{propriete}

\begin{propriete}[Critical Angle and Total Internal Reflection (TIR)]
Light travels from denser medium ($n_1$) to rarer medium ($n_2 < n_1$).
\textbf{Critical Angle $C$:} Angle of incidence in denser medium for which $r = 90^\circ$.
\[
\sin C = \frac{n_2}{n_1}.
\]
\textbf{Total Internal Reflection:} If $i > C$, no refracted ray; all energy reflected ($R=1$). Phase change occurs (depends on polarisation).
\textbf{Conditions for TIR:} (1) Light travels from optically denser to rarer medium. (2) Angle of incidence $i > C$.
\end{propriete}

\begin{savaistu}[Applications of TIR in Cameroon]
\begin{itemize}
\item \textbf{Optical Fibres:} Camtel's national backbone and FTTH (Fibre To The Home) in Yaoundé/Douala use silica glass fibres ($n_{\text{core}} \approx 1.48$, $n_{\text{cladding}} \approx 1.46$). Critical angle $\approx 80^\circ$. Light guided over km with minimal loss ($< 0.2\ \mathrm{dB/km}$ at $1550\ \mathrm{nm}$).
\item \textbf{Prism Binoculars/Periscopes:} Right-angled prisms ($45^\circ-45^\circ-90^\circ$) use TIR at $45^\circ > C$ (for glass $C \approx 42^\circ$) to erect images without mirrors.
\item \textbf{Mirage:} Hot sand heats air $\rightarrow$ $n$ decreases with height $\rightarrow$ light from sky curves upward (TIR in gradient) $\rightarrow$ appears as water on road. Common in Far North region dry season.
\end{itemize}
\end{savaistu}

\begin{propriete}[Dispersion and Prisms]
\textbf{Dispersion:} $n$ varies with $\lambda$ ($n$ decreases as $\lambda$ increases -- normal dispersion). $n_{\text{blue}} > n_{\text{red}} \Rightarrow$ blue deviated more.
\textbf{Prism Geometry:} Refracting angle $A$. Deviation $D = i + e - A$. $r_1 + r_2 = A$.
\textbf{Minimum Deviation $D_{\min}$:} Symmetric passage ($i = e$, $r_1 = r_2 = A/2$).
\[
n = \frac{\sin\frac{A + D_{\min}}{2}}{\sin\frac{A}{2}}.
\]
\textbf{Spectrometer use:} Measure $D_{\min}$ for each spectral line $\rightarrow$ determine $n(\lambda)$ $\rightarrow$ Cauchy's formula $n = A + B/\lambda^2$.
\end{propriete}

\begin{methode}[Refraction, Critical Angle and Prism Problems]
\begin{enumerate}
\item \textbf{Snell's Law:} $n_1\sin i = n_2\sin r$. Identify media, draw normal.
\item \textbf{Critical Angle:} $\sin C = n_2/n_1$ (denser $\to$ rarer). Check if $i > C$ for TIR.
\item \textbf{Prism:} Draw ray diagram. Use $r_1 + r_2 = A$. Apply Snell at each face.
\item \textbf{Minimum Deviation:} Use $n = \sin\frac{A+D_{\min}}{2} / \sin\frac{A}{2}$. If $i$ given, check if $i = e$ (then $D = D_{\min}$).
\item \textbf{Apparent Depth:} Near normal viewing: $d_{\text{app}} = d_{\text{real}}/n$.
\end{enumerate}
\end{methode}

\begin{exoresolu}[Prism and Total Internal Reflection]
A ray of monochromatic light enters a glass prism of refracting angle $A = 60^\circ$ and refractive index $n = 1.5$ at an angle of incidence $i = 45^\circ$.
\begin{enumerate}
\item Calculate the angle of refraction $r_1$ at the first face.
\item Determine the angle of incidence $r_2$ at the second face.
\item Check if Total Internal Reflection occurs at the second face.
\item If TIR does not occur, calculate the angle of emergence $e$ and the total deviation $D$.
\end{enumerate}
\tcblower
\textbf{Detailed Solution}

\textbf{1. Refraction at First Face (Air to Glass)}
$n_{\text{air}} \sin i = n_{\text{glass}} \sin r_1$
$1 \times \sin 45^\circ = 1.5 \times \sin r_1$
$\sin r_1 = 0.7071 / 1.5 = 0.4714$
$r_1 = \arcsin(0.4714) = \SI{28.1}{\ensuremath{{}^{\circ}}}$.

\textbf{2. Angle of Incidence at Second Face ($r_2$)}
Geometry: $r_1 + r_2 = A = 60^\circ$.
$r_2 = 60^\circ - 28.1^\circ = \SI{31.9}{\ensuremath{{}^{\circ}}}$.

\textbf{3. Check for TIR at Second Face (Glass to Air)}
Critical angle $C$: $\sin C = \frac{n_{\text{air}}}{n_{\text{glass}}} = \frac{1}{1.5} = 0.6667$.
$C = \arcsin(0.6667) = \SI{41.8}{\ensuremath{{}^{\circ}}}$.
Compare: $r_2 = 31.9^\circ < 41.8^\circ = C$.
\textbf{Conclusion:} $r_2 < C$, so \textbf{TIR does NOT occur}. Ray emerges.

\textbf{4. Emergence Angle $e$ and Deviation $D$}
Snell's law at second face: $n_{\text{glass}} \sin r_2 = n_{\text{air}} \sin e$
$1.5 \times \sin 31.9^\circ = \sin e$
$\sin 31.9^\circ \approx 0.528$
$\sin e = 1.5 \times 0.528 = 0.792$
$e = \arcsin(0.792) = \SI{52.4}{\ensuremath{{}^{\circ}}}$.
Total Deviation $D = i + e - A = 45^\circ + 52.4^\circ - 60^\circ = \SI{37.4}{\ensuremath{{}^{\circ}}}$.
\end{exoresolu}

\begin{exercice}[Refraction and Prisms]
\begin{enumerate}
\item A ray of light in air strikes a glass block ($n=1.5$) at $30^\circ$. Find the angle of refraction. The ray emerges from the opposite parallel face. Show that the emergent ray is parallel to the incident ray but laterally displaced.
\item Calculate the critical angle for a glass-water interface ($n_g=1.5, n_w=4/3$). Will TIR occur for light going from water to glass?
\item A prism of angle $A=60^\circ$ and $n=1.5$ is used in minimum deviation. Find $D_{\min}$ and the angle of incidence $i$.
\item A fish is at a real depth of $2\ \mathrm{m}$ in water ($n=4/3$). What is its apparent depth when viewed from directly above?
\end{enumerate}
\end{exercice}

\begin{corrige}[Refraction and Prisms]
\begin{enumerate}
\item $\sin r = \sin 30^\circ / 1.5 = 0.333 \Rightarrow r = 19.5^\circ$. At second face, incidence $= r$, emergence $i' = 30^\circ$. Parallel. Displacement $d = t \frac{\sin(i-r)}{\cos r}$.
\item $\sin C = n_w/n_g = (4/3)/1.5 = 0.889 \Rightarrow C = 62.7^\circ$. Water to glass: rarer to denser $\rightarrow$ \textbf{No TIR possible}.
\item $n = \sin\frac{60+D_{\min}}{2} / \sin 30^\circ = 2\sin\frac{60+D_{\min}}{2} = 1.5 \Rightarrow \sin\frac{60+D_{\min}}{2} = 0.75 \Rightarrow \frac{60+D_{\min}}{2} = 48.6^\circ \Rightarrow D_{\min} = 37.2^\circ$. $i = (A+D_{\min})/2 = 48.6^\circ$.
\item $d_{\text{app}} = 2 / (4/3) = 1.5\ \mathrm{m}$.
\end{enumerate}
\end{corrige}

\coursec{Section 5: Lenses, Prisms and Optical Instruments}

\begin{definition}[Thin Lenses -- Sign Convention (Cartesian -- Standard for GCE)]
Light travels Left $\to$ Right.
\begin{itemize}
\item Object distance $u$: \textbf{Negative} (real object on left).
\item Image distance $v$: \textbf{Positive} (real image on right), \textbf{Negative} (virtual image on left).
\item Focal length $f$: \textbf{Positive} (converging/convex), \textbf{Negative} (diverging/concave).
\item Heights: $h_o$ positive (up), $h_i$ positive (up/erect), negative (inverted).
\end{itemize}
\textbf{Lens Formula:} $\frac{1}{v} - \frac{1}{u} = \frac{1}{f}$.
\textbf{Magnification:} $m = \frac{h_i}{h_o} = \frac{v}{u}$.
\textbf{Power:} $P = \frac{1}{f(\mathrm{m})}$ Dioptres (D). Combination (thin lenses in contact): $P = P_1 + P_2$.
\end{definition}

\begin{propriete}[Image Formation by Converging Lens ($f>0$)]
\begin{center}
\begin{tabular}{|c|c|c|c|c|}\hline
\rowcolor{popblueL}
\textbf{Object Position} & \textbf{Image Position} & \textbf{Nature} & \textbf{Magnification $|m|$} & \textbf{Use} \\ \hline
$u > 2f$ & $f < v < 2f$ & Real, Inverted & $< 1$ (Diminished) & Camera, Eye \\ \hline
$u = 2f$ & $v = 2f$ & Real, Inverted & $= 1$ (Same size) & Symmetric copying \\ \hline
$f < u < 2f$ & $v > 2f$ & Real, Inverted & $> 1$ (Magnified) & Projector, Objective of microscope/telescope \\ \hline
$u = f$ & $v = \infty$ & Real, at infinity & $\infty$ & Collimator, Searchlight \\ \hline
$u < f$ & $v < 0$ (virtual, left) & Virtual, Erect & $> 1$ (Magnified) & Magnifying glass, Eyepiece \\ \hline
\end{tabular}
\end{center}
\end{propriete}

\begin{experience}[Determining Focal Length of a Converging Lens]
\textbf{Method 1: Distant Object (Approx $f$).} Focus image of distant object (tree, building) on screen. $v \approx f$.
\textbf{Method 2: $u$-$v$ Method (No Parallax).} Use illuminated object (cross-wire), lens, screen. Vary $u$, find sharp image $v$. Plot $1/v$ vs $1/u$ (straight line, intercept $1/f$) or $uv$ vs $u+v$ (slope $=f$). Or use formula $f = \frac{uv}{u+v}$ for several pairs.
\textbf{Method 3: Displacement Method (Two positions for fixed $D=u+v > 4f$).} $f = \frac{D^2 - d^2}{4D}$ where $d$ = distance between two lens positions giving sharp images (one magnified, one diminished). $m_1 m_2 = 1$.
\end{experience}

\courssub{Optical Instruments}

\begin{propriete}[Compound Microscope]
\textbf{Objective:} Short $f_o$ ($\approx 0.5-1\ \mathrm{cm}$). Object placed just outside $f_o$ ($u_o \approx -f_o$). Forms real, inverted, magnified image $I_1$ at distance $v_o \approx L$ (tube length, typically $16-20\ \mathrm{cm}$).
$m_o = v_o/u_o \approx -L/f_o$.
\textbf{Eyepiece:} Short $f_e$ ($\approx 2-5\ \mathrm{cm}$). Acts as simple magnifier viewing $I_1$.
\begin{itemize}
\item \textbf{Normal Adjustment (Final image at $\infty$):} $I_1$ at focal point of eyepiece ($u_e = -f_e$). Eye relaxed. $M_e = D/f_e$ ($D=25\ \mathrm{cm}$ near point). $M = m_o \times M_e = -\frac{L}{f_o} \frac{D}{f_e}$.
\item \textbf{Near Point Adjustment (Final image at $D$):} $I_1$ inside $f_e$ ($u_e = -\frac{f_e D}{f_e+D}$). Eye accommodated. $M_e = 1 + D/f_e$. $M = m_o \times M_e = -\frac{L}{f_o}(1 + \frac{D}{f_e})$.
\end{itemize}
\textbf{Resolving Power (Abbe):} $d_{\min} = \frac{0.61\lambda}{\text{NA}}$, $\text{NA} = n\sin\alpha$ (Numerical Aperture). Oil immersion ($n=1.5$) increases NA.
\end{propriete}

\begin{propriete}[Astronomical Telescope (Refracting)]
\textbf{Objective:} Long $f_o$ (large aperture for light gathering and resolution). Forms real, inverted, diminished image at its focal plane ($v_o = f_o$).
\textbf{Eyepiece:} Short $f_e$. Views $I_1$.
\begin{itemize}
\item \textbf{Normal Adjustment (Final image at $\infty$):} $I_1$ at $F_e$ ($u_e = -f_e$). Separation $= f_o + f_e$. $M = -\frac{f_o}{f_e}$. Eye relaxed.
\item \textbf{Near Point Adjustment:} $I_1$ inside $F_e$ ($u_e = -\frac{f_e D}{f_e+D}$). Separation $= f_o + \frac{f_e D}{f_e+D} < f_o+f_e$. $M = -\frac{f_o}{f_e}(1 + \frac{f_e}{D}) = -\frac{f_o}{f_e} - \frac{f_o}{D}$. Eye accommodated.
\end{itemize}
\textbf{Resolving Power (Rayleigh):} $\theta_{\min} = 1.22 \lambda / D_{\text{obj}}$ (circular aperture). Depends on \textbf{objective diameter}, not magnification.
\textbf{Light Gathering Power:} $\propto D_{\text{obj}}^2$.
\textbf{Terrestrial Telescope:} Adds erecting lens/system $\rightarrow$ final image erect. Length increased by $4f_{\text{erect}}$.
\textbf{Galilean Telescope:} Diverging eyepiece ($f_e < 0$). Final image erect. $M = f_o/|f_e|$. Short length $= f_o - |f_e|$. Narrow field of view. Opera glasses.
\end{propriete}

\begin{methode}[Lens Formula, Magnification and Optical Instruments]
\begin{enumerate}
\item \textbf{Sign Convention:} \textbf{Cartesian (Real is Positive for $v,f$; $u$ negative).} Lens formula: $\frac{1}{v} - \frac{1}{u} = \frac{1}{f}$. Magnification $m = v/u$. Stick to ONE convention.
\item \textbf{Single Lens:} Find $v$, $m$, nature (real/virtual, inverted/erect, magnified/diminished).
\item \textbf{Compound Microscope:}
\begin{enumerate}
\item Analyse Eyepiece first (given final image position $v_e = -D$ or $\infty$) $\rightarrow$ find $u_e$ (object for eyepiece = image from objective).
\item $v_o = \text{Tube Length} - |u_e|$.
\item Use Lens Formula for Objective $\rightarrow$ find $u_o$.
\item $m_o = v_o/u_o$. $M_e = D/f_e$ (normal) or $1+D/f_e$ (near point).
\item $M = m_o \times M_e$.
\end{enumerate}
\item \textbf{Astronomical Telescope:}
\begin{enumerate}
\item Objective: $v_o = f_o$ (object at $\infty$).
\item Eyepiece: Normal $\rightarrow u_e = -f_e$, $M = -f_o/f_e$. Near Point $\rightarrow$ find $u_e$ from $v_e = -D$, $M = -\frac{f_o}{f_e}(1+\frac{f_e}{D})$.
\item Length = $f_o + |u_e|$.
\end{enumerate}
\end{enumerate}
\end{methode}

\begin{exoresolu}[Compound Microscope Calculation]
A compound microscope has an objective lens of focal length $f_o = 1.0\ \mathrm{cm}$ and an eyepiece of focal length $f_e = 5.0\ \mathrm{cm}$. The distance between the lenses (tube length) is $20\ \mathrm{cm}$. The final image is formed at the near point ($D = 25\ \mathrm{cm}$).
\begin{enumerate}
\item Calculate the object distance $u_o$ from the objective lens.
\item Determine the linear magnification produced by the objective.
\item Calculate the angular magnification of the microscope.
\end{enumerate}
\tcblower
\textbf{Detailed Solution}
\textit{Sign Convention: Cartesian ($u$ negative, $v$ positive for real image on right, $f$ positive for converging). Lens formula: $\frac{1}{v} - \frac{1}{u} = \frac{1}{f}$.}

\textbf{1. Eyepiece Analysis (Final image at Near Point $v_e = -25\ \mathrm{cm}$)}
$f_e = +5.0\ \mathrm{cm}$.
$\frac{1}{v_e} - \frac{1}{u_e} = \frac{1}{f_e} \Rightarrow \frac{1}{-25} - \frac{1}{u_e} = \frac{1}{5}$
$-\frac{1}{u_e} = 0.2 + 0.04 = 0.24 \Rightarrow u_e = -\frac{1}{0.24} = -4.167\ \mathrm{cm}$.
Object for eyepiece (image $I_1$ from objective) is $4.167\ \mathrm{cm}$ to the left of eyepiece.

\textbf{2. Objective Analysis}
Tube length $L = 20\ \mathrm{cm}$.
$v_o = L - |u_e| = 20 - 4.167 = 15.833\ \mathrm{cm}$ (real image, positive).
$f_o = +1.0\ \mathrm{cm}$.
$\frac{1}{v_o} - \frac{1}{u_o} = \frac{1}{f_o} \Rightarrow \frac{1}{15.833} - \frac{1}{u_o} = 1$
$-\frac{1}{u_o} = 1 - 0.06316 = 0.93684 \Rightarrow u_o = -1.067\ \mathrm{cm}$.
\textbf{Object distance from objective:} \SI{1.07}{cm} (to the left).

\textbf{3. Linear Magnification of Objective}
$m_o = \frac{v_o}{u_o} = \frac{15.833}{-1.067} = -14.84$.
(Image real, inverted, magnified 14.8 times).

\textbf{4. Angular Magnification of Microscope}
$M_e = 1 + \frac{D}{f_e} = 1 + \frac{25}{5} = 6$.
$M = m_o \times M_e = (-14.84) \times 6 = -89.0$.
\textbf{Magnitude: 89.} (Negative $\rightarrow$ final image inverted relative to original object).
\end{exoresolu}

\begin{exoresolu}[Astronomical Telescope Adjustments]
An astronomical telescope has an objective of focal length $f_o = 100\ \mathrm{cm}$ and diameter $10\ \mathrm{cm}$, and an eyepiece of focal length $f_e = 5.0\ \mathrm{cm}$. The telescope is used to view the moon (angular diameter $\approx 0.5^\circ$). Assume $\lambda = 550\ \mathrm{nm}$.
\begin{enumerate}
\item Calculate the magnifying power and length of the telescope in normal adjustment.
\item Calculate the magnifying power and length when the final image is formed at the near point ($D = 25\ \mathrm{cm}$).
\item Determine the angular resolution limit (Rayleigh criterion) of the telescope.
\item What is the angular size of the final image of the moon in normal adjustment?
\end{enumerate}
\tcblower
\textbf{Detailed Solution}

\textbf{1. Normal Adjustment}
$M_\infty = -\frac{f_o}{f_e} = -\frac{100}{5} = \textbf{-20}$ (Inverted).
$L_\infty = f_o + f_e = 100 + 5 = \textbf{105 cm}$.

\textbf{2. Near Point Adjustment ($D = 25\ \mathrm{cm}$)}
Eyepiece: $v_e = -25\ \mathrm{cm}$, $f_e = +5\ \mathrm{cm}$.
$\frac{1}{-25} - \frac{1}{u_e} = \frac{1}{5} \Rightarrow -\frac{1}{u_e} = 0.24 \Rightarrow u_e = -4.167\ \mathrm{cm}$.
$M_D = -\frac{f_o}{f_e}\left(1 + \frac{f_e}{D}\right) = -20 \times \left(1 + \frac{5}{25}\right) = -20 \times 1.2 = \textbf{-24}$.
$L_D = f_o + |u_e| = 100 + 4.167 = \textbf{104.17 cm}$.

\textbf{3. Angular Resolution (Rayleigh Criterion)}
$D_{\text{obj}} = 10\ \mathrm{cm} = 0.10\ \mathrm{m}$.
$\theta_{\min} = 1.22 \frac{\lambda}{D_{\text{obj}}} = 1.22 \times \frac{550 \times 10^{-9}}{0.10} = 6.71 \times 10^{-6}\ \mathrm{rad}$.
In arcseconds: $6.71 \times 10^{-6} \times 206265 \approx \textbf{1.38''}$.

\textbf{4. Angular Size of Moon's Image}
$\alpha = 0.5^\circ$. $\alpha' = |M| \alpha = 20 \times 0.5^\circ = \textbf{10}^\circ$.
\end{exoresolu}

\begin{attention}[Optical Instruments -- Common Mistakes]
\begin{itemize}
\item \textbf{Sign Convention:} Mixing Cartesian and "Real is Positive" causes sign errors. \textbf{Choose Cartesian and stick to it.}
\item \textbf{Microscope Tube Length:} $L$ is distance between lenses. $v_o = L - |u_e|$. Do not assume $v_o = L$ unless final image at $\infty$ ($u_e = -f_e$).
\item \textbf{Magnification Signs:} $m_o$ negative (inverted real image). $M_e$ positive (erect virtual image relative to its object). Total $M$ negative $\rightarrow$ final image inverted relative to original object.
\item \textbf{Telescope Magnification:} $M = -f_o/f_e$ (normal). Magnitude is $f_o/f_e$. Don't forget the negative sign for inversion.
\item \textbf{Resolving Power:} Depends on \textbf{Objective Diameter} ($D$), not focal length or magnification. $\theta_{\min} = 1.22\lambda/D$.
\end{itemize}
\end{attention}

\coursec{Section 6: Energy of Waves in Free Space and the Inverse Square Law}

\begin{propriete}[Wave Intensity and the Inverse Square Law]
\textbf{Intensity $I$:} Average power per unit area perpendicular to propagation direction. $I = \frac{P}{A}$ (\si{W.m^{-2}}).
\textbf{Point Source (Isotropic) in Free Space:} Power $P$ radiated uniformly over sphere of radius $r$.
\[
I = \frac{P}{4\pi r^2} \quad \Rightarrow \quad I \propto \frac{1}{r^2} \quad \text{(Inverse Square Law)}.
\]
\textbf{Amplitude:} For EM waves $I \propto E_0^2$; for sound $I \propto p_0^2$ (pressure amplitude). Since $I \propto 1/r^2$, \textbf{Amplitude $A \propto 1/r$}.
\textbf{Comparison at two distances:} $\frac{I_1}{I_2} = \left(\frac{r_2}{r_1}\right)^2$, $\frac{A_1}{A_2} = \frac{r_2}{r_1}$.
\textbf{Non-isotropic / Directional Sources:} Gain $G$ (directivity). $I = \frac{PG}{4\pi r^2}$ in direction of max radiation.
\textbf{Plane Waves / Laser Beams:} Ideally collimated, $I$ constant with $r$. Real beams diverge: far-field $I \propto 1/r^2$.
\end{propriete}

\begin{propriete}[Decibel Scale (Logarithmic Intensity Level)]
\textbf{Sound Intensity Level:} $\beta = 10 \log_{10}\left(\frac{I}{I_0}\right)$ dB, where $I_0 = 10^{-12}\ \mathrm{W\,m^{-2}}$ (threshold of hearing at $1\ \mathrm{kHz}$).
\textbf{Change in Level:} $\Delta \beta = \beta_1 - \beta_2 = 10 \log_{10}\left(\frac{I_1}{I_2}\right) = 20 \log_{10}\left(\frac{r_2}{r_1}\right)$ (for point source).
\textbf{Useful Rules:} Double distance $\rightarrow$ $-6\ \mathrm{dB}$. Double intensity $\rightarrow$ $+3\ \mathrm{dB}$. $10\times$ distance $\rightarrow$ $-20\ \mathrm{dB}$.
\end{propriete}

\begin{experience}[Verifying Inverse Square Law for Light / Sound]
\textbf{Light:} Point source (small filament lamp or LED with aperture), light meter (Lux meter), metre rule, dark room.
\textbf{Procedure:} Measure illuminance $E$ (lux) at distances $r = 20, 30, 40, 50, 60\ \mathrm{cm}$. Plot $E$ vs $1/r^2$ $\rightarrow$ straight line through origin. Slope $= P/4\pi$ (luminous flux).
\textbf{Sound:} Signal generator + speaker (approx point source at low freq), sound level meter (dB). Measure $\beta$ vs $r$. Plot $\beta$ vs $\log r$ or check $\Delta \beta = 20\log(r_2/r_1)$.
\textbf{Corrections:} Background light/noise subtraction. Source not perfectly point-like (near field effects). Reflections from walls (use anechoic chamber or open ground / time-gating).
\end{experience}

\begin{methode}[Inverse Square Law and Decibel Calculations]
\begin{enumerate}
\item \textbf{Identify geometry:} Point source $\rightarrow$ spherical ($4\pi r^2$). Source on hard floor $\rightarrow$ hemispherical ($2\pi r^2$, intensity $\times 2$, level $+3\ \mathrm{dB}$).
\item \textbf{Intensity:} $I = P/(4\pi r^2)$.
\item \textbf{Level:} $\beta = 10\log_{10}(I/I_0)$.
\item \textbf{Distance for given level:} Use ratio $\frac{I_1}{I_2} = (r_2/r_1)^2$ or $\Delta\beta = 20\log_{10}(r_2/r_1)$.
\item \textbf{Coherent Sources (Interference):} Amplitudes add. Constructive: $A_{\text{tot}} = 2A_1 \Rightarrow I_{\text{tot}} = 4I_1$ ($+6\ \mathrm{dB}$). Destructive: $I_{\text{tot}} = 0$ ($-\infty\ \mathrm{dB}$). Incoherent: $I_{\text{tot}} = 2I_1$ ($+3\ \mathrm{dB}$).
\end{enumerate}
\end{methode}

\begin{exoresolu}[Inverse Square Law and Decibels]
A point source emits sound waves of power $P = 50\ \mathrm{W}$ uniformly in all directions. The threshold of hearing is $I_0 = 10^{-12}\ \mathrm{W\,m^{-2}}$.
\begin{enumerate}
\item Calculate the intensity and sound intensity level (in dB) at a distance of $10\ \mathrm{m}$ from the source.
\item At what distance will the sound intensity level be $60\ \mathrm{dB}$?
\item If the source is placed on a hard reflecting floor (hemispherical radiation), what is the intensity level at $10\ \mathrm{m}$?
\item Two identical sources are placed $2\ \mathrm{m}$ apart, driven in phase. Calculate the intensity at a point $10\ \mathrm{m}$ from the midpoint along the perpendicular bisector (assume constructive interference).
\end{enumerate}
\tcblower
\textbf{Detailed Solution}

\textbf{1. Intensity and Level at $r_1 = 10\ \mathrm{m}$ (Free Space / Spherical)}
$I_1 = \frac{P}{4\pi r_1^2} = \frac{50}{4\pi (10)^2} = \frac{50}{400\pi} = \frac{1}{8\pi} \approx 0.0398\ \mathrm{W\,m^{-2}}$.
$\beta_1 = 10 \log_{10}\left(\frac{I_1}{I_0}\right) = 10 \log_{10}\left(\frac{0.0398}{10^{-12}}\right) = 10 \log_{10}(3.98 \times 10^{10})$
$\beta_1 = 10 (10 + \log_{10} 3.98) = 10 (10 + 0.60) = \textbf{106 dB}$.

\textbf{2. Distance for $\beta_2 = 60\ \mathrm{dB}$}
$\beta_2 = 60 \Rightarrow I_2/I_0 = 10^6 \Rightarrow I_2 = 10^{-6}\ \mathrm{W\,m^{-2}}$.
Inverse Square Law: $\frac{I_1}{I_2} = \left(\frac{r_2}{r_1}\right)^2$.
$\frac{0.0398}{10^{-6}} = 3.98 \times 10^4 = \left(\frac{r_2}{10}\right)^2$.
$\frac{r_2}{10} = \sqrt{3.98 \times 10^4} \approx 199.5$.
$r_2 \approx \textbf{1995 m} \approx \textbf{2.0 km}$.
(Check via dB: $\Delta \beta = 106 - 60 = 46\ \mathrm{dB}$. $46 = 20 \log_{10}(r_2/10) \Rightarrow \log_{10}(r_2/10) = 2.3 \Rightarrow r_2/10 = 10^{2.3} \approx 199.5$. Matches.)

\textbf{3. Hemispherical Radiation (on hard floor)}
Power into $2\pi$ sr: $I_{\text{hemi}} = \frac{P}{2\pi r^2} = 2 I_{\text{sphere}} = 0.0796\ \mathrm{W\,m^{-2}}$.
Level increase: $\Delta \beta = 10 \log_{10}(2) \approx 3\ \mathrm{dB}$.
$\beta_{\text{hemi}} = 106 + 3 = \textbf{109 dB}$.

\textbf{4. Two Sources Constructive Interference}
Point on perpendicular bisector $\Rightarrow$ path difference $= 0 \Rightarrow$ constructive.
Amplitudes add: $A_{\text{tot}} = 2A_1 \Rightarrow I_{\text{tot}} = 4 I_1 = 4 \times 0.0398 = 0.159\ \mathrm{W\,m^{-2}}$.
Level: $\beta_{\text{tot}} = 10 \log_{10}(4 I_1 / I_0) = 10 \log_{10} 4 + \beta_1 = 6.02 + 106 \approx \textbf{112 dB}$.
(Note: Total power $100\ \mathrm{W}$, but directivity creates lobes. On axis, intensity is 4x single source).
\end{exoresolu}

\begin{exercice}[Inverse Square Law and Energy]
\begin{enumerate}
\item A laser pointer emits a beam of power $5\ \mathrm{mW}$ and diameter $2\ \mathrm{mm}$. Calculate the intensity (irradiance) of the beam. Why does the inverse square law not apply close to the laser?
\item A loudspeaker radiates $10\ \mathrm{W}$ of acoustic power. Calculate the sound level at $5\ \mathrm{m}$. If a second identical speaker is placed next to the first and driven in phase, what is the level at $5\ \mathrm{m}$ on the axis? What if they are driven incoherently?
\item The intensity of sunlight at Earth's orbit (distance $1.5 \times 10^{11}\ \mathrm{m}$) is $1360\ \mathrm{W\,m^{-2}}$ (Solar Constant). Estimate the total power output of the Sun.
\item A sound level meter reads $90\ \mathrm{dB}$ at $2\ \mathrm{m}$ from a machine. What will it read at $10\ \mathrm{m}$ assuming free-field conditions?
\end{enumerate}
\end{exercice}

\begin{corrige}[Inverse Square Law and Energy]
\begin{enumerate}
\item $A = \pi (1\times 10^{-3})^2 = \pi \times 10^{-6}\ \mathrm{m^2}$. $I = 5\times 10^{-3} / (\pi \times 10^{-6}) \approx 1590\ \mathrm{W\,m^{-2}}$. Inverse square law applies only in far field (beyond Rayleigh range $z_R = \pi w_0^2/\lambda$). Near field: beam roughly collimated.
\item Single: $I = 10/(4\pi 25) = 0.0318\ \mathrm{W\,m^{-2}}$. $\beta = 10\log(0.0318/10^{-12}) = 105\ \mathrm{dB}$.
Coherent, in phase, on axis: $I_{\text{tot}} = 4I = 0.127\ \mathrm{W\,m^{-2}}$, $\beta = 105 + 6 = 111\ \mathrm{dB}$.
Incoherent: $I_{\text{tot}} = 2I = 0.0637\ \mathrm{W\,m^{-2}}$, $\beta = 105 + 3 = 108\ \mathrm{dB}$.
\item $P = I \times 4\pi r^2 = 1360 \times 4\pi (1.5\times 10^{11})^2 \approx 3.85 \times 10^{26}\ \mathrm{W}$.
\item $\Delta r = 5\times$. $\Delta \beta = -20\log_{10}(5) = -14\ \mathrm{dB}$. Reading $= 90 - 14 = 76\ \mathrm{dB}$.
\end{enumerate}
\end{corrige}

\begin{aretenir}[Chapter Summary -- Electromagnetic Waves]
\begin{itemize}
\item \textbf{Sound speed:} Resonance tube $v = 2f(L_2-L_1)$. Temp correction $v \propto \sqrt{T}$.
\item \textbf{EM Waves:} Transverse, $c=3\times10^8\ \mathrm{m/s}$, $E_0=cB_0$, $I = \frac{1}{2}\varepsilon_0 c E_0^2$. Spectrum: Radio $\to$ $\gamma$.
\item \textbf{X-rays:} Coolidge tube. Bremsstrahlung (continuous, $\lambda_{\min}=hc/eV$) + Characteristic lines. Hardness controlled by kVp, intensity by mA.
\item \textbf{Polarisation:} Malus $I=I_0\cos^2\theta$. Brewster $\tan i_p=n$, reflected ray perfectly polarised ($\vec{E}\parallel$ interface).
\item \textbf{Diffraction Grating:} $d\sin\theta=n\lambda$. $n_{\max}=\lfloor d/\lambda\rfloor$. Resolving Power $R=nN_{\text{total}}$.
\item \textbf{Refraction:} Snell $n_1\sin i=n_2\sin r$. Critical angle $\sin C=n_2/n_1$ (denser$\to$rarer). TIR if $i>C$. Prism $n=\sin\frac{A+D_{\min}}{2}/\sin\frac{A}{2}$.
\item \textbf{Lenses:} Cartesian: $\frac{1}{v}-\frac{1}{u}=\frac{1}{f}$, $m=v/u$. Power $P=1/f$ (D).
\item \textbf{Microscope:} $M = -\frac{L}{f_o}(1+\frac{D}{f_e})$ (near point). $M = -\frac{L}{f_o}\frac{D}{f_e}$ (normal).
\item \textbf{Telescope:} $M = -f_o/f_e$ (normal). $M = -\frac{f_o}{f_e}(1+\frac{f_e}{D})$ (near point). Resolution $\theta_{\min}=1.22\lambda/D_{\text{obj}}$.
\item \textbf{Inverse Square Law:} $I=P/4\pi r^2$, $A\propto 1/r$. $\beta=10\log(I/I_0)$. $\Delta\beta=20\log(r_2/r_1)$.
\end{itemize}
\end{aretenir}

\newpage

\coursec{Exercises}

\courssub{I check my knowledge}

\begin{exercice}[Multiple choice questions (GCE Paper 1 style)]
For each question, choose the single correct answer.
\begin{enumerate}
\item Which of the following lists electromagnetic waves in order of \emph{increasing} wavelength?
\begin{enumerate}
\item visible light, X-rays, radio waves, gamma rays
\item gamma rays, X-rays, visible light, radio waves
\item radio waves, visible light, X-rays, gamma rays
\item X-rays, gamma rays, radio waves, visible light
\end{enumerate}
\item Plane polarisation is a property that can be demonstrated with:
\begin{enumerate}
\item sound waves only
\item longitudinal waves only
\item transverse waves only
\item all types of waves
\end{enumerate}
\item Total internal reflection can occur only when light travels:
\begin{enumerate}
\item from air into glass at any angle of incidence
\item from glass into air at an angle of incidence greater than the critical angle
\item from glass into air at an angle of incidence less than the critical angle
\item from air into glass at the critical angle
\end{enumerate}
\item For a diffraction grating of slit spacing $d$ used at normal incidence, the maxima satisfy:
\begin{enumerate}
\item $d\sin\theta = n\lambda$
\item $d\sin\theta = (n+\tfrac{1}{2})\lambda$
\item $\lambda\sin\theta = nd$
\item $d = n\lambda\sin\theta$
\end{enumerate}
\item The intensity of a wave emitted by a point source in free space varies with distance $r$ from the source as:
\begin{enumerate}
\item $I \propto r$
\item $I \propto r^{2}$
\item $I \propto \dfrac{1}{r}$
\item $I \propto \dfrac{1}{r^{2}}$
\end{enumerate}
\end{enumerate}
\end{exercice}

\begin{exercice}[True or false — justify your answer]
State whether each statement is true or false, and justify your answer in one or two sentences.
\begin{enumerate}
\item Sound is an electromagnetic wave because it can travel through air.
\item In a resonance tube experiment, the speed of sound in air is given by $v = 2f(l_{2}-l_{1})$, where $l_{1}$ and $l_{2}$ are the lengths of the air column for the first two successive resonance positions.
\item X-rays are deflected by a strong magnetic field because they carry electric charge.
\item When light passes from air into glass, its frequency decreases while its wavelength stays constant.
\item The power of a converging lens of focal length $10\ \mathrm{cm}$ is $+10$ dioptres.
\end{enumerate}
\end{exercice}

\begin{exercice}[Definitions]
Define each of the following terms precisely, giving the appropriate formula or unit where relevant:
\begin{enumerate}
\item the \cle{refractive index} of a medium;
\item the \cle{critical angle} for a boundary between two media;
\item a \cle{plane polarised} electromagnetic wave;
\item the \cle{dioptre};
\item the \cle{intensity} of a wave at a point in free space.
\end{enumerate}
\end{exercice}

\begin{exercice}[Direct calculations]
\begin{enumerate}
\item A tuning fork of frequency $512\ \mathrm{Hz}$ is held over a resonance tube. Successive resonances occur for air columns of lengths $16.2\ \mathrm{cm}$ and $49.8\ \mathrm{cm}$. Calculate the speed of sound in air.
\item Light enters a block of glass of refractive index $1.52$ at an angle of incidence of $45^{\circ}$. Calculate the angle of refraction.
\item A lamp produces an intensity of $8.0\ \mathrm{W\,m^{-2}}$ at a distance of $2.0\ \mathrm{m}$. Calculate the intensity at a distance of $5.0\ \mathrm{m}$, assuming the lamp is a point source radiating uniformly in all directions.
\end{enumerate}
\end{exercice}

\courssub{I practise}

\begin{exercice}[The diffraction grating]
A plane transmission grating having $500$ lines per millimetre is illuminated at normal incidence by sodium light.
\begin{enumerate}
\item Calculate the slit spacing $d$ of the grating in metres.
\item The second-order maximum is observed at an angle of $36.1^{\circ}$ to the normal. Calculate the wavelength of the sodium light.
\item Determine the highest order of maximum that can be observed with this grating and this light.
\end{enumerate}
\end{exercice}

\begin{exercice}[Refraction and critical angle]
A ray of light in air strikes the plane surface of a glass block of refractive index $1.52$ at an angle of incidence of $45^{\circ}$.
\begin{enumerate}
\item Calculate the angle of refraction inside the glass.
\item Calculate the speed of light in the glass. (Speed of light in vacuum: $c = 3.00\times10^{8}\ \mathrm{m\,s^{-1}}$.)
\item Calculate the critical angle for a glass--air boundary.
\item State what happens to a ray inside the glass that strikes the boundary at an angle of incidence of $50^{\circ}$.
\end{enumerate}
\end{exercice}

\begin{exercice}[The converging lens]
An object of height $4.0\ \mathrm{cm}$ is placed $30\ \mathrm{cm}$ from a thin converging lens of focal length $10\ \mathrm{cm}$.
\begin{enumerate}
\item Using the lens formula, calculate the position of the image.
\item Calculate the linear magnification and hence the height of the image.
\item State the nature of the image (real or virtual, upright or inverted, magnified or diminished).
\item Calculate the power of the lens in dioptres.
\item Describe one practical situation in everyday life in which a single converging lens is used to produce a magnified image.
\end{enumerate}
\end{exercice}

\begin{exercice}[The prism]
A glass prism of refracting angle $A = 60^{\circ}$ produces a minimum deviation $D_{\min} = 37^{\circ}$ for a certain monochromatic light.
\begin{enumerate}
\item Write down the relation between the refractive index $n$, the prism angle $A$ and the minimum deviation $D_{\min}$.
\item Calculate the refractive index of the glass for this light.
\item Show that at minimum deviation the ray passes symmetrically through the prism, and calculate the angle of incidence at the first face.
\item Explain briefly why white light is dispersed by the prism.
\end{enumerate}
\end{exercice}

\courssub{I prepare for the GCE}

\begin{exercice}[Structured question: measuring the speed of sound and the inverse square law]
\begin{enumerate}
\item Describe, with the aid of a labelled diagram, the resonance tube experiment used to measure the speed of sound in air. Your description should include the apparatus, the procedure, and the measurements taken. \hfill (6 marks)
\item Show that the speed of sound is given by $v = 2f(l_{2}-l_{1})$, where $f$ is the frequency of the tuning fork and $l_{1}$, $l_{2}$ are the lengths of the air column at two successive resonance positions. \hfill (3 marks)
\item State one precaution you would take to obtain an accurate result, and one possible source of error. \hfill (2 marks)
\item In such an experiment, a tuning fork of frequency $480\ \mathrm{Hz}$ gives successive resonances at $17.5\ \mathrm{cm}$ and $52.9\ \mathrm{cm}$. Calculate the speed of sound in air. \hfill (2 marks)
\item A small loudspeaker emits sound uniformly in all directions in free space. Explain what is meant by the \emph{inverse square law} for the intensity of the sound, and derive it from the principle of conservation of energy. \hfill (4 marks)
\item The intensity is measured as $2.0\times10^{-3}\ \mathrm{W\,m^{-2}}$ at $4.0\ \mathrm{m}$ from the loudspeaker. Calculate the total power emitted by the loudspeaker and the distance at which the intensity falls to $5.0\times10^{-4}\ \mathrm{W\,m^{-2}}$. State the assumptions you make. \hfill (3 marks)
\end{enumerate}
\end{exercice}

\begin{exercice}[Structured question: X-rays and the electromagnetic spectrum]
\begin{enumerate}
\item Draw a labelled diagram of a Coolidge tube (modern X-ray tube) and explain how X-rays are produced in it. \hfill (6 marks)
\item State two properties of X-rays that distinguish them from visible light. \hfill (2 marks)
\item State two uses of X-rays, one in medicine and one in industry, and explain in each case which property of X-rays makes the use possible. \hfill (4 marks)
\item State one safety precaution that must be taken by people who work with X-rays, and explain why it is necessary. \hfill (2 marks)
\item An X-ray tube operates at an accelerating voltage of $50\ \mathrm{kV}$. Calculate the maximum energy acquired by an electron, and hence the minimum wavelength of the X-rays produced. (Electronic charge $e = 1.60\times10^{-19}\ \mathrm{C}$; Planck constant $h = 6.63\times10^{-34}\ \mathrm{J\,s}$; $c = 3.00\times10^{8}\ \mathrm{m\,s^{-1}}$.) \hfill (4 marks)
\item X-rays can be plane polarised. Explain what is meant by \emph{plane polarisation} of an electromagnetic wave, and explain why the phenomenon of polarisation is evidence that X-rays are transverse waves. \hfill (2 marks)
\end{enumerate}
\end{exercice}

\begin{exercice}[Structured question: optical instruments]
An astronomical telescope consists of an objective lens of focal length $100\ \mathrm{cm}$ and an eyepiece of focal length $5.0\ \mathrm{cm}$.
\begin{enumerate}
\item Explain what is meant by \emph{normal adjustment} of an astronomical telescope, and draw a ray diagram showing the passage through the telescope of two rays from a distant point object off the axis. \hfill (5 marks)
\item Show that in normal adjustment the magnifying power is $M = \dfrac{f_{o}}{f_{e}}$, and calculate its value for this telescope. \hfill (3 marks)
\item Calculate the distance between the two lenses (the length of the telescope) in normal adjustment. \hfill (2 marks)
\item Explain the advantage of using an objective of large diameter in an astronomical telescope. \hfill (2 marks)
\item A compound microscope uses an objective of focal length $1.0\ \mathrm{cm}$ and an eyepiece of focal length $5.0\ \mathrm{cm}$. State, with reasons, why the objective of a microscope must have a very short focal length whereas that of an astronomical telescope must be long. \hfill (4 marks)
\item The lenses of these instruments are made of glass of refractive index $1.50$. Calculate the critical angle at a glass--air surface, and explain why a right-angled prism of this glass can be used as a total internal reflection prism in optical instruments such as binoculars. \hfill (4 marks)
\end{enumerate}
\end{exercice}

\newpage

\coursec{Solutions to the exercises}

\begin{corrige}[Exercise 1 — Multiple choice questions]
\begin{enumerate}
\item \textbf{Correct answer: (b)} gamma rays, X-rays, visible light, radio waves.\par
In the electromagnetic spectrum, frequency decreases (and therefore wavelength $\lambda = c/f$ increases) in the order: gamma rays $\to$ X-rays $\to$ ultraviolet $\to$ visible light $\to$ infrared $\to$ microwaves $\to$ radio waves. Hence the order of \emph{increasing wavelength} is: gamma rays, X-rays, visible light, radio waves.

\item \textbf{Correct answer: (c)} transverse waves only.\par
Polarisation is the restriction of the vibrations of a wave to one plane. A longitudinal wave vibrates only along the direction of propagation, so there is no transverse direction to select and it cannot be polarised. Only transverse waves (light, X-rays, waves on a rope) can be plane polarised. The fact that light can be polarised is in fact the experimental evidence that light is a transverse wave.

\item \textbf{Correct answer: (b)} from glass into air at an angle of incidence greater than the critical angle.\par
Total internal reflection requires two conditions simultaneously: (i) the light must travel from the \emph{optically denser} medium (glass) towards the \emph{less dense} medium (air), i.e.\ $n_{1}>n_{2}$; (ii) the angle of incidence must exceed the critical angle, $i>C$, where $\sin C = n_{2}/n_{1}$.

\item \textbf{Correct answer: (a)} $d\sin\theta = n\lambda$.\par
For a transmission grating at normal incidence, the path difference between rays from adjacent slits is $d\sin\theta$; constructive interference (principal maxima) occurs when this equals a whole number of wavelengths: $d\sin\theta_{n}=n\lambda$, with $n=0,1,2,\dots$

\item \textbf{Correct answer: (d)} $I\propto \dfrac{1}{r^{2}}$.\par
By conservation of energy, the power $P$ of a point source radiating uniformly in free space is spread over the surface of a sphere of area $4\pi r^{2}$, so $I = \dfrac{P}{4\pi r^{2}} \propto \dfrac{1}{r^{2}}$: the inverse square law.
\end{enumerate}
\end{corrige}

\begin{corrige}[Exercise 2 — True or false]
\begin{enumerate}
\item \textbf{False.} Sound is a \emph{mechanical longitudinal} wave: it needs a material medium (air, water, solids) whose particles vibrate. Electromagnetic waves are transverse, do not need a medium, and travel in vacuum at $c = 3.00\times10^{8}\ \mathrm{m\,s^{-1}}$. Travelling through air does not make a wave electromagnetic.

\item \textbf{True.} Successive resonance positions of the air column differ by half a wavelength: $l_{2}-l_{1} = \dfrac{\lambda}{2}$, so $\lambda = 2(l_{2}-l_{1})$ and $v = f\lambda = 2f(l_{2}-l_{1})$. Using two successive positions also eliminates the end-correction of the tube, which is the main systematic error of the method.

\item \textbf{False.} X-rays are electromagnetic waves (oscillating electric and magnetic fields), not charged particles; they carry no electric charge and are therefore \emph{not} deflected by electric or magnetic fields. It is the cathode rays — the electrons inside the tube — that are deflected by fields.

\item \textbf{False.} The \emph{frequency} of light is fixed by the source and does not change on refraction. It is the speed that decreases ($v = c/n$) and therefore the wavelength decreases: $\lambda_{\text{glass}} = \lambda_{\text{air}}/n$.

\item \textbf{True.} The power of a lens is $P = \dfrac{1}{f}$ with $f$ in metres: $f = 10\ \mathrm{cm} = 0.10\ \mathrm{m}$, so $P = \dfrac{1}{0.10} = +10\ \mathrm{D}$ (dioptres). The sign is positive because the lens is converging.
\end{enumerate}
\end{corrige}

\begin{corrige}[Exercise 3 — Definitions]
\begin{enumerate}
\item \textbf{Refractive index.} The (absolute) refractive index $n$ of a medium is the ratio of the speed of light in vacuum to its speed in the medium:
\[ n = \frac{c}{v}. \]
It is a pure number (no unit), always greater than 1 for material media. Equivalently, for light entering the medium from air, $n = \dfrac{\sin i}{\sin r}$ (Snell's law).

\item \textbf{Critical angle.} For light travelling from a denser medium (index $n_{1}$) towards a less dense medium (index $n_{2}<n_{1}$), the critical angle $C$ is the angle of incidence in the denser medium for which the angle of refraction is $90^{\circ}$ (the refracted ray grazes the surface):
\[ \sin C = \frac{n_{2}}{n_{1}} \qquad (\text{for glass--air: } \sin C = \tfrac{1}{n}). \]

\item \textbf{Plane polarised wave.} An electromagnetic wave is plane polarised when the oscillations of its electric field vector $\vec{E}$ are confined to a single fixed direction (one plane) perpendicular to the direction of propagation; the magnetic field $\vec{B}$ then oscillates in the perpendicular plane.

\item \textbf{Dioptre.} The dioptre (symbol D) is the SI unit of the power of a lens (or curved surface): $P = \dfrac{1}{f}$ with $f$ in metres. One dioptre is the power of a lens of focal length one metre: $1\ \mathrm{D} = 1\ \mathrm{m^{-1}}$.

\item \textbf{Intensity.} The intensity $I$ of a wave at a point is the power received per unit area perpendicular to the direction of propagation:
\[ I = \frac{P}{A}, \qquad \text{unit: watt per square metre } (\mathrm{W\,m^{-2}}). \]
For a point source radiating uniformly in free space, $I = \dfrac{P}{4\pi r^{2}}$.
\end{enumerate}
\end{corrige}

\begin{corrige}[Exercise 4 — Direct calculations]
\begin{enumerate}
\item \textbf{Speed of sound.} Successive resonances differ by half a wavelength:
\[ \frac{\lambda}{2} = l_{2}-l_{1} = 49.8 - 16.2 = 33.6\ \mathrm{cm} = 0.336\ \mathrm{m}. \]
\[ v = f\lambda = 2f(l_{2}-l_{1}) = 2\times 512 \times 0.336 = 344\ \mathrm{m\,s^{-1}}. \]
This is consistent with the accepted value of about $340\ \mathrm{m\,s^{-1}}$ in air at room temperature.

\item \textbf{Angle of refraction.} Snell's law at the air--glass boundary:
\[ n_{\text{air}}\sin i = n\sin r \;\Rightarrow\; \sin r = \frac{\sin 45^{\circ}}{1.52} = \frac{0.7071}{1.52} = 0.4652. \]
\[ r = \sin^{-1}(0.4652) \approx 27.7^{\circ}. \]
As expected, the ray bends \emph{towards} the normal on entering the denser medium.

\item \textbf{Intensity at 5.0 m.} Inverse square law: $I r^{2} = \text{constant}$, so
\[ I_{2} = I_{1}\left(\frac{r_{1}}{r_{2}}\right)^{2} = 8.0\times\left(\frac{2.0}{5.0}\right)^{2} = 8.0\times 0.16 = 1.28\ \mathrm{W\,m^{-2}} \approx 1.3\ \mathrm{W\,m^{-2}}. \]
\end{enumerate}
\end{corrige}

\begin{corrige}[Exercise 5 — The diffraction grating]
\begin{enumerate}
\item \textbf{Slit spacing.} With $N = 500$ lines per millimetre:
\[ d = \frac{1}{N} = \frac{1}{500\ \mathrm{mm^{-1}}} = \frac{10^{-3}}{500}\ \mathrm{m} = 2.0\times10^{-6}\ \mathrm{m}. \]

\item \textbf{Wavelength.} Grating equation at normal incidence, $d\sin\theta_{n} = n\lambda$, with $n = 2$ and $\theta_{2} = 36.1^{\circ}$:
\[ \lambda = \frac{d\sin\theta_{2}}{2} = \frac{2.0\times10^{-6}\times\sin 36.1^{\circ}}{2} = \frac{2.0\times10^{-6}\times 0.589}{2} \approx 5.89\times10^{-7}\ \mathrm{m} = 589\ \mathrm{nm}. \]
This is the well-known yellow light of a sodium lamp.

\item \textbf{Highest observable order.} A maximum is observable only if $\sin\theta_{n}\leq 1$:
\[ n \leq \frac{d}{\lambda} = \frac{2.0\times10^{-6}}{5.89\times10^{-7}} = 3.40. \]
Since $n$ must be a whole number, the highest order is $n = 3$ (observed at $\sin\theta_{3} = 3\lambda/d = 0.884$, i.e.\ $\theta_{3}\approx 62.1^{\circ}$). The fourth order would require $\sin\theta_{4} = 1.18 > 1$: impossible. Counting the central maximum and both sides, seven bright lines are seen in total ($n = 0, \pm1, \pm2, \pm3$).
\end{enumerate}
\end{corrige}

\begin{corrige}[Exercise 6 — Refraction and critical angle]
\begin{enumerate}
\item \textbf{Angle of refraction.} Snell's law: $\sin r = \dfrac{\sin 45^{\circ}}{1.52} = 0.4652$, so $r \approx 27.7^{\circ}$.

\item \textbf{Speed of light in the glass.}
\[ v = \frac{c}{n} = \frac{3.00\times10^{8}}{1.52} = 1.97\times10^{8}\ \mathrm{m\,s^{-1}}. \]

\item \textbf{Critical angle.} For a glass--air boundary:
\[ \sin C = \frac{1}{n} = \frac{1}{1.52} = 0.658 \;\Rightarrow\; C \approx 41.1^{\circ}. \]

\item \textbf{Ray incident at $50^{\circ}$ inside the glass.} Since $50^{\circ} > C \approx 41.1^{\circ}$, and the ray travels from the denser medium (glass) towards air, \emph{total internal reflection} occurs: no refracted ray emerges; the ray is entirely reflected back into the glass, obeying the laws of reflection (angle of reflection $= 50^{\circ}$).
\end{enumerate}
\end{corrige}

\begin{corrige}[Exercise 7 — The converging lens]
We use the lens formula with the ``real is positive'' convention: $\dfrac{1}{f} = \dfrac{1}{u} + \dfrac{1}{v}$, with $u = +30\ \mathrm{cm}$, $f = +10\ \mathrm{cm}$.
\begin{enumerate}
\item \textbf{Image position.}
\[ \frac{1}{v} = \frac{1}{f} - \frac{1}{u} = \frac{1}{10} - \frac{1}{30} = \frac{3-1}{30} = \frac{2}{30} \;\Rightarrow\; v = +15\ \mathrm{cm}. \]
The image forms $15\ \mathrm{cm}$ from the lens on the side opposite the object.

\item \textbf{Magnification and image height.}
\[ m = \frac{v}{u} = \frac{15}{30} = 0.50, \qquad h_{i} = m\,h_{o} = 0.50\times 4.0 = 2.0\ \mathrm{cm}. \]

\item \textbf{Nature of the image.} Since $u > 2f$ (here $u = 30\ \mathrm{cm} > 2f = 20\ \mathrm{cm}$), the image is \textbf{real} ($v>0$), \textbf{inverted} and \textbf{diminished} ($m<1$) — the situation of a camera.

\item \textbf{Power of the lens.}
\[ P = \frac{1}{f\ (\mathrm{m})} = \frac{1}{0.10} = +10\ \mathrm{D}. \]

\item \textbf{Practical magnifying situation.} A \emph{magnifying glass} (simple microscope): when the object is placed \emph{inside} the focal length ($u < f$), the lens produces an upright, virtual, magnified image — used to read small print, to examine a specimen, or by a watchmaker. The eyepieces of compound microscopes and astronomical telescopes are used in exactly this way.
\end{enumerate}
\end{corrige}

\begin{corrige}[Exercise 8 — The prism]
\begin{enumerate}
\item \textbf{Relation at minimum deviation.}
\[ n = \frac{\sin\left(\dfrac{A + D_{\min}}{2}\right)}{\sin\left(\dfrac{A}{2}\right)}. \]

\item \textbf{Refractive index.}
\[ n = \frac{\sin\left(\dfrac{60^{\circ}+37^{\circ}}{2}\right)}{\sin 30^{\circ}} = \frac{\sin 48.5^{\circ}}{0.500} = \frac{0.749}{0.500} \approx 1.50. \]

\item \textbf{Symmetry at minimum deviation.} The deviation is $D = i_{1} + i_{2} - A$. For a given prism and a given wavelength, the same deviation $D$ is obtained for two angles of incidence, $i_{1}$ and $i_{2}$ interchanged (principle of reversibility of light). As $D$ passes through its minimum, these two solutions merge into one, which requires $i_{1} = i_{2} = i$; then by Snell's law $r_{1} = r_{2} = r$, and since $r_{1}+r_{2} = A$, we get $r = A/2 = 30^{\circ}$: the ray passes symmetrically through the prism, parallel to the base. Then
\[ D_{\min} = 2i - A \;\Rightarrow\; i = \frac{A + D_{\min}}{2} = \frac{60^{\circ}+37^{\circ}}{2} = 48.5^{\circ}. \]
(Check: $\sin i = n\sin r = 1.50\times\sin 30^{\circ} = 0.75$, so $i = 48.6^{\circ}$ — consistent with the measured value.)

\item \textbf{Dispersion of white light.} The refractive index of glass depends on wavelength (it is larger for violet than for red). Each colour of white light is therefore deviated by a different amount ($D_{\min}$ is larger for violet than for red), so the prism separates white light into its spectrum: this is \emph{dispersion}.
\end{enumerate}
\end{corrige}

\begin{corrige}[Exercise 9 — Structured question: speed of sound and inverse square law]
\textbf{(a) Description of the experiment (6 marks).}\par
\emph{Apparatus:} a tall glass resonance tube open at the top, connected at the bottom by a rubber hose to a water reservoir whose height can be adjusted (so the water level — and hence the length of the air column — can be varied); a set of tuning forks of known frequency $f$; a metre rule fixed alongside the tube.
\begin{center}
\begin{tikzpicture}[scale=0.85]
% tube
\draw[thick] (0,0) -- (0,6.5);
\draw[thick] (1.2,0) -- (1.2,6.5);
\draw[thick] (0,0) -- (1.2,0);
% water
\fill[popblueL] (0.02,0.02) rectangle (1.18,3.4);
\draw[popblue,thick] (0.02,3.4) -- (1.18,3.4);
\node[popblue] at (0.6,1.6) {water};
% air column label
\draw[<->,popdark] (1.35,3.4) -- (1.35,6.5) node[midway,right]{$l$};
% tuning fork
\draw[thick,poporange] (0.45,7.0) -- (0.45,7.9);
\draw[thick,poporange] (0.75,7.0) -- (0.75,7.9);
\draw[thick,poporange] (0.45,7.0) -- (0.75,7.0);
\draw[thick,poporange] (0.6,7.0) -- (0.6,6.6);
\node[poporange] at (2.6,7.4) {tuning fork ($f$)};
% reservoir
\draw[thick] (3.4,0.4) -- (3.4,2.6);
\draw[thick] (4.6,0.4) -- (4.6,2.6);
\draw[thick] (3.4,0.4) -- (4.6,0.4);
\fill[popblueL] (3.42,0.42) rectangle (4.58,2.0);
\node at (4.0,3.0) {reservoir};
% hose
\draw[thick] (0.6,0) .. controls (1.4,-0.7) and (3.2,-0.5) .. (4.0,0.4);
% ruler
\draw[thick,popmuted] (-0.5,0) -- (-0.5,6.5);
\foreach \y in {0,1,2,3,4,5,6} \draw[popmuted] (-0.5,\y) -- (-0.35,\y);
\node[popmuted,rotate=90] at (-0.85,3.2) {metre rule};
\end{tikzpicture}
\legende{Resonance tube apparatus}
\end{center}
\emph{Procedure:} strike the tuning fork and hold it above the open end of the tube. Starting with a high water level (short air column), lower the reservoir slowly until the sound becomes suddenly much louder: this is the first resonance; measure the length $l_{1}$ of the air column. Lower the water further until a second, louder resonance is heard; measure $l_{2}$. Repeat with other tuning forks and average. \hfill [diagram 2; apparatus 1; procedure 2; measurements 1]

\textbf{(b) Derivation of $v = 2f(l_{2}-l_{1})$ (3 marks).}\par
At resonance the air column carries a stationary wave with a node at the water surface (closed end) and an antinode at the open end. The first resonance corresponds to a quarter wavelength in the tube: $l_{1} + e = \dfrac{\lambda}{4}$, where $e$ is the small end-correction at the open end. The second resonance corresponds to three quarters of a wavelength: $l_{2} + e = \dfrac{3\lambda}{4}$. Subtracting eliminates $e$:
\[ l_{2} - l_{1} = \frac{3\lambda}{4} - \frac{\lambda}{4} = \frac{\lambda}{2} \;\Rightarrow\; \lambda = 2(l_{2}-l_{1}), \qquad v = f\lambda = 2f(l_{2}-l_{1}). \]
\hfill [node/antinode idea 1; the two equations 1; subtraction and result 1]

\textbf{(c) Precaution and source of error (2 marks).}\par
\emph{Precaution:} lower the water level slowly and locate the position of \emph{maximum} loudness carefully, approaching it from above and below and averaging (or repeat each measurement several times). \emph{Source of error:} the end-correction at the open end (eliminated by using $l_{2}-l_{1}$); the difficulty of judging the exact resonance position by ear; or the temperature of the air differing from the assumed value, since the speed of sound depends on temperature. \hfill [1 + 1]

\textbf{(d) Numerical value (2 marks).}
\[ v = 2f(l_{2}-l_{1}) = 2\times 480\times(0.529 - 0.175) = 960\times 0.354 = 339.8\ \mathrm{m\,s^{-1}} \approx 340\ \mathrm{m\,s^{-1}}. \]
\hfill [substitution 1; value with unit 1]

\textbf{(e) Inverse square law (4 marks).}\par
Statement: the intensity of the sound from a point source in free space is inversely proportional to the square of the distance from the source: $I\propto 1/r^{2}$.\par
\emph{Derivation from conservation of energy:} the loudspeaker emits sound power $P$ uniformly in all directions. In free space there is no absorption, so energy is conserved: all the power crossing one sphere centred on the source also crosses any larger sphere. At distance $r$, this power is spread over the surface of a sphere of area $A = 4\pi r^{2}$. Hence
\[ I = \frac{P}{A} = \frac{P}{4\pi r^{2}} \;\propto\; \frac{1}{r^{2}}. \]
\hfill [statement 1; conservation of energy 1; area $4\pi r^{2}$ 1; conclusion 1]

\textbf{(f) Power and distance (3 marks).}
\[ P = 4\pi r^{2} I = 4\pi (4.0)^{2}\times 2.0\times10^{-3} = 4\pi\times 16\times 2.0\times10^{-3} \approx 0.40\ \mathrm{W}. \]
For the new intensity, $I_{1}r_{1}^{2} = I_{2}r_{2}^{2}$:
\[ r_{2} = r_{1}\sqrt{\frac{I_{1}}{I_{2}}} = 4.0\times\sqrt{\frac{2.0\times10^{-3}}{5.0\times10^{-4}}} = 4.0\times\sqrt{4} = 8.0\ \mathrm{m}. \]
\emph{Assumptions:} the loudspeaker behaves as a point source radiating uniformly in all directions, and no energy is absorbed by the medium or reflected by obstacles (free space conditions). \hfill [$P$ 1; $r_{2}$ 1; assumptions 1]
\end{corrige}

\begin{corrige}[Exercise 10 — Structured question: X-rays]
\textbf{(a) The Coolidge tube (6 marks).}
\begin{center}
\begin{tikzpicture}[scale=0.9]
% glass envelope
\draw[thick] (-5,0) ellipse (0.5 and 1.6);
\draw[thick] (-5,1.6) -- (4,1.6);
\draw[thick] (-5,-1.6) -- (4,-1.6);
\draw[thick] (4,0) ellipse (0.5 and 1.6);
\node at (-0.5,2.1) {evacuated glass envelope};
% cathode
\draw[thick,poporange] (-5,0) -- (-3.6,0);
\draw[thick,poporange] (-3.6,-0.5) arc (210:150:1.0);
\node[poporange] at (-4.2,-1.1) {heated cathode (filament)};
% filament supply
\draw[thick] (-6.6,0.4) -- (-5.5,0.4); \draw[thick] (-6.6,-0.4) -- (-5.5,-0.4);
\draw[thick] (-6.6,0.4) -- (-6.6,-0.4);
\node at (-7.5,0) {LT supply};
% anode target
\draw[thick,popdark] (3.2,0) -- (4,0);
\draw[thick,popdark,fill=popgold] (3.2,0) -- (2.6,0.9) -- (2.9,0.9) -- (3.5,0) -- cycle;
\node[popdark] at (3.3,-1.1) {tungsten target (anode)};
% electron beam
\draw[->,thick,popblue,dashed] (-3.3,0) -- (2.4,0.55);
\node[popblue] at (-0.6,-0.45) {electron beam};
% X-rays
\foreach \a in {35,55,75} \draw[->,thick,popgreen] (2.75,0.85) -- ++(\a:1.5);
\node[popgreen] at (4.6,2.2) {X-rays};
% HT
\draw[thick] (-4.2,1.6) -- (-4.2,2.9); \draw[thick] (3.6,1.6) -- (3.6,2.9);
\draw[thick] (-4.2,2.9) -- (3.6,2.9);
\node at (-0.3,3.35) {high voltage ($\sim$ kV)};
\end{tikzpicture}
\legende{Coolidge X-ray tube}
\end{center}
\emph{Production:} a low-voltage current heats the tungsten filament (cathode), which emits electrons by \emph{thermionic emission}. A high voltage (tens of kilovolts) between cathode and anode accelerates the electrons to very high speed in the evacuated tube. The electrons strike the tungsten target and are suddenly decelerated; most of their kinetic energy becomes heat, but a small fraction is converted into electromagnetic radiation of very short wavelength: X-rays. The target is set in a copper block (and often cooled) to conduct the heat away. \hfill [diagram with labels 2; thermionic emission 1; acceleration by HV 1; sudden deceleration at target 1; heat/cooling 1]

\textbf{(b) Two distinguishing properties (2 marks).}
\begin{itemize}
\item X-rays have a much shorter wavelength (much higher frequency) than visible light, hence much greater penetrating power: they pass through flesh, paper and thin aluminium, which are opaque to visible light.
\item X-rays ionise gases (and blacken photographic plates / cause fluorescence) much more strongly than visible light.
\end{itemize}
\hfill [any two, 1 each]

\textbf{(c) Two uses (4 marks).}
\begin{itemize}
\item \emph{Medicine — radiography:} X-rays are absorbed more by bone than by soft tissue (differential absorption linked to their penetrating power), so a shadow image of bones and internal organs is formed on a detector — used to diagnose fractures.
\item \emph{Industry — detecting flaws in metal castings and welds:} their penetrating power lets them pass through metal; cracks and air bubbles absorb less and appear on the film, revealing hidden defects.
\end{itemize}
\hfill [use 1 + property 1, twice]

\textbf{(d) Safety precaution (2 marks).}
Workers stand behind a lead screen or wear lead-lined aprons, and limit exposure time (film badge dosimeters are worn to monitor the dose received). X-rays are ionising radiation: they damage living cells and can cause burns, cancers and genetic mutations, so exposure must be minimised. \hfill [precaution 1; reason 1]

\textbf{(e) Minimum wavelength (4 marks).}
Energy gained by an electron accelerated through $V = 50\ \mathrm{kV}$:
\[ E = eV = 1.60\times10^{-19}\times 50\times10^{3} = 8.0\times10^{-15}\ \mathrm{J}. \]
The shortest-wavelength photon is produced when \emph{all} this energy is converted into a single photon:
\[ eV = hf_{\max} = \frac{hc}{\lambda_{\min}} \;\Rightarrow\; \lambda_{\min} = \frac{hc}{eV} = \frac{6.63\times10^{-34}\times 3.00\times10^{8}}{8.0\times10^{-15}} \approx 2.5\times10^{-11}\ \mathrm{m}. \]
\hfill [$E = eV$ 1; value 1; $eV = hc/\lambda_{\min}$ 1; value with unit 1]

\textbf{(f) Polarisation of X-rays (2 marks).}
Plane polarisation means that the electric field vector of the wave oscillates in one fixed plane only (perpendicular to the direction of propagation). Polarisation is possible only for \emph{transverse} waves: a longitudinal wave vibrates along the propagation direction and has no transverse direction to select. Since X-rays can be polarised (for example by scattering), they must be transverse waves — like all electromagnetic waves. \hfill [definition 1; transverse argument 1]
\end{corrige}

\begin{corrige}[Exercise 11 — Structured question: optical instruments]
\textbf{(a) Normal adjustment (5 marks).}\par
An astronomical telescope is in \emph{normal adjustment} when it is focused so that the final image is at infinity: the image formed by the objective lies in the common focal plane of both lenses, i.e.\ the second focal point $F_{o}$ of the objective coincides with the first focal point $F_{e}$ of the eyepiece. The eye then receives parallel rays and views the image without accommodation (relaxed eye).
\begin{center}
\begin{tikzpicture}[scale=0.9]
% axis
\draw[popmuted] (-1,0) -- (11.5,0);
% objective
\draw[thick,popblue] (1,-1.6) -- (1,1.6);
\draw[thick,popblue] (1,1.6) -- (0.85,1.3); \draw[thick,popblue] (1,1.6) -- (1.15,1.3);
\draw[thick,popblue] (1,-1.6) -- (0.85,-1.3); \draw[thick,popblue] (1,-1.6) -- (1.15,-1.3);
\node[popblue] at (1,-2.0) {objective $f_{o}$};
% eyepiece
\draw[thick,poporange] (9,-1.0) -- (9,1.0);
\draw[thick,poporange] (9,1.0) -- (8.87,0.8); \draw[thick,poporange] (9,1.0) -- (9.13,0.8);
\draw[thick,poporange] (9,-1.0) -- (8.87,-0.8); \draw[thick,poporange] (9,-1.0) -- (9.13,-0.8);
\node[poporange] at (9,-1.4) {eyepiece $f_{e}$};
% common focal point
\fill (8,0) circle (0.06); \node[below] at (8,-0.05) {$F_{o}=F_{e}$};
% incoming parallel rays (off axis)
\draw[thick,popgreen] (-1,1.9) -- (1,1.2);
\draw[thick,popgreen] (-1,1.2) -- (1,0.6);
% refracted through objective to focal plane
\draw[thick,popgreen] (1,1.2) -- (8,-0.42);
\draw[thick,popgreen] (1,0.6) -- (8,-0.42);
% through eyepiece, emerge parallel
\draw[thick,popgreen] (8,-0.42) -- (9,-0.55) -- (11.5,-1.15);
\draw[thick,popgreen] (8,-0.42) -- (9,-0.28) -- (11.5,-0.88);
% intermediate image arrow
\draw[->,thick,popdark] (8,0) -- (8,-0.42);
\node[popgreen] at (-0.6,2.3) {rays from distant object};
\node at (11.2,0.4) {to eye};
\end{tikzpicture}
\legende{Astronomical telescope in normal adjustment}
\end{center}
\hfill [definition of normal adjustment 2; ray diagram: parallel rays in, image at common focus, parallel rays out 3]

\textbf{(b) Magnifying power (3 marks).}\par
The magnifying (angular) power is $M = \dfrac{\beta}{\alpha}$, where $\alpha$ is the angle subtended at the unaided eye by the distant object and $\beta$ the angle subtended by the final image. If $h$ is the height of the intermediate image in the common focal plane, then for small angles:
\[ \alpha \approx \tan\alpha = \frac{h}{f_{o}}, \qquad \beta \approx \tan\beta = \frac{h}{f_{e}} \;\Rightarrow\; M = \frac{\beta}{\alpha} = \frac{f_{o}}{f_{e}}. \]
Numerically: $M = \dfrac{100}{5.0} = 20$. \hfill [angles 1; ratio 1; value 1]

\textbf{(c) Length of the telescope (2 marks).}\par
In normal adjustment the lenses are separated by the sum of their focal lengths:
\[ L = f_{o} + f_{e} = 100 + 5.0 = 105\ \mathrm{cm}. \]
\hfill [formula 1; value 1]

\textbf{(d) Large-diameter objective (2 marks).}\par
A large objective collects more light (light-gathering power proportional to its area, hence to $D^{2}$), so faint, distant stars and galaxies become visible; it also gives greater resolving power (less diffraction at the aperture), allowing fine detail and close pairs of stars to be distinguished. \hfill [brightness 1; resolution 1]

\textbf{(e) Microscope objective vs telescope objective (4 marks).}\par
For a compound microscope, the magnifying power is larger when the objective focal length is \emph{short}: the object is placed just outside $f_{o}$ so that the objective forms a real, highly magnified intermediate image (the linear magnification of the objective is large when $f_{o}$ is small). Moreover, the object is very close to the lens, so a short-focus lens can be used.\par
For a telescope, $M = f_{o}/f_{e}$, so a \emph{long} focal length objective gives high magnifying power; it also forms a larger real image of the distant object at its focal plane, which the eyepiece then magnifies. \hfill [microscope reason 2; telescope reason 2]

\textbf{(f) Critical angle and the TIR prism (4 marks).}
\[ \sin C = \frac{1}{n} = \frac{1}{1.50} = 0.667 \;\Rightarrow\; C \approx 41.8^{\circ}. \]
In a right-angled ($45^{\circ}$--$45^{\circ}$--$90^{\circ}$) prism, light enters one short face normally (without refraction) and strikes the long (hypotenuse) face at $45^{\circ}$. Since $45^{\circ} > C \approx 41.8^{\circ}$, total internal reflection occurs at the glass--air surface: the prism acts as a perfect mirror, turning the rays through $90^{\circ}$ (or through $180^{\circ}$ in a pair of Porro prisms) with no loss of light and no silvering to deteriorate — which is why such prisms are used in binoculars and periscopes. \hfill [$C$ value 1; incidence at $45^{\circ}$ 1; comparison $45^{\circ}>C$ 1; application/advantage 1]
\end{corrige}

\newpage

\coursec{Summary sheet}

\begin{popfigure}
\begin{center}\tcbox[colback=popdark,colframe=popdark,coltext=white,arc=9pt,boxsep=4pt,left=6pt,right=6pt]{\bfseries\small Electromagnetic Waves}\end{center}
\vspace{-2pt}
\begin{tcbraster}[raster columns=3,raster equal height=rows,raster column skip=6pt,raster row skip=6pt,enhanced,blankest]
\begin{tcolorbox}[enhanced,colback=white,colframe=popblue,boxrule=0.9pt,arc=3pt,title={Nature of EM waves},fonttitle=\bfseries\scriptsize,coltitle=white,colbacktitle=popblue,left=4pt,right=4pt,top=2pt,bottom=2pt,boxsep=1pt,toptitle=1.5pt,bottomtitle=1.5pt,halign=left]
\begin{itemize}[nosep,leftmargin=*,itemsep=1pt,font=\scriptsize]
\item $c=f\lambda$
\item Transverse $\vec{E}\perp\vec{B}$
\item Speed of sound $\approx\SI{340}{m/s}$
\end{itemize}
\end{tcolorbox}
\begin{tcolorbox}[enhanced,colback=white,colframe=poporange,boxrule=0.9pt,arc=3pt,title={EM spectrum \& X-rays},fonttitle=\bfseries\scriptsize,coltitle=white,colbacktitle=poporange,left=4pt,right=4pt,top=2pt,bottom=2pt,boxsep=1pt,toptitle=1.5pt,bottomtitle=1.5pt,halign=left]
\begin{itemize}[nosep,leftmargin=*,itemsep=1pt,font=\scriptsize]
\item Radio $\to$ gamma
\item X-ray tube
\item Polarisation
\end{itemize}
\end{tcolorbox}
\begin{tcolorbox}[enhanced,colback=white,colframe=popgreen,boxrule=0.9pt,arc=3pt,title={Diffraction \& gratings},fonttitle=\bfseries\scriptsize,coltitle=white,colbacktitle=popgreen,left=4pt,right=4pt,top=2pt,bottom=2pt,boxsep=1pt,toptitle=1.5pt,bottomtitle=1.5pt,halign=left]
\begin{itemize}[nosep,leftmargin=*,itemsep=1pt,font=\scriptsize]
\item $d\sin\theta=n\lambda$
\item Multiple slits
\end{itemize}
\end{tcolorbox}
\begin{tcolorbox}[enhanced,colback=white,colframe=poppurple,boxrule=0.9pt,arc=3pt,title={Reflection \& refraction},fonttitle=\bfseries\scriptsize,coltitle=white,colbacktitle=poppurple,left=4pt,right=4pt,top=2pt,bottom=2pt,boxsep=1pt,toptitle=1.5pt,bottomtitle=1.5pt,halign=left]
\begin{itemize}[nosep,leftmargin=*,itemsep=1pt,font=\scriptsize]
\item Snell: $n_1\sin i=n_2\sin r$
\item $n=c/v$
\item Dispersion
\end{itemize}
\end{tcolorbox}
\begin{tcolorbox}[enhanced,colback=white,colframe=poppink,boxrule=0.9pt,arc=3pt,title={Total internal reflection},fonttitle=\bfseries\scriptsize,coltitle=white,colbacktitle=poppink,left=4pt,right=4pt,top=2pt,bottom=2pt,boxsep=1pt,toptitle=1.5pt,bottomtitle=1.5pt,halign=left]
\begin{itemize}[nosep,leftmargin=*,itemsep=1pt,font=\scriptsize]
\item $\sin C=1/n$
\item Fibres, prisms
\end{itemize}
\end{tcolorbox}
\begin{tcolorbox}[enhanced,colback=white,colframe=popgold,boxrule=0.9pt,arc=3pt,title={Lenses \& instruments},fonttitle=\bfseries\scriptsize,coltitle=white,colbacktitle=popgold,left=4pt,right=4pt,top=2pt,bottom=2pt,boxsep=1pt,toptitle=1.5pt,bottomtitle=1.5pt,halign=left]
\begin{itemize}[nosep,leftmargin=*,itemsep=1pt,font=\scriptsize]
\item $\frac{1}{f}=\frac{1}{u}+\frac{1}{v}$
\item Microscope, telescope
\end{itemize}
\end{tcolorbox}
\begin{tcolorbox}[enhanced,colback=white,colframe=popmuted,boxrule=0.9pt,arc=3pt,title={Energy of waves},fonttitle=\bfseries\scriptsize,coltitle=white,colbacktitle=popmuted,left=4pt,right=4pt,top=2pt,bottom=2pt,boxsep=1pt,toptitle=1.5pt,bottomtitle=1.5pt,halign=left]
\begin{itemize}[nosep,leftmargin=*,itemsep=1pt,font=\scriptsize]
\item $I\propto 1/r^{2}$
\item $I=P/(4\pi r^{2})$
\end{itemize}
\end{tcolorbox}
\end{tcbraster}

\legende{Mind map of the chapter \emph{Electromagnetic Waves}: the seven key ideas to master for the GCE Advanced Level.}
\end{popfigure}

\begin{bilan}
\textbf{Key definitions}
\begin{itemize}
\item An \cle{electromagnetic wave} is a self-propagating disturbance of perpendicular electric ($\vec{E}$) and magnetic ($\vec{B}$) fields, both perpendicular to the direction of travel: EM waves are \cle{transverse} and need no material medium.
\item All EM waves travel in vacuum at $c=\SI{3.00e8}{m/s}$, with $c=f\lambda$.
\item The \cle{refractive index} of a medium: $n=\dfrac{c}{v}=\dfrac{\sin i}{\sin r}$ (air $\to$ medium).
\item \cle{Total internal reflection} occurs when light passes from dense to less dense medium with $i>C$, where the \cle{critical angle} satisfies $\sin C=\dfrac{n_{2}}{n_{1}}$ ($n_1>n_2$).
\item \cle{Plane polarisation}: vibrations of $\vec{E}$ confined to one plane; only possible for transverse waves — proof that light is transverse.
\item \cle{Intensity} at distance $r$ from a point source radiating power $P$ uniformly: $I=\dfrac{P}{4\pi r^{2}}$ (inverse square law).
\end{itemize}

\textbf{Laws and formulas to know}
\begin{itemize}
\item Wave relation: $v=f\lambda$; for sound in air $v\approx\SI{340}{m/s}$ at room temperature (measured by resonance tube or echo timing).
\item Snell's laws: incident ray, refracted ray and normal are coplanar; $n_1\sin i=n_2\sin r$.
\item Dispersion: $n$ depends on $\lambda$ ($n_{\text{violet}}>n_{\text{red}}$), so a prism splits white light; deviation is greatest for violet.
\item Diffraction grating (normal incidence): $d\sin\theta_n=n\lambda$, with $d=\dfrac{1}{N}$ the slit spacing ($N$ = lines per metre). More slits $\Rightarrow$ sharper, brighter maxima.
\item Thin converging lens: $\dfrac{1}{f}=\dfrac{1}{u}+\dfrac{1}{v}$ (real-is-positive convention); magnification $m=\dfrac{v}{u}$; power $P=\dfrac{1}{f(\text{m})}$ in \cle{dioptres} ($\delta$).
\item Compound microscope: $M=M_{\text{obj}}\times M_{\text{eyepiece}}$; astronomical telescope in normal adjustment: $M=\dfrac{f_{\text{obj}}}{f_{\text{eye}}}$.
\item Inverse square law: $\dfrac{I_1}{I_2}=\dfrac{r_2^{2}}{r_1^{2}}$; amplitude falls as $1/r$.
\end{itemize}

\textbf{Methods}
\begin{itemize}
\item \emph{Grating problems:} convert lines/mm to $d$ in metres, then apply $d\sin\theta=n\lambda$; check $\sin\theta\le 1$ for the maximum order.
\item \emph{Refraction problems:} draw the normal first, measure angles \emph{from the normal}, apply Snell, check the ray bends towards the normal entering a denser medium.
\item \emph{Lens problems:} sketch a ray diagram (ray through optical centre undeviated; ray parallel to axis through $F'$), then confirm numerically with the lens formula.
\item \emph{Speed of sound:} $v=\dfrac{2d}{t}$ for echoes, or $v=4fL$ (with end correction) for the resonance tube.
\end{itemize}

\textbf{Common mistakes}
\begin{itemize}
\item Measuring angles from the surface instead of from the \emph{normal}.
\item Forgetting that frequency $f$ is unchanged on refraction — only $v$ and $\lambda$ change.
\item Using $\sin C=n$ instead of $\sin C=1/n$; TIR is impossible going from less dense to more dense.
\item Mixing units: $d$ in mm, $\lambda$ in nm in the grating formula.
\item Sign errors in the lens formula for virtual images ($v<0$) and diverging lenses ($f<0$).
\item Believing the inverse square law applies to amplitude: it applies to \emph{intensity}.
\end{itemize}

\textbf{GCE tips}
\begin{itemize}
\item Always state the condition (normal incidence, dense-to-rare, point source in free space) before quoting a formula — examiners award marks for conditions.
\item In structured questions on X-rays, mention: hot-cathode tube, high vacuum, accelerating p.d., metal target; uses in radiography and crystal analysis.
\item Give every numerical answer with its SI unit and 2--3 significant figures; powers of lenses in dioptres require $f$ in metres.
\item Practise drawing neat ray diagrams for the microscope and telescope: they recur almost every year.
\end{itemize}
\end{bilan}

\findecours

\end{document}
